\documentclass[11pt,reqno]{amsart}
\usepackage[T1]{fontenc}
\usepackage[utf8]{inputenc}
\usepackage{lmodern}
\usepackage[margin=1.08in]{geometry}
\usepackage{amsmath,amssymb,amsthm,mathtools}
\usepackage[expansion=false]{microtype}
\usepackage{enumitem}
\usepackage{graphicx}
\usepackage{tikz-cd}
\usepackage{booktabs}
\usepackage{needspace}
\usepackage[colorlinks=true,linkcolor=blue!50!black,citecolor=blue!50!black,urlcolor=blue!50!black]{hyperref}
\makeatletter
\g@addto@macro\UrlBreaks{\do\-}
\makeatother
\hypersetup{
 pdftitle={The reverse logarithmic Kodaira inequality and additivity},
 pdfauthor={OpenAI},
 pdfsubject={Reverse logarithmic Kodaira inequality and additivity for stratum-smooth reduced SNC projective pairs},
 pdfkeywords={logarithmic Kodaira dimension, Hodge theory, period maps, canonical bundle formula}
}
\usepackage[capitalise,nameinlink,noabbrev]{cleveref}

\newcommand{\C}{\mathbb C}
\newcommand{\Q}{\mathbb Q}
\newcommand{\Z}{\mathbb Z}
\newcommand{\R}{\mathbb R}
\newcommand{\OO}{\mathcal O}
\DeclareMathOperator{\Supp}{Supp}
\DeclareMathOperator{\rk}{rk}
\DeclareMathOperator{\Gr}{Gr}
\DeclareMathOperator{\Sym}{Sym}
\DeclareMathOperator{\End}{End}
\DeclareMathOperator{\Hom}{Hom}
\DeclareMathOperator{\Aut}{Aut}
\DeclareMathOperator{\Lie}{Lie}
\DeclareMathOperator{\PGL}{PGL}
\DeclareMathOperator{\GL}{GL}
\DeclareMathOperator{\Pic}{Pic}
\DeclareMathOperator{\tr}{tr}
\DeclareMathOperator{\ord}{ord}
\DeclareMathOperator{\lct}{lct}
\DeclareMathOperator{\im}{im}
\DeclareMathOperator{\codim}{codim}

\theoremstyle{plain}
\newtheorem{theorem}{Theorem}[section]
\newtheorem{proposition}[theorem]{Proposition}
\newtheorem{lemma}[theorem]{Lemma}
\newtheorem{corollary}[theorem]{Corollary}
\newtheorem{claim}[theorem]{Claim}
\theoremstyle{definition}
\newtheorem{definition}[theorem]{Definition}
\newtheorem{convention}[theorem]{Convention}
\theoremstyle{remark}
\newtheorem{remark}[theorem]{Remark}
\newtheorem{example}[theorem]{Example}
\crefname{equation}{Equation}{Equations}
\crefname{section}{Section}{Sections}
\numberwithin{equation}{section}
\setlist[enumerate]{label=\textup{(\roman*)},leftmargin=*,itemsep=3pt}
\setlist[itemize]{leftmargin=*,itemsep=3pt}
\allowdisplaybreaks[2]
\emergencystretch=1.5em

\title[Reverse logarithmic Kodaira inequality]{The reverse logarithmic Kodaira inequality and additivity}
\author{OpenAI}
\date{September 26, 2026}
\subjclass[2020]{14D07, 14E30, 14C30, 14J10}
\keywords{logarithmic Kodaira dimension, algebraic fiber space, variation of Hodge structure, period map, canonical bundle formula}

\newcommand{\FoundationLogNumber}{6.2}


\begin{document}
\begin{abstract}
We prove the reverse logarithmic Kodaira inequality for a surjective
connected-fiber morphism $f:(X,E)\to(Y,D)$ of smooth projective
reduced simple-normal-crossing pairs, with
$\Supp(f^*D)\subseteq\Supp E$, such that $X$ and every boundary
stratum are smooth over $Y\setminus\Supp D$. Together with logarithmic
subadditivity, the inequality gives additivity, including both
negative-infinity cases. This resolves Popa's logarithmic additivity
conjecture positively in the projective reduced-SNC, stratum-smooth setting.
\end{abstract}
\maketitle
\tableofcontents

\section{Introduction}\label{sec:introduction}


Let a smooth family be completed to a morphism of projective pairs.
How many logarithmic pluricanonical forms can its total space have?
Subadditivity gives a lower bound from the base and a general fiber.
We prove the complementary upper bound when the total space and all
boundary strata are smooth over the open base. The boundary may have
horizontal components, and the conclusion includes the case in which
the base has no logarithmic pluricanonical forms at all.

Throughout, varieties are over $\C$. Let $f:X\to Y$ be a surjective
morphism with connected fibers between smooth connected projective
varieties. Let $E$ and $D$ be reduced simple-normal-crossing divisors
on $X$ and $Y$, respectively, with
\begin{equation}\label{eq:support-inclusion}
 \Supp(f^*D)\subseteq\Supp(E).
\end{equation}
Either divisor may be zero. Set $V=Y\setminus\Supp D$. We assume that
$X$ and every irreducible component of every nonempty intersection of
components of $E$, restricted to $f^{-1}(V)$, are smooth over $V$.
We call this \emph{stratum smoothness over $V$}. For a very general
$y\in V$, set $F=X_y$ and $E_F=E|_F$. Write
\[
 L_X=K_X+E,\qquad L_Y=K_Y+D,\qquad L_F=K_F+E_F.
\]
We give a point Kodaira dimension zero and use
$(-\infty)+a=-\infty$, also when $a=-\infty$.

\begin{theorem}[Reverse logarithmic Kodaira inequality]\label{thm:upper-bound}
Under these hypotheses,
\begin{equation}\label{eq:upper-bound}
 \kappa(X,L_X)\leq\kappa(Y,L_Y)+\kappa(F,L_F).
\end{equation}
If either term on the right is $-\infty$, then
$H^0(X,mL_X)=0$ for every $m>0$.
\end{theorem}

There is no abundance, semiampleness, good-minimal-model, or general-type
hypothesis. Condition \eqref{eq:support-inclusion} is a condition on supports;
the multiplicities in $f^*D$ remain in the proof. No smoothness is
assumed over $D$.

The additional input for equality is the logarithmic subadditivity
theorem of the companion article \cite[Corollary~\FoundationLogNumber]{Foundation}.
It applies to any connected-fiber surjective morphism of smooth
projective reduced-SNC pairs satisfying \eqref{eq:support-inclusion}, without
the stratum-smoothness assumption. Its fiber is the same very general
pair $(F,E_F)$, and it gives the opposite inequality. Therefore:

\begin{corollary}[Logarithmic additivity]\label{cor:additivity}
For the pairs in Theorem~\ref{thm:upper-bound},
\begin{equation}\label{eq:additivity}
 \kappa(X,K_X+E)=\kappa(Y,K_Y+D)+\kappa(F,K_F+E_F).
\end{equation}
\end{corollary}

This gives a positive resolution of Popa's logarithmic additivity
conjecture in its projective
reduced-SNC, stratum-smooth formulation
\cite[Conjecture 3.9 and its defining footnote]{PopaSurvey}.
The proof of the upper inequality is independent of logarithmic
subadditivity; the companion theorem is used only for this corollary.

\subsection{The section-construction result}

The difficult case is a base with logarithmic Kodaira dimension zero.
The space of logarithmic pluricanonical sections then has dimension
at most one in each degree. Choose a very general point outside the zero divisors of all
these sections. If a total-space pluriform vanished on its fiber, we
would like to produce a base pluriform vanishing at that point. This
would be a contradiction, and restriction to the fiber would be
injective in every degree.

The following result carries out that construction. It also treats a
base of logarithmic Kodaira dimension $-\infty$.

\begin{proposition}\label{prop:section-construction}
Let $f,E,D$ satisfy the hypotheses of Theorem~\ref{thm:upper-bound}, and
let $0\ne s\in H^0(X,mL_X)$ for some $m>0$. Then
$\kappa(Y,L_Y)\geq0$. If $\kappa(Y,L_Y)=0$ and $s$ vanishes identically
on $X_y$ for some $y\in V$, there are an integer $a>0$ and a nonzero
$\tau\in H^0(Y,aL_Y)$ with $\tau(y)=0$.
\end{proposition}

The point $y$ in this statement belongs to the original open base
$V$. It need not belong to a smaller open set chosen in constructing
a cyclic cover from $s$. This uniformity is what allows the restriction
argument to handle all total-space sections at one very general fiber.
In particular, if $\kappa(Y,L_Y)=0$, restriction to one very general
fiber is injective in every positive degree, so
\[
 h^0(X,mL_X)\leq h^0(F,mL_F)
 \qquad\text{for every positive integer }m.
\]
\Cref{sec:completion} proves \Cref{thm:upper-bound} from this proposition,
including both negative-infinity cases.

\subsection{Two outcomes of the Hodge construction}

A pluriform on the total space first produces a cohomology class:
take a root of the pluriform on a finite cover, then project its class
to a nonzero pure weight quotient. This gives a vector in the bundle
of highest Hodge classes. The chosen vector need not span that bundle.
Thus positivity of the determinant of the whole bundle does not yet
control the coordinates of this particular vector, or retain its zero
on a chosen fiber.

We keep track of these coordinates in the logarithmic differential
frames of the original pair $(Y,D)$. The root cover may acquire extra
singular fibers inside $V$, but its new discriminant must not enlarge
the poles allowed on the base. A local calculation shows both that
no new poles occur and that a zero on any original smooth fiber is
retained. Homogeneous tensor operations preserve these properties
while making the cohomological data descend to a projective parameter
space $S$. This space records the varying Hodge structure and the
additional vector directions that its period map alone can miss.
There are two possibilities.

If $S$ is a point, the resulting cohomology bundle is constant.
Its vector has ordinary scalar coordinates, at least one of which
is nonzero. The coefficient calculation makes each coordinate an
actual base pluriform. The same calculation preserves the prescribed
zero. This proves base nonvanishing, excluding a base of logarithmic Kodaira
dimension $-\infty$. The preserved zero gives the additional conclusion
needed when its logarithmic Kodaira dimension is zero.

If $S$ has positive dimension, restrict first to the generic fiber
of the map from the base to $S$. The coefficient system is constant
there, so a nonzero scalar coordinate permits a relative Iitaka
fibration. It produces a divisor on an intermediate variety with the
same sufficiently divisible section spaces as $K_Y+D$. Period
positivity makes this divisor big after adding a pullback from $S$.
The remaining obstacle is to remove that added divisor. We do so by
comparing the iterated Higgs images of the chosen vector with the
period directions and applying a numerical subtraction argument.
The divisor itself is then big, which forces
$\kappa(Y,K_Y+D)>0$. Thus this alternative cannot occur when the
base has logarithmic Kodaira dimension zero or $-\infty$.

\subsection{Earlier results and the proof route}

Iitaka's subadditivity problem asks whether
the Kodaira dimension of an algebraic fiber space is at least the sum
of those of its base and general fiber. Viehweg proved the
relative-dimension-one case \cite[Introduction]{Viehweg1977}.
Logarithmic pluriforms bring open varieties into the same birational
framework \cite[Introduction]{IitakaLogForms1976}.

For smooth families, the stronger question is whether the lower bound
becomes equality. Popa's logarithmic formulation specifies smoothness
of the total space and every boundary stratum over the open base
\cite[Conjecture 3.9 and footnote 5]{PopaSurvey}. The reverse inequality
constrains the extra pluricanonical growth that variation might
otherwise produce. Its negative-base case makes the distinction from
subadditivity especially clear: the lower bound is then automatic,
whereas the upper bound asserts vanishing of every positive
pluricanonical space on the total pair.

Popa--Schnell proved the upper inequality for smooth projective
families under the Campana--Peternell conjecture and obtained a
degreewise plurigenus bound when the base has rationally trivial
canonical bundle \cite[Theorems C and H]{PopaSchnellSmooth}.
Their reduction through the base Iitaka fibration also organizes our
argument. The section-construction proposition adds the specific
conclusion needed here: an actual base pluriform, with a prescribed
zero whenever the original section vanishes on a chosen smooth fiber.

Fujino--Fujisawa's logarithmic results already allow relatively
normal-crossing horizontal boundary. Their upper inequality for a base of nonnegative logarithmic Kodaira
dimension assumes generalized abundance for sufficiently general fibers of the base's
logarithmic Iitaka fibration \cite[Theorem 1.12]{FF}; their Hodge-theoretic construction also
descends pseudoeffectivity to the base
\cite[Theorems 1.4 and 4.5]{FF}.
Park proves the logarithmic upper inequality, in the reduced
inverse-image boundary setting and for a base of nonnegative
logarithmic Kodaira dimension, assuming good minimal models for the
very general fibers of that base Iitaka fibration
\cite[Theorem 1.12]{Park}. These hypotheses concern fibers of the
\emph{base's} Iitaka fibration, rather than the original fibers of $f$.

A different route is Campana's logarithmic additivity theorem for
smooth projective families over a quasiprojective base with semiample
canonical bundles on the original fibers
\cite[Theorem 1]{CampanaAdditivity}. Its proof uses birational
isotriviality and the core. The extension to original fibers with good
minimal models uses Taji's birational-isotriviality theorem over a
special base \cite[Theorem 1.1]{Taji}, as explained in
\cite[Remark 2(1)]{CampanaAdditivity}.
Theorem~\ref{thm:upper-bound} imposes neither of these minimal-model
conditions.

The differential construction has an earlier model in Kov\'acs's
iteration of cohomology maps from logarithmic cotangent sequences
\cite[Section 1, Lemma 1.1]{KovacsFineModuli}. The passage from a
pluriform to a cyclic cover and iterated Higgs maps belongs to the
Viehweg--Zuo method
\cite[Introduction, Lemma 1.1, and Section 3]{ViehwegZuoIsotriviality}.
We use its Hodge-theoretic development in Popa--Schnell and the
logarithmic construction of Fujino--Fujisawa, while retaining the
chosen vector and its original coefficient lattice. Further inputs
are full integral period-image algebraicity \cite{BBT}, orbifold
cotangent positivity \cite{CP}, and b-nefness of the moduli part of
an lc-trivial fibration \cite{FG}. For the last input we verify the
rank condition for an auxiliary subpair which may have negative
coefficients. The new task addressed by the section construction is
the passage from these positivity statements to actual base sections
with the stated control at every point of the original open base.

\paragraph{Organization.}
\Cref{sec:completion} deduces the upper inequality in all sign cases from
\Cref{prop:section-construction}. The rest of the paper proves that
proposition. \Cref{sec:logarithmic-input} constructs the geometric vector,
its coefficient lattice, and the zero test. \Cref{sec:marked-periods}
constructs the quotient that retains its periods and compact-factor
directions. \Cref{sec:boundary-rank} proves the two comparisons that remove
the boundary and control the chosen vector by its iterated Higgs images.
\Cref{sec:relative-iitaka} produces an intermediate base with exactly the
original logarithmic section spaces. \Cref{sec:nonvanishing} first proves
the constant-parameter case, then removes the period twist in the
positive-dimensional case and completes the section construction. Metric and
flat-connection estimates appear immediately before their first uses.

\section{Reduction to the section construction}\label{sec:completion}

The sole input from the remaining sections is
\Cref{prop:section-construction}. We first give the reduction in every
sign case. This identifies exactly what the geometric construction must
produce, including its assertion at a prescribed point of the original
open base.

\subsection{Conventions and birational changes}
All varieties are over $\C$, except when a generic fiber is explicitly
considered over a characteristic-zero function field.  We use additive
notation for line bundles and rational Cartier divisors.  A statement about
sections of a rational divisor is understood after taking a positive
Cartier multiple.  For numerical divisor classes, $P\ge_{\mathrm{pe}}Q$
means that $P-Q$ is pseudoeffective.

We make projective resolutions and resolve rational maps throughout.
Resolution may be taken to preserve a specified smooth open set and to be
equivariant for a finite group; see
\cite[Theorems 3.35 and 3.36]{KollarResolution}.
Here is the boundary convention that accompanies a birational
modification.  If $\mu:Z'\to Z$ is a birational morphism between smooth
projective varieties and $\Delta$ is reduced SNC, put
\[
 \Delta'=\mu^{-1}_*\Delta+\operatorname{Exc}(\mu)_{\mathrm{red}},
\]
after further resolution when necessary.  The log canonicity of the
smooth pair $(Z,\Delta)$ gives
\begin{equation}\label{eq:prelim-birational-log}
 K_{Z'}+\Delta'=\mu^*(K_Z+\Delta)+A_\mu,
 \qquad A_\mu\geq0\quad\text{$\mu$-exceptional}.
\end{equation}
Consequently pullback identifies the spaces of logarithmic pluricanonical
sections in every degree.  Indeed, pullback supplies one inclusion; in the
other direction a rational section regular away from a set of codimension
at least two on the smooth variety $Z$ extends across that set.  Equivalently,
$\mu_*\OO_{Z'}(mA_\mu)=\OO_Z$.  This convention also applies to a generic
fiber after a characteristic-zero extension of the ground field.

We use the standard numerical properties of big and nef divisors
\cite[Chapters 1--2]{Lazarsfeld}.  In particular, the big cone is the
interior of the pseudoeffective cone, and adding a pseudoeffective class
to a big class preserves bigness.  If $P$ is nef on an $n$-dimensional
projective variety, its numerical dimension is
\begin{equation}\label{eq:prelim-numerical-dimension}
 \nu(P)=\max\{k\in\{0,\ldots,n\}:P^k A^{n-k}>0\},
 \qquad A\ \text{ample},
\end{equation}
independently of $A$, and $P$ is big exactly when $\nu(P)=n$.
Bigness is preserved in both directions by a dominant generically finite
map.  For descent, factor the map through its finite normalization and
apply the field norm to an effective multiple of the pulled-back divisor
minus a sufficiently small pullback of an ample divisor.  Birational
invariance then gives the assertion on a resolution.  On a point the empty
intersection product is $1$, and numerical and Iitaka dimensions of the
zero line are both zero.

\subsection{The two negative-infinity cases}

\phantomsection\label{proof:fiber-vanishing}
Suppose first that $\kappa(F,L_F)=-\infty$. A nonzero section of $mL_X$
cannot vanish identically on every very general fiber of a dominant map:
its nonvanishing locus is a nonempty open subset of $X$, whose image
contains a dense open subset of $Y$. Restriction and adjunction would
therefore give a nonzero section of $mL_F$ on a very general fiber. This is
impossible. Hence $\kappa(X,L_X)=-\infty$.

\phantomsection\label{proof:base-vanishing}
Suppose next that $\kappa(Y,L_Y)=-\infty$. A hypothetical nonzero section
of $mL_X$ contradicts the first assertion of
\Cref{prop:section-construction}.
Thus the total Kodaira dimension is again $-\infty$. This is an actual
vanishing conclusion, not a consequence of an automatic lower inequality.

\subsection{A base of Kodaira dimension zero}

Assume $\kappa(Y,L_Y)=0$. If $Y$ is a point, the conclusion is immediate.
Otherwise choose $y\in V$ very general so that all fiber pluriform spaces
have their very general dimensions and $y$ lies outside the zero divisor
of every nonzero base pluriform. This last condition is possible:
\begin{equation}\label{eq:base-one-section}
 h^0(Y,aL_Y)\leq 1\qquad(a>0).
\end{equation}
Indeed, two linearly independent sections in a common degree have a
nonconstant ratio and hence give Iitaka dimension at least one. There are
only countably many degrees to consider.

For every $m>0$, restriction is injective:
\begin{equation}\label{eq:fiber-injective}
 H^0(X,mL_X)\longrightarrow H^0(F,mL_F).
\end{equation}
To prove this, suppose that a nonzero section on the left vanishes on $F$.
By \Cref{prop:section-construction}, this gives a nonzero section of a
positive multiple of $L_Y$ vanishing at $y$. That proposition applies at
every point of the original $V$, regardless of the auxiliary open sets
used for this particular section. This contradicts the choice of
$y$. Thus \eqref{eq:fiber-injective} holds in every degree.

The usual growth characterization of Iitaka dimension, applied to
\eqref{eq:fiber-injective}, gives
\begin{equation}\label{eq:zero-base-bound}
 \kappa(X,L_X)\leq\kappa(F,L_F).
\end{equation}
One may use sufficiently divisible degrees on both sides; the elementary
growth statement does not require finite generation of either section
ring, cf.\ \cite[Section~2.1]{Lazarsfeld}.

\subsection{Positive and intermediate base dimensions}

Let $k=\kappa(Y,L_Y)>0$. We use the reduction through the base Iitaka
fibration of Popa--Schnell \cite[\S2, Step~2]{PopaSchnellSmooth}, retaining
the reduced boundary and its strata. Resolve a logarithmic Iitaka fibration
$q\colon Y'\to T$, with $Y',T$ smooth projective and connected fibers, so
that $\dim T=k$. Put $D'$ equal to the strict transform of $D$ plus the
reduced exceptional divisor, with further resolution to SNC. The section
spaces of $K_{Y'}+D'$ agree with those of $L_Y$.

The very general fiber $G$ of $q$, with $D_G=D'|_G$, has
\begin{equation}\label{eq:iitaka-fiber-zero}
 \kappa(G,K_G+D_G)=0.
\end{equation}
Here the Iitaka map is taken with the relative algebraic closure of its
image field. The standard section-ratio argument proving
\eqref{eq:iitaka-fiber-zero} is the same one used in
\Cref{lem:relative-iitaka-construction}: nonzero forms restrict to the general fiber,
and an additional independent ratio there would spread after clearing
vertical poles with a sufficiently positive twist from the image. The
resolved original system dominates that twist, contradicting maximality
of the Iitaka image. No abundance statement is being used.

Take the main-component birational base change of $X$ to $Y'$ and resolve
it compatibly with the boundary. Denote the resulting morphism by
$f'\colon X'\to Y'$ and its boundary by $E'$, the strict transform of $E$
plus the reduced exceptional divisor. We can preserve the smooth pair
over $Y'\setminus D'$; the support inclusion holds because over $V$ the
pullback of a center of codimension at least two still has codimension at
least two, whereas the original inverse image of $D$ already lies in the
boundary. Thus all necessary new divisors are exceptional on the source
or lie over the old boundary. The morphism has connected fibers: it has
geometrically connected generic fiber, and its Stein factor is finite
birational over the normal base. These modifications preserve
$\kappa(X,L_X)$ and $\kappa(Y,L_Y)$.

For a very general point $t\in T$, put $G=q^{-1}(t)$ and
$H=(q\circ f')^{-1}(t)$. By generic smoothness applied to the finitely
many boundary strata, these are smooth projective varieties with their
induced reduced SNC boundaries.  Locally, surjectivity on the deepest
stratum lets the boundary parameters and the parameters pulled back from
$T$ be part of one coordinate system; restriction therefore retains
distinct reduced normal-crossing boundary branches.  The induced map
\[
 f_t\colon (H,E'|_H)\longrightarrow(G,D'|_G)
\]
has connected fibers and satisfies the original stratum-smoothness
hypothesis over $G\setminus D'|_G$. The condition over that open is
preserved by base change. Adjunction gives the restrictions of the two
log canonical divisors. The very general fiber invariant of $f_t$ is
$\kappa(F,L_F)$: choose $t$ and then the point of $G$ outside the
countably many loci governing all section dimensions. Applying
\eqref{eq:zero-base-bound}, or the already proved negative-infinity fiber
case, yields
\[
 \kappa(H,K_H+E'|_H)\leq\kappa(F,L_F).
\]
For completeness, take any nonempty total pluriform system, with
rational image map $\phi$. The image of the combined map
$(q\circ f',\phi)$ has generic fiber over $T$ equal to the image of the
restricted subsystem. Its fiber dimension is at most
$\kappa(H,K_H+E'|_H)$. Projection to the image of $\phi$ therefore gives
$\dim\operatorname{im}\phi\leq\dim T+\kappa(H,K_H+E'|_H)$.
Taking the supremum over degrees gives
\[
 \kappa(X',K_{X'}+E')
 \leq \dim T+\kappa(H,K_H+E'|_H)
 \leq k+\kappa(F,L_F).
\]
If the middle fiber Kodaira dimension is $-\infty$, restriction excludes
all total sections instead. This proves \Cref{thm:upper-bound}, including
the case $k=\dim Y$ and hence a point Iitaka fiber.

\section{A logarithmic Hodge vector and its coefficient lattice}
\label{sec:logarithmic-input}

Let $f:(X,E)\to(Y,D)$ satisfy the hypotheses of
\Cref{thm:upper-bound}, and put $d=\dim X-\dim Y$ and $L_Y=K_Y+D$.
Fix one total-space logarithmic pluriform
\begin{equation}\label{eq:log-input-section}
  0\ne s\in H^0\bigl(X,\OO_X(m(K_X+E))\bigr),\qquad m>0.
\end{equation}
We turn $s$ into a highest Hodge vector whose coefficients retain
the original base cotangent lattice. The open set on which a cyclic
cover constructed from $s$ is a smooth family will usually be smaller
than $V=Y\setminus D$. The differential lattice nevertheless remains
$\Omega_Y^1(\log D)$, including at divisors in $V$ where that cover
degenerates. Retaining this lattice will later turn the Hodge vector
into actual base sections.

\phantomsection\label{sec:preliminaries}
\subsection{Variations and extension conventions}

A polarized rational variation of pure Hodge structure of weight $w$ on a
smooth quasi-projective variety $B^\circ$ consists of a rational local
system $\mathbb V_{\Q}$, its flat holomorphic bundle
$\mathcal V=\mathbb V_{\Q}\otimes_{\Q}\OO_{B^\circ}$ with connection
$\nabla$, a decreasing holomorphic filtration $F^\bullet\mathcal V$,
and a flat polarization.  The filtration defines polarized pure Hodge
structures on fibers and satisfies Griffiths transversality
\[
 \nabla(F^p\mathcal V)\subseteq
 F^{p-1}\mathcal V\otimes\Omega^1_{B^\circ}.
\]
The initial variations in this paper have invariant lattices; torsion is
discarded when taking such a lattice.  A complex variation used below is
obtained, after a specified finite level if necessary, as a polarized
complex Hodge summand of a tensor construction in these rational
variations and their duals.  In particular its projector is flat and
preserves the Hodge decomposition.  Such a summand need not have its own
real form.  It has the induced positive Hodge metric and a flat Hermitian
polarization.  Constant changes of the Hodge indices in these constructions
do not change the highest Hodge spaces or their differentials.

We record precisely the standard extension results used here.  Choose a
smooth projective compactification $B$ with SNC boundary
$\Delta=B\setminus B^\circ$.  The local boundary monodromies of the
integral polarized variation are quasi-unipotent, by the monodromy theorem
\cite[(6.1), pp.~245--246]{Schmid}.  They become unipotent on a finite normal level.  One
can choose the kernel of reduction of the lattice modulo an integer
$M\geq3$: if a quasi-unipotent integral matrix $T$ is congruent to $1$
modulo $M$, then for every eigenvalue $\zeta$ the number
$(\zeta-1)/M$ is an algebraic integer.  All its conjugates have absolute
value at most $2/M<1$, so its algebraic norm forces it to be zero.
Thus every eigenvalue of $T$ is $1$.  The cover exists algebraically by
Riemann existence; normalize its compactification and resolve it.  New
boundary monodromies are products of powers of commuting unipotent
monodromies and remain unipotent.  This level construction may be combined
with a finite level killing a monodromy component group.

At unipotent level, let $\overline{\mathcal V}$ denote Deligne's canonical
extension, with nilpotent logarithmic residues.  The nilpotent orbit
theorem extends $F^\bullet$ by holomorphic subbundles of
$\overline{\mathcal V}$; the induced logarithmic Higgs field is
\[
 \theta:\Gr_F^p\overline{\mathcal V}\longrightarrow
 \Gr_F^{p-1}\overline{\mathcal V}\otimes\Omega^1_B(\log\Delta).
\]
We use the nilpotent orbit and limiting mixed Hodge structure theorems
of \cite[(4.9), (4.12), and (6.16)]{Schmid}.

Here and later, pullback compatibility means compatibility at unipotent
level.  Suppose $g:(B',\Delta')\to(B,\Delta)$ is a morphism of smooth
compactifications whose open part maps to $B^\circ$, and that
$g^{-1}\Delta$ is supported on the SNC divisor $\Delta'$.  Locally write
\[
 g^*z_i=v_i\prod_j t_j^{a_{ij}},\qquad v_i\ \text{a holomorphic unit}.
\]
If $N_i$ are the commuting nilpotent residue matrices, the residues of the
pulled-back logarithmic connection are $\sum_i a_{ij}N_i$ and are
nilpotent.  Uniqueness of the canonical extension therefore identifies it
with $g^*\overline{\mathcal V}$.  The pulled-back Hodge subbundles are
subbundles and agree with the Hodge filtration on the open set, so
uniqueness of the extending Grassmann maps identifies the extended
filtrations as well.  The units $v_i$ contribute only invertible
holomorphic exponential factors.  This proves compatibility also under
the modifications and ramified covers used below.  Tensor operations and
flat Hodge projectors commute with these extensions: the latter commute
with the local monodromies and hence with the logarithmic exponentials.
Thus all the preceding extension assertions apply to the
complex Hodge summands specified above.

The regular-singular Riemann--Hilbert correspondence
\cite[Chapter II, \S\S5--6]{DeligneRSE} makes the flat bundles, their
tensor morphisms, and their pullbacks algebraic on the open set.  On a
projective compactification the extending Hodge subbundles are algebraic
by GAGA; this gives algebraicity on the original open set as well, using
finite descent when a level was introduced.  We also use the theorem of
the fixed part: the monodromy-invariant subspace of a polarized variation
on a smooth quasi-projective base is a constant Hodge substructure, and
its inclusion is a morphism of variations.  This applies to tensor
constructions and finite étale covers
\cite[\S4.1]{DeligneHodgeII} and \cite[(7.22)]{Schmid}.
The assertions about connected algebraic
monodromy needed for the marked-period construction will be stated at
their use.

\subsection{The vector and the coefficient statement}
We now specify regularity in the logarithmic cotangent lattice on a
higher model.
Let $\beta:Y'\to Y$ be a projective birational morphism with $Y'$ smooth,
and let $D'$ be the reduced union of the strict transform of $D$ and all
$\beta$-exceptional divisors, with simple normal crossings.  Set
\[
 L_{Y'}=K_{Y'}+D',\qquad B'=\Omega^1_{Y'}(\log D').
\]
All rational identifications of canonical lines use the canonical
pullback of rational differential forms.  At the generic point of a
prime $A\subset Y'$, choose local frames of $\OO_{Y'}(-aL_{Y'})$ and
$(B')^{\otimes j}$.  After a sufficiently divisible transverse
ramification, a variation with quasi-unipotent monodromy has unipotent
monodromy.  A tensor has \emph{regular coefficients} if, in the
Schmid extension of its Hodge factor, its coefficients relative to the
\emph{pulled-back chosen frames} are regular.  Thus this convention
pulls back the vector bundle $(B')^{\otimes j}$; it does not replace
that bundle by cotangent tensors on the testing curve.

\begin{proposition}[Logarithmic Hodge input]\label{prop:logarithmic-vector}
There are a dense smooth open $Y^\circ\subset V$, a polarizable pure
rational variation $\mathbb V$ on $Y^\circ$ with an invariant lattice
and quasi-unipotent boundary monodromy, and a nonzero algebraic map
\begin{equation}\label{eq:log-input-vector}
 u:\OO_{Y^\circ}(-L_Y)\longrightarrow
 \mathcal E^d,
 \qquad
 \mathcal E^p=\Gr_F^p(\mathbb V\otimes\OO_{Y^\circ}),
 \qquad F^{d+1}\mathbb V=0.
\end{equation}
On every model $(Y',D')$ just described, the rationally identified map
and all its Higgs iterates
\begin{equation}\label{eq:log-input-iterates}
 \theta^j(u):\OO_{Y'}(-L_{Y'})\dashrightarrow
 \mathcal E^{d-j}\otimes(B')^{\otimes j}
\end{equation}
have regular coefficients at every prime of $Y'$ in the preceding
sense.  Here the variation and its Hodge bundles are understood over
the common dense open, and a grade outside their range is zero.

The same statement holds for the homogeneous tensor replacement used
in \Cref{prop:marked-period}.  More precisely, if $u_*$ is obtained
from $u$ by homogeneous tensor operations of positive degree $a$ and
flat Hodge-compatible projections, then
\begin{equation}\label{eq:log-input-tensor}
 u_*:\OO(-aL_{Y'})\dashrightarrow E_0,
 \qquad
 \theta^j(u_*):\OO(-aL_{Y'})\dashrightarrow
 E_j\otimes(B')^{\otimes j}
\end{equation}
have the same coefficient property. Here $\mathbb A$ is the resulting
pure variation, $p_*$ is its highest nonzero Hodge index, and
$E_j=\Gr_F^{p_*-j}\mathbb A$. Finite levels and a common
further unipotent ramification may be used in these assertions.
\end{proposition}

\begin{proof}
We give the construction and then prove the assertion for the full
coefficient lattice.

\smallskip\noindent\emph{Step 1: the cyclic root and the pure vector.}
The cyclic-cover construction and its Higgs iteration follow the
Viehweg--Zuo method \cite[\S3]{ViehwegZuoIsotriviality}; here we must also
retain every coefficient in the original base lattice.
Put $N=K_X+E$. Normalize the cyclic cover defined by an $m$-th root
of $s$, and resolve it to obtain a projective generically finite map
$\psi:Z\to X$.  The source may be disconnected.  The tautological
root gives a regular pairing
\begin{equation}\label{eq:log-root-pairing}
 \psi^*\OO_X(-N)\longrightarrow\OO_Z.
\end{equation}
Choose a reduced simple-normal-crossing divisor on $Z$ containing
the inverse image of $E$ and the exceptional divisors needed for the
chosen compactification.  Shrink the base to $Y^\circ\subset V$ so
that $Z\to Y$ and every stratum of this divisor are smooth there.
Write $H$ for its restriction over $Y^\circ$ and
$q:Z\setminus H\to Y^\circ$ for the resulting open family.  Properness
of the compactification and smoothness of all its strata give local
topological triviality of this family.

The logarithmic mixed Hodge construction gives a variation on
$R^d q_*\Q$ with weight filtration $W$ and Hodge filtration $F$.
We recall the precise ingredients used here.  For a smooth projective
compactification with a simple-normal-crossing boundary, residues
identify the weight quotients of the logarithmic complex with
complexes on the closed boundary intersections, shifted and Tate
twisted.  The associated weight spectral sequence degenerates at
$E_2$, the Hodge-to-de Rham sequence degenerates at $E_1$, and the
same residue filtration on a fixed form sheaf degenerates at $E_2$;
see \cite[\S3.1.5, Theorem~3.2.5 and Corollary~3.2.13]{DeligneHodgeII}.
These statements apply relatively on $Y^\circ$: cohomology and base
change give the Hodge bundles, and the proper smooth boundary strata
give the pure Gauss--Manin systems occurring on the residue page.
Their Gysin maps, including the Tate twists, are morphisms of pure
variations.  Their kernels and images split by polarization.  Thus
the weight quotients are polarizable rational variations, and their
integral cohomological constructions give invariant lattices after
discarding torsion.  Their monodromy is quasi-unipotent.  The proper
stratum Gauss--Manin connections are algebraic and regular singular
by the regularity theorem for Gauss--Manin connections
\cite[Chapter~II, \S6.14 and Theorem~7.9]{DeligneRSE}; residues and their cohomological maps
are algebraic.  Consequently all the Hodge and weight maps used
below have their usual algebraic structures.

Over $Y^\circ$ there is a canonical identification
\[
 \Omega^d_{X/Y}(\log E)\otimes\OO_X(-N)
       =f^*\OO_Y(-L_Y).
\]
Pulling back the top relative logarithmic form and applying
\eqref{eq:log-root-pairing} therefore gives
\begin{equation}\label{eq:log-input-mixed-vector}
 u^{\mathrm{mix}}:\OO_{Y^\circ}(-L_Y)
   \longrightarrow F^d(R^d q_*\C\otimes\OO_{Y^\circ}).
\end{equation}
The form is nonzero: the root is nonzero generically, and pullback
of a top relative differential under a generically finite map in
characteristic zero is injective at the generic point.  A nonzero
global top logarithmic form represents a nonzero top Hodge class
by Hodge-to-de Rham degeneration.  Choose the smallest $w$ for which
the image of \eqref{eq:log-input-mixed-vector} lies generically in
$W_w$.  Projection to $\Gr^W_w$ is nonzero.  Let $\mathbb V=\Gr^W_w$
and let $u$ be this projection.  Since the fibers of $q$ have
dimension $d$, $F^{d+1}=0$.  Thus $u$ is in the highest nonzero
Hodge grade of this pure variation, even when it does not span that
grade.

For use with this construction, the Higgs field has the following
description.  Filter absolute logarithmic forms by the number of
base differentials.  The short exact sequence consisting of base
degrees zero and one induces on relative de Rham cohomology the
Gauss--Manin connection.  This follows either from the logarithmic
de Rham comparison or by lifting a flat class to an absolute class
on a topologically trivial base disk.  Taking its Hodge symbol
gives the connecting map for the corresponding short exact sequence
of individual form sheaves.  Hodge-to-de Rham degeneration identifies
the resulting cohomology sheaves with the Hodge grades.  Since $W$
is flat and all the filtrations are strict, these connecting maps
preserve weight and induce the pure Higgs field on $\Gr^W_w$.

\smallskip\noindent\emph{Step 2: coefficient maps on higher models.}
We now work on a higher model $Y'$.  Resolve the main component of
$X\times_Y Y'$ to obtain $f':X'\to Y'$, and let $E'$ be the strict
transform of $E$ together with the reduced divisors exceptional
over $X$.  Resolve the pair, preserving the smooth relative pair
over $Y'\setminus D'$.  Then
\begin{equation}\label{eq:log-input-support}
 \Supp f'^*D'\subset\Supp E'.
\end{equation}
Indeed, over $D$ this follows from the original support inclusion.
Outside $D$, the image in $Y$ of a $\beta$-exceptional divisor has
codimension at least two.  Smoothness of $f$ there implies that its
inverse image has codimension at least two in $X$.  Every divisor
above it is consequently exceptional over $X$ and belongs to $E'$.
The logarithmic discrepancy formula gives a pulled-back section
\[
 s'\in H^0(X',\OO_{X'}(mN')),
 \qquad N'=K_{X'}+E',
\]
whose rational identification with $s$ is fixed by pullback.
We may dominate both cyclic-root diagrams on a common resolution,
preserving the chosen family over a common dense open.

Fix a prime $A\subset Y'$.  On the inverse image of a sufficiently
small Zariski neighborhood of its generic point, the natural map
\begin{equation}\label{eq:log-input-subbundle}
 f'^*B'\longrightarrow\Omega^1_{X'}(\log E')
\end{equation}
is an injection of subbundles.  If $A\not\subset D'$, this is the
relative smooth-pair condition.  For a component of $D'$, take its
parameter $t$ and tangential base coordinates $y_2,\ldots,y_b$.
At a point over a general point of $A$, write
$f'^*t=v\prod z_\alpha^{n_\alpha}$ in boundary coordinates, with
$v$ a unit and at least one $n_\alpha>0$.  Its logarithmic
differential has the nonzero residue vector $(n_\alpha)$.
On each boundary stratum dominating $A$, generic smoothness gives
independent pullbacks of $dy_2,\ldots,dy_b$ in the ordinary
tangential directions.  Strata not dominating $A$, and the
nonsmooth loci on the dominating strata, have images missing a
dense open of $A$; properness permits their removal on the base.
The residue vector and these tangential differentials are linearly
independent.  This proves \eqref{eq:log-input-subbundle} at every
point over that open of $A$, without any smoothness assumption
over the original boundary.

Let $\Omega^i_{\mathrm{rel}}$ denote the exterior powers of the
locally free quotient in \eqref{eq:log-input-subbundle}.  Its top
determinant gives
\begin{equation}\label{eq:log-input-top-identity}
 \Omega^d_{\mathrm{rel}}\otimes\OO_{X'}(-N')
       =f'^*\OO_{Y'}(-L_{Y'}).
\end{equation}
For $i\ge1$, the exterior-power filtration gives the exact sequence
\begin{equation}\label{eq:log-input-absolute-extension}
 \begin{split}
 0\longrightarrow f'^*B'\otimes\Omega^{i-1}_{\mathrm{rel}}
  \longrightarrow
  \Omega^i_{X'}(\log E')/\mathrm{Fil}_{\mathrm{base}}^2
  \longrightarrow\Omega^i_{\mathrm{rel}}
  \longrightarrow0.
 \end{split}
\end{equation}
The following connecting-map construction is the logarithmic cotangent
procedure of Kov\'acs \cite[\S1, Lemma~1.1]{KovacsFineModuli}, applied
before passing to the root family so that the coefficient lattice remains
fixed. Tensor this sequence by $\OO_{X'}(-N')$, apply higher direct
images, and use the projection formula.  Starting with the unit
map into the direct image of \eqref{eq:log-input-top-identity},
its successive connecting maps give honest morphisms of coherent
sheaves
\begin{equation}\label{eq:log-input-coherent-iterates}
 \alpha_j:\OO_{Y'}(-L_{Y'})\longrightarrow
 A_j\otimes(B')^{\otimes j},\qquad
 A_j=R^j f'_*(\Omega^{d-j}_{\mathrm{rel}}\otimes\OO_{X'}(-N')).
\end{equation}
They are defined before passing to the root family.  On the good
open, differential pullback and root multiplication map each
short exact sequence \eqref{eq:log-input-absolute-extension} to
the corresponding sequence on that family.  Naturality of
connecting homomorphisms identifies the images of $\alpha_j$
with $(\theta^{\mathrm{mix}})^j(u^{\mathrm{mix}})$.  This argument
uses morphisms of short exact sequences of $\OO$-modules; it
does not require root multiplication to commute with the de Rham
differential or a connection on $\OO(-N')$.

\smallskip\noindent\emph{Step 3: the semistable coefficient test.}
Take a general algebraic curve transverse to $A$ and then a disk
about its intersection point.  Resolve the root diagram over
that disk, keeping its smooth pair over the puncture, and include
the special fiber in the resolved boundary.  One-parameter
semistable reduction, after finite ramification, gives a
projective semistable family with smooth total space
\[
 g:Z_\Delta\longrightarrow\Delta
\]
with reduced simple-normal-crossing special fiber $V_0$, and a
horizontal divisor $H$ such that $V_0+H$ has simple normal
crossings.  We use the toroidal construction with the horizontal
divisor included in the boundary
\cite[Chapter~IV, \S3]{KKMS}; an explicit statement allowing
the prescribed boundary is \cite[Theorem~7.17]{KollarMori}.
After resolving the special fiber, its parameter is monomial;
the resolution centers over the bad fiber, followed by base
change and toroidal subdivisions for the vertical
monomial preserve the transverse horizontal coordinate
hyperplanes.  Thus the local equations on the resulting model
are
\begin{equation}\label{eq:log-input-semistable-coordinates}
 \tau=z_1\cdots z_r,
 \qquad H:\ z_{r+1}\cdots z_{r+h}=0.
\end{equation}
In particular every closed horizontal intersection is itself a
proper semistable family over $\Delta$.  The model maps to $X'$
and to the chosen root family.  Its boundary contains the
inverse image of $E'$.

Let $i:\Delta\to Y'$ denote the ramified testing map.  Pullback
of forms, followed by the root pairing, gives
\[
 p^*(\Omega^{d-j}_{\mathrm{rel}}\otimes\OO_{X'}(-N'))
 \longrightarrow
 \Omega^{d-j}_{Z_\Delta/\Delta}(\log(V_0+H)),
\]
where $p:Z_\Delta\to X'$ is the diagram map.  This descends to
relative forms because every pulled-back base logarithmic form
from $Y'$ is a form from $\Delta$ and vanishes in the relative
quotient.  Its pole at $0$ is at most logarithmic.  The root
pairing is regular on this model; alternatively, its monic
power equation and normality prove this regularity.
The natural coherent-cohomology pullback maps therefore give
\begin{equation}\label{eq:log-input-coefficient-test}
 i^*A_j\longrightarrow
 R^j g_*\Omega^{d-j}_{Z_\Delta/\Delta}(\log(V_0+H)).
\end{equation}
No cohomology base-change isomorphism at $0$ is needed.

To see exactly what \eqref{eq:log-input-coefficient-test} tests,
choose a local source generator $e$ and an entire local frame
$(b_\nu)_\nu$ of $(B')^{\otimes j}$, and write
\[
 \alpha_j(e)=\sum_\nu a_\nu\otimes b_\nu.
\]
Every $i^*a_\nu$ maps to a regular relative log-form cohomology
class on $\Delta$.  We keep each $i^*b_\nu$ as a basis vector
of the external bundle $i^*(B')^{\otimes j}$.  On $\Delta^*$,
smooth-family comparison identifies these classes with the
individual coefficients of the original mixed Higgs iterate.
Each coefficient belongs to $W_w$ there, because the mixed
Higgs field preserves $W_w$.

\Cref{lem:weight-extension} shows that its weight projection is
a regular section of the Schmid-extended pure Hodge grade.
This proves the required bound for every coefficient of
\eqref{eq:log-input-iterates}.  In particular, a tangential
factor annihilated by the differential of $i$ is still retained
as an external frame vector.  At an additional root-cover
discriminant outside $D'$, the transverse frame is $dt$.
Although $d(\tau^e)=e\tau^e d\log\tau$, no coefficient is
divided by $e\tau^e$, since this differential substitution is
never made on the external frame.  This proves the bound with
the original $D'$, without adding that discriminant to it.

For completeness, these tests detect divisorial regularity.
Take a finite level on which the variation is unipotent and a
smooth compactification with normal-crossing boundary.  Its
canonical Hodge extensions are algebraic, and the rational
coefficient maps are meromorphic in their frames.  A pole at
a prime persists on a general transverse disk; the preceding
argument excludes it after divisible ramification.  Further
ramified tests preserve regularity by compatibility of the
unipotent canonical extension with pullback.  Since $Y'$ and
$A$ were arbitrary, this proves the assertion on every stated
higher model.

\smallskip\noindent\emph{Step 4: homogeneous replacements.}
Finally, the tensor product Higgs field satisfies the Leibniz
rule.  Each coefficient of $\theta^j(u^{\otimes a})$ is a finite
sum of tensor products of coefficients of iterates of $u$,
with total iteration degree $j$.  Tensor products and flat
Hodge-compatible projections commute with unipotent canonical
extension.  The same applies to the conjugate factors in a
finite-orbit tensor construction on a common finite level.
Thus every homogeneous tensor construction used in
\Cref{prop:marked-period} preserves the asserted regularity,
with source $\OO(-aL_{Y'})$.  Those constructions use polynomial
tensor operations and projections, so no division introduces
a pole.  This proves \eqref{eq:log-input-tensor}.
\end{proof}

The preceding construction reduces regularity to one precise extension
fact. It is important that the weight steps extend as subbundles: this
allows a scalar zero to be divided out before taking its pure projection.

\begin{lemma}[Extension of a weight projection]\label{lem:weight-extension}
Suppose $g:Z\to\Delta$ is projective and has the local form
\eqref{eq:log-input-semistable-coordinates}, with $Z$ smooth,
reduced special fiber $V_0$, and horizontal boundary $H$.
For fixed $i,j$, the relative log-form cohomology
\[
 \mathcal A^{i,j}=R^j g_*\Omega^i_{Z/\Delta}(\log(V_0+H))
\]
is locally free and carries a filtration by subbundles extending
the weight filtration on its Hodge grade over $\Delta^*$.
Its successive quotients are the Hodge-graded unipotent
canonical extensions of the corresponding pure weight
variations.  In particular, a regular class that belongs
generically to a weight step projects regularly to that
step's pure weight quotient.
\end{lemma}

\begin{proof}
Filter $\Omega^i_{Z/\Delta}(\log(V_0+H))$ by the number $k$ of
horizontal logarithmic factors.  In
\eqref{eq:log-input-semistable-coordinates}, taking a horizontal
residue commutes with quotienting by
$d\log\tau=\sum_{a=1}^r d\log z_a$.  Therefore the residue
quotient of filtration degree $k$ is
\begin{equation}\label{eq:log-input-relative-residues}
 \bigoplus_{|I|=k}
 (\iota_I)_*\Omega^{i-k}_{H_I/\Delta}(\log(V_0|_{H_I})),
 \qquad H_I=\bigcap_{\alpha\in I}H_\alpha.
\end{equation}
The terms are zero when their form degree is negative.
The maps $g_I:H_I\to\Delta$ are proper projective semistable
families.  The semistable logarithmic de Rham comparison
\cite[\S2, Theorem~2.7, Corollaries~2.9--2.10 and
Theorem~2.11]{Steenbrink77}, in the canonical-extension form explained
in \cite[\S4.9 and the proof of Lemma~4.10, Step~1]{FF14}, with its
strictness argument completed in \cite{FF14Remark}, identifies
\[
 R^q(g_I)_*\Omega^p_{H_I/\Delta}(\log(V_0|_{H_I}))
 \quad\text{with}\quad
 \Gr_F^p\overline{R^{p+q}(g_I|_{\Delta^*})_*\C},
\]
where the bar means unipotent canonical extension with its
extended Hodge filtration.  In particular these are locally
free sheaves.

Consider the spectral sequence of the finite residue filtration
of this single relative form sheaf.  We use the indexing
\[
 E_1^{-k,j+k}=
 \bigoplus_{|I|=k}R^j(g_I)_*
 \Omega^{i-k}_{H_I/\Delta}(\log(V_0|_{H_I})).
\]
On $\Delta^*$, $d_1$ is the Hodge grade of the signed Gysin
map from the $k$-fold intersections to the $(k-1)$-fold
intersections, with the Tate twists included.  Source and
target have the same pure weight $i+j+k$.  A morphism of
polarizable pure variations has a flat Hodge kernel and image
and admits flat Hodge complements.  Consequently it factors
as a projection to a complement, an isomorphism onto its image,
and an inclusion.  Canonical extension and extension of $F$
preserve this direct-sum factorization.  Its extended
Hodge-graded map has constant rank at $0$.

The coherent $d_1$ and that extended Gysin map are holomorphic
maps between the same locally free $E_1$ terms and agree on
$\Delta^*$.  They therefore agree everywhere.  It follows that
$E_2$ is locally free and is precisely the canonical
Hodge-graded extension of the pure residue cohomology.
By the fixed-form version of Deligne's residue theorem
\cite[Corollary~3.2.13(iii)]{DeligneHodgeII}, every $d_r$ for
$r\ge2$ vanishes on $\Delta^*$.  Inductively its target is
locally free on $\Delta$, so a holomorphic map vanishing on
the puncture vanishes everywhere.  Thus $E_2=E_\infty$ on
the disk.  This is also the constant-rank extension mechanism
in \cite[Proposition~3.12 and the proof of Lemma~4.10,
Step~1]{FF14}; here the residue calculation
\eqref{eq:log-input-relative-residues} supplies its particular
geometric complex.

The abutment has a finite filtration with these locally free
successive quotients.  It and every filtration step are
therefore locally free, and the filtration steps are
subbundles.  On $\Delta^*$ this filtration, with the usual
cohomological shift, is the induced weight filtration on
$\Gr_F^i H^{i+j}(Z_t\setminus H_t)$: this is exactly the
compatibility of the two filtrations in the logarithmic mixed
Hodge complex, or equivalently the fixed-form assertion in
the cited corollary.  Its extended quotients were identified
on $E_2$ above.  Finally, the image of a regular class in the
quotient by a weight subbundle is a section of a locally free
sheaf.  If it is generically zero it is zero.  The class hence
lies in that subbundle on the entire disk and has a regular
image in its successive quotient, as claimed.
\end{proof}

We now retain a zero on a specified original fiber. The source line in
this calculation is the pulled-back original line, whereas the preceding
higher-model estimates used the new logarithmic canonical line. Keeping
these frames distinct is what detects vanishing at the chosen point.

\begin{lemma}[Zero test in the original source frame]\label{lem:zero-test}
Suppose $\dim Y>0$ and $y\in V$, and suppose the section $s$ in
\eqref{eq:log-input-section} vanishes identically on $f^{-1}(y)$.
Blow up $y$ on the base if $\dim Y>1$, and let $A$ be its
exceptional divisor; if $\dim Y=1$, put $A=\{y\}$.
Along a general transverse test at $A$, after ramification
of order $e$ divisible by $m$, the coefficients of $u$ in
the pulled-back \emph{original} $\OO_Y(-L_Y)$ frame vanish
to order at least $e/m$.  Every homogeneous replacement
$u_*$ of degree $a>0$ vanishes to order at least $ae/m$
in the corresponding original $\OO_Y(-aL_Y)$ frame.
\end{lemma}

\begin{proof}
Near $y$ the original morphism is smooth and its boundary is
relative simple normal crossing.  Its fiber is reduced.
Consequently, in any local frame of $\OO_X(m(K_X+E))$, the
coefficient of $s$ belongs to
$\mathfrak m_y\OO_X$.  On the blowup the ideal
$\mathfrak m_y$ becomes the invertible ideal of $A$.
Along a transverse parameter $t$ at its general point the
pulled coefficient is therefore divisible by $t$.

Use throughout the pulled-back old line $\OO_X(K_X+E)$ and
the old relative logarithmic forms.  After $t=\tau^e$, the
tautological root $r$ satisfies, in this old line frame,
\begin{equation}\label{eq:log-input-zero-equation}
 \left(r/\tau^{e/m}\right)^m=s/\tau^e.
\end{equation}
The right-hand side is regular.  The left-hand root is a
rational section integral over the normal root model and
is therefore regular there.  It stays regular on the
semistable model used above.  Thus the map from old relative
top logarithmic forms into the semistable relative forms
has the factor $\tau^{e/m}$.

Passing to coherent cohomology retains that factor.  Divide
it out first: the resulting regular class still lies
generically in $W_w$, since multiplication by a nonzero
scalar on the puncture does not change membership in a
weight step.  By \Cref{lem:weight-extension}, its projection
to $\Gr^W_w$ is regular.  Multiplying the projected class
back by $\tau^{e/m}$ proves the claimed order for $u$.
Every degree-$a$ homogeneous tensor operation multiplies
the order by $a$, and every flat Hodge-compatible projection
preserves this divisibility.  The result follows for $u_*$.
This calculation never replaces the old pulled-back source
line by $\OO(-L_{Y'})$ on the blown-up base; the asserted
positive order is in the original source frame.
\end{proof}

\section{Marked periods of a highest Hodge vector}
\label{sec:marked-periods}

The vector supplied by the covering construction need not generate the
highest Hodge space.  We therefore retain the whole highest space and add
flags determined by the vector on certain compact factors.  This section
constructs the quotient and its line bundles.  The boundary and Higgs-tail
arguments in \Cref{sec:boundary-rank} complete the positivity statement.

\begin{proposition}[Marked-period construction]\label{prop:marked-period}
Let $B^\circ$ be a smooth irreducible complex quasi-projective variety,
and let $\mathbb V$ be a polarizable rational pure variation with an
invariant integral lattice.  Suppose that a line bundle $\mathcal L$ has
a generically nonzero algebraic morphism
\[
 u:\mathcal L\longrightarrow F^p\mathbb V_{\C},
 \qquad F^{p+1}\mathbb V_{\C}=0.
\]
There are a positive integer $a$, a homogeneous tensor construction
followed by flat Hodge-compatible projections, and a generically nonzero
morphism
\begin{equation}\label{eq:marked-vector}
 u_*:\mathcal L^{\otimes a}\longrightarrow E_0,
 \qquad E_0=F^{p_*}\mathbb A,
 \qquad F^{p_*+1}\mathbb A=0,
\end{equation}
where $\mathbb A$ is a polarized complex pure variation.  The operations
may first be performed equivariantly at finite monodromy level.  They
descend to $B^\circ$ after shrinking it.  On smooth projective
birational models there is a dominant morphism
$b:B\to S$ with connected fibers and the following properties.
\begin{enumerate}
\item \emph{Descent.} On dense smooth opens, $\mathbb A$ is the pullback of an actual
variation on $S^\circ$.  Its connected representation is irreducible
and factors through a rational semisimple adjoint group $G$.  The full
monodromy acts through its projected action on
$\mathfrak g=\Lie G$.  In particular, finite scalar monodromy in the
kernel of that action has been eliminated.
\item \emph{The marks.} There is an algebraic compact-flag section $j_{\mathrm u}$,
possibly empty, determined pointwise and equivariantly by $[u_*]$.
In flat coordinates the joint marks
\[
 j=(E_0,j_{\mathrm u})
\]
take values in a closed product flag orbit $J$ of $G_{\C}$.  The two
kinds of marks use disjoint complex simple factors.  Every rational
simple factor of $G$ has a complex factor acting nontrivially on $J$.
The full Lie-period map together with $j_{\mathrm u}$ is generically
immersive on $S^\circ$.
\item \emph{The nef lines and adjoint positivity.} If $S$ is a point, $\mathbb A$ is constant and we set
$\mathcal H=\mathcal H_{\mathrm u}=0$.  Otherwise $S$ has
nef rational line bundles $\mathcal H$ and $\mathcal H_{\mathrm u}$,
with $\mathcal H-\mathcal H_{\mathrm u}$ nef.  On a smooth unipotent
connected-monodromy cover their pullbacks are respectively
\[
 \det E_0+\mathcal H_{\mathrm u}
 \quad\hbox{and the positive compact-flag line.}
\]
The generic curvature rank of $\mathcal H$ is the generic rank of
$dj$, and equals $\nu(\mathcal H)$.  If $D_0$ denotes the reduced
union of boundary components with nonidentity local monodromy on
$\mathfrak g$, then
\begin{equation}\label{eq:marked-initial-big}
 K_S+D_0+c\mathcal H_{\mathrm u}\quad\text{is big for }c\gg0.
\end{equation}
Furthermore,
\begin{equation}\label{eq:marked-final-big}
 K_S+b_0\mathcal H\quad\text{is big for }b_0\gg0.
\end{equation}
\end{enumerate}
All assertions about extensions are compatible with further smooth
projective modifications and pullback at unipotent level.  In
particular, the iterated Higgs fields of $u_*$ are polynomial tensor
expressions in those of $u$; the construction introduces no division
by a local function.
\end{proposition}

We prove the construction and \eqref{eq:marked-initial-big} below.
Equation \eqref{eq:marked-final-big} is
\Cref{prop:adjoint-bigness}.  The other comparison needed later concerns
a chosen line in a pullback of $E_0$: \Cref{prop:tail-rank} constructs
a nef line $P$ from its iterated Higgs images and proves that adding
the pulled-back $\mathcal H$ does not increase $\nu(P)$, provided the
chosen line induces the same compact marks.  Both comparisons are
proved in the next section, after the quotient is constructed.

\subsection{Monodromy, representations, and homogeneous replacements}

Write $M$ for the full rational algebraic monodromy group of
$\mathbb V$, and $G_c=M^\circ$.  A connected monodromy level means a
connected finite étale cover on which the monodromy lies in $G_c$;
its image is still Zariski dense there.

\begin{lemma}\label{lem:marked-monodromy}
The group $G_c$ is semisimple.  At connected monodromy level the Hodge
circles normalize $G_c$, are locally conjugate by $G_c(\mathbb R)^+$,
and have locally constant action on every $G_c$-invariant tensor
space.  The horizontal period differential therefore has the form
\begin{equation}\label{eq:marked-period-tangent}
 t(v)\in\mathfrak g_{c,\C}^{-1,1},
 \qquad \theta(v)=d\rho\bigl(t(v)\bigr),
 \qquad [t(v),t(w)]=0.
\end{equation}
The Lie algebra, and every monodromy-stable sum of its rational
simple ideals, is a polarized rational integral variation of weight
zero.  Its Hodge metric makes $x\mapsto-x^*$ a compact conjugation,
preserving each complex simple ideal.
\end{lemma}

\begin{proof}
André's monodromy normality theorem
\cite[\S5, Theorem~1]{Andre92} applies to this polarizable pure
variation over a smooth connected algebraic variety: the connected
monodromy is normal in the derived generic Mumford--Tate group.
It is therefore semisimple.  The same statement applies after a
finite étale cover and after restriction to a smooth algebraic
subvariety, with its own generic Mumford--Tate group.

The theorem of the fixed part for polarized variations says that
the invariant part of each tensor variation is a constant Hodge
structure; see \cite[Theorem~7.22 and Corollary~7.23]{Schmid}.
These tensor invariants determine
the reductive group $G_c$.  To see the latter assertion in the form
needed here, a $G_c$-stable subspace in a tensor representation has
a $G_c$-equivariant complementary subspace, so it can be encoded by
its invariant projector.  The tensor-stabilizer description of an
algebraic subgroup then applies
\cite[Theorem~9.2 and \S\S9.6--9.7]{Milne}.  Thus, in a fixed flat frame, the
Hodge circles $h_x$ normalize $G_c$, and
\[
 h_x(z)h_{x_0}(z)^{-1}\in G_c(\mathbb R).
\]
In particular their images in the quotient by $G_c$ are constant.
For completeness, local conjugacy follows directly from compactness
of the circle.  Along a smooth path of these homomorphisms, the
logarithmic derivative is a Lie-algebra-valued cocycle for the
circle action.  Averaging a cocycle over the circle writes it as a
coboundary.  Integrating the resulting time-dependent element of
$\Lie G_c(\mathbb R)$ conjugates the circles along the path.  One
can do this first in the adjoint group and lift through the central
isogeny; the remaining central discrepancy is locally constant.

The tangent to the filtration orbit is obtained from the negative
degrees in $\mathfrak g_{c,\C}$.  Griffiths transversality makes its
only possible degree equal to $-1$, giving $t(v)$ and its indicated
action in every representation.  Higgs integrability gives the
commutation relation in \eqref{eq:marked-period-tangent}.

The Lie algebra sits as a rational flat subbundle of
$\End\mathbb V$.  Normalization by the Hodge circles makes it a
Hodge subvariation, and restriction of the endomorphism polarization
polarizes it.  Intersecting its rational fiber with the endomorphism
lattice supplies an invariant full lattice.  These assertions also
hold for the stated rational sums of ideals.  Finally, preservation
of the polarization and normalization by the Hodge grading imply
stability under Hodge-metric adjoint.  On the real Lie algebra the
negative adjoint is the Cartan involution obtained from the Weil
operator.  Its complex antilinear extension is the claimed compact
conjugation.  A compact real form of a semisimple complex algebra
is the direct sum of compact real forms of its simple ideals, so
this conjugation preserves each ideal.
\end{proof}

\begin{lemma}[Homogeneous replacement]\label{lem:marked-replacement}
The replacement \eqref{eq:marked-vector} can be chosen so that its
connected representation is irreducible and its full representation
factors through the action on the retained rational adjoint Lie
factors.  Among all such homogeneous replacements, choose one with
the fewest retained rational simple factors.  Write $G$ for their
product as an adjoint group and $\mathfrak g=\Lie G$.

Every representation factor used in splitting $\mathbb A$ according
to a product decomposition of $G$ is a Hodge-compatible summand of
tensor constructions on the corresponding Lie variation, up to a
constant shift of the Hodge indices.  It can be equipped with the
metric for which representation adjoints are induced by Lie
adjoints.  The highest space $E_0$ is the tensor product of the
highest spaces of these factors and has a closed projective
homogeneous orbit.
\end{lemma}

\begin{proof}
At connected monodromy level decompose the complex representation
into $G_{c,\C}$-isotypical summands.  Each is stable under the Hodge
circle, since that circle acts on the connected group by inner
automorphisms.  An isotypical summand is the tensor product of an
irreducible representation with its multiplicity space.  After a
finite covering of the circle, lift its inner action to the first
factor; the remaining circle action is on the multiplicity space.
Its eigenspace decomposition permits a decomposition into
group-irreducible summands stable under the circle.  The corresponding
flat projectors are Hodge-compatible at a basepoint and, by
\Cref{lem:marked-monodromy}, everywhere.  At least one projector
has generically nonzero value on $u$.

Conjugate this projector by the full monodromy group.  Its orbit is
finite because $G_c$ acts trivially on the projector.  Tensor the
projected vectors over that finite orbit.  On the normal connected
level all factors are generically nonzero: the deck transformations
carry one projected vector to the others.  Permutation of the
factors makes the tensor construction equivariant for full
monodromy, so it descends.  Its source is a positive tensor power
of $\mathcal L$.

For the connected group take the Cartan component, the irreducible
summand of highest weight equal to the sum of the highest weights
of the factors.  It has multiplicity one, and its projection is
preserved by simultaneous automorphisms and permutations of the
factors.  This projection never kills a pure tensor of nonzero
vectors.  Indeed, for a nonzero vector in an irreducible
representation its coefficient along an extremal weight line
after translation by $g\in G_{c,\C}$ is a nonzero regular function
of $g$.  Finitely many such functions are simultaneously nonzero
on a nonempty open set.  Their product is the corresponding
extremal coefficient in the Cartan component.  Since the
projection commutes with translation, its value on the original
tensor was nonzero as well.  The projector preserves the Hodge
grading.  The projected vector is in the highest Hodge degree of
the tensor, which, because its projection is nonzero, is also the
highest degree of the retained Cartan component.

Retain each rational simple adjoint factor having at least one
complex simple ideal acting nontrivially on this irreducible
representation.  This collection is full-monodromy invariant.
Let $q:M\to\operatorname{Aut}(\mathfrak g)$ be the projected
adjoint action, and put $K=\ker q$.  Every element of $K$
commutes with the represented Lie algebra.  Schur's lemma implies
that its action on the irreducible representation is scalar.
Every ideal in $\Lie K$ is unrepresented by the definition of the
retained collection, so $K^\circ$ acts trivially.  Consequently
the image of $K$, although $K$ itself need not be finite, is a
finite group of scalars.  Choose a positive integer killing that
image and take the Cartan component in the corresponding tensor
power.  The resulting representation factors through $q(M)$,
and the preceding nonvanishing argument applies again.  In
particular its connected representation factors through the
adjoint group $G$.

These constructions, as well as subsequent invariant projections,
are homogeneous of positive degree in $u$.  The number of retained
rational factors is a nonnegative integer, so it has a minimum
among these choices.  Fix a minimizing choice.  Its purpose is to
prevent a further nonzero invariant projection from discarding a
factor.  When a factor fixes the highest space throughout its orbit,
this minimality will force the instability that supplies an additional
flag.

Connected irreducibility does not require the highest Hodge step
$E_0$ to have rank one; the construction keeps the entire step.
The adjoint representation is faithful on $G$.  Tensor generation
for a reductive group, \cite[Theorem~4.14]{Milne}, realizes each
of its representations as a summand of finite sums of mixed
tensors of this faithful representation.  For a factor of $G$
the Hodge grading on its Lie algebra is a cocharacter of the
adjoint group.  Its action in an irreducible representation and
the grading inherited from $\mathbb A$ induce the same action on
the Lie algebra; their quotient is scalar.  They thus differ
only by a constant Hodge shift.  Equivariant projectors on the
Lie tensors preserve the grading, and the induced compact-form
metrics have the asserted compatibility with adjoints.  Applying
this separately to the factors gives the tensor decomposition of
$E_0$.  Finally the parabolic associated with the Hodge grading
preserves its maximal-weight subspace.  The stabilizer of $E_0$
therefore contains a parabolic, proving that its orbit is closed
and projective.

All projectors used above are flat Hodge morphisms.  They and all
tensor operations extend on unipotent canonical extensions, by
the functoriality recalled in \Cref{sec:preliminaries}.  The
Leibniz rule for the Higgs field proves the last assertion of
\Cref{prop:marked-period} about iterates.
\end{proof}

The homogeneous replacement has removed the scalar kernel while retaining
a nonzero polynomial expression in the original vector. Full periods
can still miss a direction of that vector on a compact factor. We now
record those directions as additional algebraic flags. They will allow
the joint quotient to retain the vector directions missed by the full
period map; a point quotient will give a constant coefficient system.

\begin{remark}\label{rem:prelim-algebraic-monodromy-torsor}
We will use an algebraic monodromy torsor both for the compact-flag
construction below and for the later horizontal-germ estimate
\Cref{lem:connection-dimension}. Here is its construction.
Restrict a variation to a smooth algebraic base and pass to a finite
étale cover killing the component group of its algebraic monodromy.
Suppose its connected algebraic monodromy $H\subseteq\operatorname{GL}(V)$
is semisimple.  By the stabilizer-of-a-line theorem for algebraic groups
and the tensor realization of representations of $\operatorname{GL}(V)$
\cite[Theorems 4.27 and 4.14]{Milne}, there is a finite tensor construction containing
a line whose stabilizer is exactly $H$.  Since a connected semisimple
group has no nontrivial character, $H$ fixes a nonzero vector $t$ in
that line, and the stabilizer of $t$ is still exactly $H$.  This vector
defines a global monodromy-invariant flat tensor section.  It is
algebraic by the full faithfulness of the regular-singular
Riemann--Hilbert correspondence \cite[Chapter II, \S6]{DeligneRSE}.
Frames carrying $t$ to this section form an algebraic $H$-torsor inside
the frame bundle; on a dense open it is a principal bundle, and its
connection is algebraic and preserves the reduction.  Its monodromy is
Zariski dense in $H$ by definition.  Thus regular singularity is used
to obtain the algebraic reduction, while the dimension estimate itself
only needs the flat algebraic connection in its statement.
\end{remark}


\subsection{Compact flags attached to the vector}

Call a retained rational simple factor \emph{invisible} if all its
complex factors act trivially on the orbit of $E_0$.  This condition
is invariant under the component group of the full monodromy.

\begin{lemma}\label{lem:marked-flags}
For the minimizing replacement of \Cref{lem:marked-replacement}
there is a generically algebraic section $j_{\mathrm u}$ with the
compact-flag and factor-visibility properties in
part~\textup{(ii)} of \Cref{prop:marked-period}.
Its factors have flat positive definite metrics, and their flag
varieties carry positive Plücker lines with invariant metrics.
The construction is pointwise equivariant in $[u_*]$.
\end{lemma}

\begin{proof}
On a nontrivially represented complex ideal in an invisible rational
factor the maximal Hodge subspace is invariant under the whole
simple group.  Irreducibility makes it the whole representation
factor.  The grading on that representation is scalar, and
therefore its induced Lie grading is zero.  The Hodge metric on
this Lie ideal is consequently flat and positive definite.
Monodromy acts isometrically for this compact structure.

Fix an orbit of invisible rational factors under the full
monodromy, and let $U$ be the complexification of their product.
Consider the action of $U$ on a general nonzero value of $u_*$.
If this vector were not unstable for the origin, some homogeneous
invariant polynomial of positive degree would be nonzero on it.
Equivalently, its tensor power would have a nonzero projection to
the $U$-invariant subspace.  This elementary invariant-theoretic
equivalence follows because invariants separate the disjoint
closed sets $\{0\}$ and a closed orbit in the orbit closure;
averaging over a maximal compact subgroup gives the invariant
projection.  The projection is Hodge-compatible: $U$ is normalized
by the Hodge circle.  It is also compatible with full monodromy,
because the chosen product of rational factors is invariant under
it.  After repeating \Cref{lem:marked-replacement}, the nonzero
invariant projection eliminates all factors in this orbit and
introduces none.  This contradicts minimality.  Thus the vector
is generically unstable for each such $U$.

Choose on the cocharacter lattice of $U$ an integral positive
length form invariant under its Weyl group and under the finite
factor automorphisms in question.  One obtains it by summing
positive invariant forms over those automorphisms.  Kempf's
optimal one-parameter subgroup theorem
\cite[\S3]{Kempf78}, in its origin-instability form
\cite[Theorem~6]{Kempf76}, applied to the affine representation
and the closed origin, assigns to an unstable vector a unique
optimal parabolic.  It is equivariant under $U$ and under every
automorphism preserving the representation and the length form.
It is unchanged by multiplying the vector by a nonzero scalar.
It is proper, since $U$ is semisimple and an optimal positive
slope requires a noncentral cocharacter.

We spell out why this pointwise assignment supplies an algebraic
section on a dense open.  Fix a maximal torus.  There are finitely
many subsets of the weight set of the representation.  For each
subset, maximizing the positive normalized minimum pairing is
the closest-point problem for its convex hull, with respect to
the chosen rational positive form.  It determines a rational
optimal ray when the minimum is positive.  Thus there are only
finitely many optimal cocharacter directions and slopes up to
conjugacy.  The conditions that a translate of a vector have a
specified nonzero weight support are constructible.  Their
images, and the conditions selecting the maximal slope, are
constructible by Chevalley's theorem.  The resulting graph of
optimal parabolics on the unstable locus is constructible and single
valued.  Use the algebraic principal monodromy reduction of
\Cref{rem:prelim-algebraic-monodromy-torsor}; its associated projective
representation and flag bundles are algebraic.  Equivariance makes the
graph an algebraic constructible subset of their product over the base.
Pull this graph back along the actual algebraic section $[u_*]$, on the
dense open where that section is unstable.  The resulting constructible
graph has exactly one point over each point of this open base.  Over an
irreducible base in characteristic zero, its dominant component is
birational to the base: after restricting to dense opens, its projection
is a generically injective morphism and therefore has degree one.
It consequently defines a rational section, regular after shrinking.
Applying this argument to the actual section avoids any assumption that
it misses the indeterminacy locus of a rational map on the ambient
projective representation bundle.  The construction uses algebraic associated
bundles, not algebraic local trivializations by holomorphic flat frames.

An optimal cocharacter has zero component on every ineffective
factor, since deleting such a component preserves its slope
numerator and decreases its length.  All nontrivial flags
therefore occur on the isometric complex ideals identified
above.  On the irreducible connected level the generic parabolic
types are fixed.  The component action permutes these types
along with the rational factors.  Properness on at least one
factor in an orbit implies properness on some complex factor of
each rational factor in that orbit.  Repeating this for all
invisible orbits gives $j_{\mathrm u}$.

The highest-space marks already see every other rational factor.
Stabilizers of highest spaces and of the new flags are parabolic,
and the two constructions use disjoint factors.  Thus their
joint orbit is the asserted product $J$.  A tensor product of
highest subspaces determines each factor subspace, so this
description is also valid when $E_0$ is presented as a single
subspace of $\mathbb A$.  Embed a parabolic flag by the top
exterior power of its parabolic subalgebra.  The dual tautological
Plücker line is positive, with the metric induced by the flat
compact metric.  Tensoring these lines gives a positive line
invariant also under the relevant factor permutations.
\end{proof}

The marks now see every retained rational factor. We next construct an
algebraic parameter space and descend the actual coefficient variation
to it. Descent of the filtration alone would leave a possible finite
scalar obstruction, which is why that kernel was removed first.

\Needspace{10\baselineskip}
\subsection{An algebraic quotient and actual descent}

\begin{lemma}[Generic Lie stabilizer]\label{lem:generic-lie-stabilizer}
Let a polarized rational Lie variation on a smooth connected complex
quasi-projective variety have an invariant integral lattice, semisimple
fiber $\mathfrak a$, and connected algebraic monodromy
$\operatorname{Int}(\mathfrak a)$.  At a Hodge-generic lifted
period, a rational Lie automorphism preserving the Hodge
filtration is the identity.
\end{lemma}

\begin{proof}
A rational endomorphism preserving the filtration of a pure Hodge
structure also preserves its conjugate filtration, hence is a
Hodge endomorphism.  It commutes with the Mumford--Tate group.
At a Hodge-generic point this is the generic Mumford--Tate group,
which contains the connected monodromy by
\cite[\S5, Theorem~1]{Andre92}.  If $\alpha$ is also a Lie
automorphism, then for every $x\in\mathfrak a$,
\[
 \operatorname{ad}(\alpha x)
   =\alpha\operatorname{ad}(x)\alpha^{-1}
   =\operatorname{ad}(x).
\]
The Lie algebra has zero center, so $\alpha x=x$.  Such points
may be chosen while imposing finitely many generic differential
rank conditions: in a flat chart the proper loci where any one
of the countably many rational tensors acquires an additional
Hodge condition can be avoided together with those conditions.
\end{proof}

\begin{lemma}[Period quotient and descent]\label{lem:period-descent}
There is an algebraic quotient with connected generic fibers
which remembers the full integral period of $\mathfrak g$ and
the compact marks.  Both $\mathfrak g$ and $\mathbb A$ descend
as variations on a dense smooth open of the quotient.  Its full
Lie period and compact marks give a generically immersive joint
map in flat coordinates.  The same construction is available for
any monodromy-stable collection of rational normal factors, at
connected level if necessary.
\end{lemma}

\begin{proof}
Use the full filtration on the rational integral polarized Lie
variation, not only the highest subspace of $\mathbb A$.  Let
$\Gamma$ be an arithmetic group of isometries of its polarized
lattice containing its monodromy, and let $\Omega$ be its period
domain, or the associated generic Mumford--Tate domain.
The algebraicity theorem for pure integral period maps
\cite[Theorem~1.1]{BBT} factors the period map through a
dominant algebraic map to a quasi-projective algebraic image.
Its hypotheses hold here by \Cref{lem:marked-monodromy}; in
particular a lattice has been retained in the full Lie system.
Normalize this image in the relative algebraic closure of its
function field in $\C(B)$.  The resulting map
$b_0:B^\circ\dashrightarrow T^\circ$ has geometrically connected
generic fiber.  Shrink to smooth opens where it is a morphism.
The dimension of $T^\circ$ is the generic rank of the full
period map.  Indeed, the real period stabilizer is compact,
arithmetic orbits are discrete, and the factorization preserves
the local analytic image dimension.

Here is the local-system descent, including possible finite
monodromy.  Compactify the graph properly over $T^\circ$, resolve
the complement, and shrink $T^\circ$ so that the compactification
and all strata of its SNC complement are smooth over it.
Lifting vector fields tangent to the strata gives a local
differentiable trivialization of this proper smooth pair and
hence of its open part.  We may thus assume that $b_0$ is a
topological fiber bundle with connected fibers.  Removing a
proper algebraic subset of a smooth complex variety induces a
surjection on its fundamental group, so these shrinkings do not
decrease the relevant monodromy images.

Choose a Hodge-generic point of $B^\circ$.  Along the universal
cover of its $b_0$-fiber the lifted Lie-period map has values in
one discrete arithmetic orbit; it is therefore constant.
Every vertical loop acts by a rational Lie automorphism fixing
that lifted period.  \Cref{lem:generic-lie-stabilizer} makes its
action trivial.  The exact sequence
\[
 \pi_1\bigl(b_0^{-1}(t)\bigr)\longrightarrow
 \pi_1(B^\circ)\longrightarrow\pi_1(T^\circ)
 \longrightarrow1
\]
now descends the Lie representation.  The representation on
$\mathbb A$ descends as well: by
\Cref{lem:marked-replacement} it factors through the full
projected Lie action, including its component group.  This is
where removal of the finite scalar image is necessary.

The Lie filtration is constant in flat coordinates along each
connected fiber.  The derivative of the filtration on
$\mathbb A$ is the action of
\eqref{eq:marked-period-tangent}; it too vanishes on those
fibers.  Consequently both filtrations descend.  Holomorphicity,
Griffiths transversality, and polarization can be checked after
the smooth surjective pullback.  The descended local systems
have quasi-unipotent local monodromy: for a boundary valuation
of $T$ choose a valuation above it on a proper model of $B$;
a positive power of its monodromy is a monodromy of the original
variation.  Quasi-unipotence of that power implies
quasi-unipotence of the original operator.  Canonical extension
and GAGA, or regular-singular Riemann--Hilbert, therefore give
the algebraic flat bundles and algebraic Hodge bundles on these
opens; see \cite[Chapter~II]{DeligneRSE} and
\Cref{sec:preliminaries}.

Over $T^\circ$ consider the algebraic flag bundle of the
descended compact factors.  The section $j_{\mathrm u}$ gives
an algebraic map from $B^\circ$ into that bundle.  Take its
image and again normalize in the relative algebraic closure
inside $\C(B)$.  This is the desired $S^\circ$.  Variations
pull back from $T^\circ$, and the compact flag descends by its
tautological definition.  On a dense open the finite map from
$S^\circ$ to its image is étale, and the period-image
factorization preserves tangent rank.  Hence the full period
and compact marks have joint differential of rank $\dim S$.
Resolve projective compactifications and the map from $B$ to
obtain $b:B\to S$ with connected fibers.  The monodromy images
on the open source and target agree by the fundamental-group
sequence just used.

For a monodromy-stable collection of rational factors the
projected Lie system again has its rational polarization and
lattice.  The identical argument gives its quotient and
descent.  At connected level every rational normal factor
collection is permitted.
\end{proof}

\phantomsection\label{proof:point-period}
If $S$ is a point, the full projected period is constant.
Its monodromy lies in a compact period stabilizer and in the
discrete integral isometry group, and is thus finite.  Its
connected algebraic monodromy $G$ is trivial.  The representation
on $\mathbb A$ factors through the projected Lie action, whose
kernel has already been removed, so $\mathbb A$ is constant.
When there are no retained factors the construction simply takes
$S$ to be a point.

\subsection{The metric comparison used for nefness}

The quotient has now been constructed and the coefficient variation
descends to it.  Two different differentials enter the positivity
argument.  The full joint differential $(t,dj_{\mathrm u})$ has generic
rank $\dim S$, by \Cref{lem:period-descent}; it will control the
logarithmic adjoint.  The differential of the marks
$j=(E_0,j_{\mathrm u})$ can have smaller rank.  It is this second
differential that the line $\det E_0+\mathcal H_{\mathrm u}$ measures
on the open set.  We first compute its curvature and then use the
boundary growth estimates to identify curvature rank with numerical
dimension on a projective model.

In the smooth Hodge splitting $E=\bigoplus_p E^p$, write
\[
 \nabla=\mathcal D+\theta+\theta^\dagger,
\]
where $\mathcal D$ is the graded Chern connection and
$\theta^\dagger$ is the Hodge adjoint of the degree-lowering Higgs field.
If $\theta_p(v):E^p\to E^{p-1}$, flatness and the polarization give,
with positive curvature normalization,
\begin{equation}\label{eq:prelim-hodge-curvature}
 R_{E^p}(v,\bar v)=
 \theta_p(v)^*\theta_p(v)
 -\theta_{p+1}(v)\theta_{p+1}(v)^*.
\end{equation}
The irrelevant common positive normalization factor is suppressed here.
For the highest nonzero step $F^p=E^p$, the second term vanishes, and hence
\begin{equation}\label{eq:prelim-top-curvature}
 c_1(\det F^p,h)(v,\bar v)
 =c\,\|\theta_p(v)\|_{\mathrm{HS}}^2\quad(c>0).
\end{equation}
In a local flat frame the differential of the subspace map
$b\mapsto F^p_b$ is exactly $\theta_p$; transversality and $F^{p+1}=0$
ensure that its image lies in $F^{p-1}/F^p$.  Thus the curvature kernel
in \eqref{eq:prelim-top-curvature} is exactly the kernel of that
subspace differential.  This statement concerns the whole highest
step, not a chosen line inside it.  Adding the positive compact-flag
form gives kernel $\ker(dE_0,dj_{\mathrm u})$.

To extend this computation across the boundary, we use the
several-variable norm estimates of \cite[Theorem 5.21]{CKS}, also stated
and proved in \cite[Theorem 5.1]{CK}.  Work at unipotent level with the
canonical Hodge extensions of \Cref{sec:preliminaries}.  In a boundary
chart $\Delta=\{z_1\cdots z_\ell=0\}$, shrink so that $|z_i|<e^{-1}$
and put $\Lambda=\prod_{i=1}^{\ell}(-\log|z_i|)$.  Relative to any
smooth positive metric $h_0$ on the extending bundle there are constants
$C,N>0$ such that its Hodge metric satisfies
\begin{equation}\label{eq:prelim-hodge-growth}
 C^{-1}\Lambda^{-N}h_0\ \leq\ h\ \leq\ C\Lambda^N h_0.
\end{equation}
The same assertion holds for the dual, every extending filtration step,
and every Hodge graded quotient, with their induced metrics.  To pass
from the usual sector estimates to \eqref{eq:prelim-hodge-growth}, divide
the logarithmic variables into their finitely many order sectors.  Each
ratio power in the sector estimate is bounded by a power of $\Lambda$.
The finite logarithmic exponential that converts flat multivalued frames
to canonical frames has entries polynomial in the logarithms.  Applying
the same bounds to the dual gives the lower estimate.  Restriction to a
subbundle and passage to its quotient preserve a two-sided comparison
with a smooth metric, which gives the assertions for $F^p$ and $\Gr_F^p$.
The extension compatibility with flat Hodge projectors proved in
\Cref{sec:preliminaries} gives the same estimates for the complex Hodge
summands used here.

The induced line metrics therefore have two-sided logarithmic-logarithmic
weight bounds.  A curvature computation on the open set alone would not
identify their intersection numbers.  The following full-mass argument
makes that comparison.


\begin{lemma}\label{lem:log-metric}
Let $Z$ be a smooth projective variety of dimension $n$, let $\Delta$ be
an SNC divisor, and let $L$ be a rational line bundle.  Suppose that on
$Z^\circ=Z\setminus\Delta$ it has a smooth Hermitian metric with
semipositive curvature form $\Theta$, normalized to represent $c_1(L)$.
Suppose that in regular nonvanishing local frames of a Cartier multiple
of $L$ the metric weights satisfy
\begin{equation}\label{eq:prelim-loglog-bound}
 |\varphi(z)|\leq C\left(1+
       \sum_{i=1}^{\ell}\log(-\log|z_i|)\right)
\end{equation}
near every boundary chart $\Delta=\{z_1\cdots z_\ell=0\}$.
Then $L$ is nef.  For every Kähler form $\eta$ on $Z$ and
$0\leq k\leq n$,
\begin{equation}\label{eq:prelim-mass-identity}
 \int_{Z^\circ}\Theta^k\wedge\eta^{n-k}
       =c_1(L)^k[\eta]^{n-k}.
\end{equation}
In particular,
\begin{equation}\label{eq:prelim-curvature-rank}
 \nu(L)=\max_{z\in Z^\circ}\rk\Theta_z.
\end{equation}
When the curvature is real analytic, as for the induced Hodge and
subspace metrics used below, this maximum is its generic rank.
\end{lemma}

\begin{proof}
Taking a Cartier multiple and dividing its weights and curvature by that
multiple reduces to a line bundle.  We use $dd^c$ normalized so that the
curvature of a metric of weight $\varphi$ is $dd^c\varphi$.

First we extend its local plurisubharmonic weights.  For $\epsilon>0$,
the function
\[
 \varphi_\epsilon=\varphi+
               \epsilon\sum_{i=1}^{\ell}\log|z_i|
\]
is plurisubharmonic off the coordinate divisor and is locally bounded
above near it, by \eqref{eq:prelim-loglog-bound}.  The removable
singularity theorem extends it plurisubharmonically to the polydisk.
On a smaller polydisk, the maximum principle, applied successively in
the coordinate variables, bounds it above by its values on a fixed
distinguished boundary torus disjoint from $\Delta$.  These bounds are
uniform as $\epsilon\downarrow0$.  Its increasing limit, with upper
semicontinuous regularization, is therefore a plurisubharmonic extension
of $\varphi$.  The extensions agree under the pluriharmonic frame changes
and give a closed positive current $T\in c_1(L)$ on $Z$.

This current has zero Lelong number at every point.  At a boundary point
approach along a path on which all the vanishing boundary coordinates
have absolute values comparable to a common parameter $r\downarrow0$.
The two-sided bound in \eqref{eq:prelim-loglog-bound} is then
$O(\log\log(1/r))=o(|\log r|)$.  The growth characterization of the
Lelong number gives an upper bound of zero; positivity gives the reverse
bound.  Away from $\Delta$ the current is smooth.  Demailly's
regularization theorem, in its zero-Lelong-number consequence
\cite[Corollary 6.4]{Demailly92}, now shows that $c_1(L)$ is nef.

We next prove the mass assertion, since zero Lelong numbers alone would
not suffice for it.  Fix $\epsilon>0$.  Nefness makes
$\alpha_\epsilon=c_1(L)+\epsilon[\eta]$ a Kähler class.  Choose a Kähler
representative $\omega_\epsilon$ and write
\[
 T+\epsilon\eta=\omega_\epsilon+dd^c\phi.
\]
Here $\phi$ is a global $\omega_\epsilon$-plurisubharmonic function,
smooth on $Z^\circ$; changing the smooth background has not changed its
logarithmic growth bound.  Write the irreducible components of $\Delta$
as $\Delta_i$, choose their defining sections $s_i$ with smooth metrics,
and rescale these metrics so that
\[
 u_i=-\log\|s_i\|^2\geq1\quad\text{on }Z^\circ.
\]
For any fixed $0<b<1$ and sufficiently small $\delta>0$, the function
\begin{equation}\label{eq:prelim-mass-barrier}
 \psi=-\delta\sum_i u_i^b
\end{equation}
is $\omega_\epsilon$-plurisubharmonic.  Indeed, on $Z^\circ$ each
$dd^c u_i$ is a fixed smooth form and
\[
 dd^c(-\delta u_i^b)
   =-\delta b u_i^{b-1}dd^c u_i
     +\delta b(1-b)u_i^{b-2}du_i\wedge d^c u_i.
\]
The second summand is positive and the first is bounded below by
$-\delta C\omega_\epsilon$, because $u_i\geq1$.  Decrease $\delta$ so
that the sum of these lower bounds is at least
$-\tfrac12\omega_\epsilon$.  The same extension argument applies across
$\Delta$.  Since $\log u=o(u^b)$, the lower growth bound for $\phi$
also gives a constant $C_\epsilon$ such that
\begin{equation}\label{eq:prelim-less-singular}
 \phi\geq\psi-C_\epsilon.
\end{equation}

We verify directly that $\psi$ has full Monge--Ampère mass.  Take a smooth
decreasing function $\chi:[0,\infty)\to[0,1]$ that is $1$ on $[0,1]$
and $0$ on $[2,\infty)$.  For $R>1$ set
\[
 f_R(t)=1+\int_1^t b s^{b-1}\chi(s/R)\,ds\qquad(t\geq1).
\]
Then $f_R=t^b$ for $t\leq R$, $f_R$ is increasing and concave, and it
is constant for $t\geq2R$.  Its derivative bounds, uniformly in $R$, are
\begin{equation}\label{eq:prelim-barrier-truncation}
 0\leq f_R'(t)\leq C t^{b-1},\qquad
 0\leq-f_R''(t)\leq C t^{b-2}.
\end{equation}
Moreover $f_R\uparrow t^b$ as $R\to\infty$.  The functions
$\psi_R=-\delta\sum_i f_R(u_i)$ extend smoothly over $Z$, are
$\omega_\epsilon$-plurisubharmonic for the same choice of $\delta$,
and equal $\psi$ locally on $Z^\circ$ once $R$ is sufficiently large.

In a fixed boundary coordinate chart, comparison of
$du_i=-d\log|z_i|^2+$ a smooth form and
\eqref{eq:prelim-barrier-truncation} gives a bound of positive forms
\begin{equation}\label{eq:prelim-local-majorant}
 0\leq\omega_\epsilon+dd^c\psi_R
 \leq C\left(\eta_0+
       \sum_{i=1}^{\ell}
       \frac{\sqrt{-1}\,dz_i\wedge d\bar z_i}
       {|z_i|^2(-\log|z_i|)^{2-b}}\right),
\end{equation}
where $\eta_0$ is a smooth positive coordinate form and $C$ is independent
of $R$.  The right-hand side has an integrable $n$th power.  Each
singular summand has rank one and its square is zero, and its normal
integral is bounded by a constant times
\[
 \int_0^{r_0}\frac{dr}{r(-\log r)^{2-b}}<\infty.
\]
Products involving different normal coordinates are integrable by
Fubini.  A finite chart covering and dominated convergence now imply
\begin{align*}
 \int_{Z^\circ}(\omega_\epsilon+dd^c\psi)^n
  &=\lim_{R\to\infty}
        \int_Z(\omega_\epsilon+dd^c\psi_R)^n\\
  &=\int_Z\omega_\epsilon^n.
\end{align*}
The last equality is Stokes' theorem for the smooth approximants.  Since
the non-pluripolar Monge--Ampère product does not charge $\Delta$, this
is exactly full mass for $\psi$.

The full-mass class is preserved on passing to a less singular
plurisubharmonic potential.  Apply
\cite[Proposition 1.6]{GZ} to the Kähler form $\omega_\epsilon$ and
\eqref{eq:prelim-less-singular}.  All the hypotheses are satisfied:
$Z$ is compact Kähler, both potentials are
$\omega_\epsilon$-plurisubharmonic, and the lower potential has just
been shown to have full mass.  It follows that
\begin{equation}\label{eq:prelim-epsilon-mass}
 \int_{Z^\circ}(\Theta+\epsilon\eta)^n
       =(c_1(L)+\epsilon[\eta])^n
       \qquad(\epsilon>0).
\end{equation}
All mixed integrands on the left are nonnegative.  The identity for one
positive $\epsilon$ therefore proves that each has finite integral.
Expansion of \eqref{eq:prelim-epsilon-mass} and comparison of polynomial
coefficients proves \eqref{eq:prelim-mass-identity}.

Finally, a smooth semipositive $(1,1)$-form has
$\Theta^k\wedge\eta^{n-k}>0$ at a point exactly when its rank there is
at least $k$.  If this happens at one point it happens on an open
neighborhood, so its integral is positive.  Combining this observation
with \eqref{eq:prelim-mass-identity} and
\eqref{eq:prelim-numerical-dimension} proves
\eqref{eq:prelim-curvature-rank}.  For a real-analytic form the nonzero
maximal minors cannot vanish on an open set, so the maximal-rank locus
is dense.  This proves the final assertion.  When
$n=0$ the statements follow from the empty-product convention.
\end{proof}

\begin{remark}\label{rem:prelim-metric-operations}
Estimate \eqref{eq:prelim-hodge-growth} supplies
\eqref{eq:prelim-loglog-bound} for determinants of extending Hodge
bundles, their holomorphic subbundles and quotients, and duals and tensor
products of these lines.  For a subspace given initially only on a dense
open set, first resolve its rational Grassmann map; the resulting
tautological subbundle is an actual subbundle of the pulled-back
extending bundle.  Its inclusion is uniformly nondegenerate relative
to smooth metrics, so the two-sided estimate applies.  Multiplication
by a smooth positive metric factor only adds a bounded weight.  The
semipositivity required by \Cref{lem:log-metric} must still be proved for
each line by its curvature calculation.

The rank computation may also be made on a smooth generically finite
cover $q:Z'\to Z$.  If the lemma applies to $q^*L$, it gives nefness of
$q^*L$, hence of $L$, by testing a curve using a curve lying over it.
To compute intersections downstairs, the mass identity upstairs remains
valid with $q^*\eta$ in place of the Kähler form: apply it first to
$q^*\eta+t\eta'$ for $t>0$, where $\eta'$ is Kähler, and compare
coefficients in $t$.  The projection formula gives
\[
 (q^*c_1(L))^k(q^*[\eta])^{n-k}
       =\deg(q)\,c_1(L)^k[\eta]^{n-k}.
\]
The form $q^*\eta$ is positive definite on the étale locus, which is
dense.  A nonempty open maximal-curvature-rank locus meets it.  Thus the
same curvature rank computes $\nu(L)$ downstairs.
\end{remark}


\subsection{The nef line and an initial logarithmic adjoint}

Assume $\dim S>0$.  Choose a normal finite monodromy level of
$S^\circ$ contained in the connected group and in a principal
integral congruence subgroup of sufficiently high level, for
example a level divisible by $3$.  Quasi-unipotent elements at
this level are unipotent.  One elementary way to see the
assertion is that if $A\equiv1\pmod m$, then
$(\zeta-1)/m$ is an algebraic integer for each root-of-unity
eigenvalue $\zeta$.  For $m\ge3$ all its conjugates have
absolute value less than one unless it is zero; hence
$\zeta=1$.

Take an equivariant smooth projective compactification
$\widetilde S$ of this level, with SNC boundary.  On the compact
complex ideals the now-unipotent local monodromy is unitary and
therefore trivial.  Their flat bundles extend across the
boundary with their smooth positive metrics.  Resolve the
compact-flag maps in these bundles equivariantly.  The line
\begin{equation}\label{eq:marked-level-line}
 \widetilde{\mathcal H}
   =\det\widetilde E_0+\widetilde{\mathcal H}_{\mathrm u}
\end{equation}
is defined on this model: $\widetilde E_0$ is the highest
Schmid extension, and
$\widetilde{\mathcal H}_{\mathrm u}$ is the product Plücker
line from \Cref{lem:marked-flags}.  Both have semipositive
curvature.  For $\det\widetilde E_0$ the curvature is the
squared norm of the differential of the highest subspace; for
the compact summand it is the pullback of the positive flag
form.  Their kernels therefore intersect in the kernel of $dj$.
The Hodge norm and dual norm have logarithmic growth in
canonical-extension frames; the compact metric is smooth.
\Cref{lem:log-metric} gives nefness and
\begin{equation}\label{eq:marked-nef-rank}
 \nu(\widetilde{\mathcal H})
   =\max\operatorname{rank}(dj).
\end{equation}

The construction is linearized for the deck group.  A common
positive power kills the finite stabilizer actions on the
fibers of each line, and that power descends to the finite
quotient.  For clarity, descent can be checked locally by
averaging a local frame after its stabilizer acts trivially;
the invariant frame is nonzero after shrinking and supplies
the required quotient frame.  Let $S$ now be a smooth
projective log resolution of that quotient, replacing the
earlier compactification, and pull the descended rational
lines to $S$.  Denote them by $\mathcal H$ and
$\mathcal H_{\mathrm u}$.  Nefness descends under a finite
surjection and is preserved by pullback.  On the resolution
of the corresponding base change upstairs the lines are
exactly the pullbacks of \eqref{eq:marked-level-line}, by
functoriality of unipotent extensions and of the resolved flag
maps.  In particular
\begin{equation}\label{eq:marked-downstairs-rank}
 \nu(\mathcal H)=\max\operatorname{rank}(dj),
 \qquad
 \mathcal H_{\mathrm u}\text{ and }
 \mathcal H-\mathcal H_{\mathrm u}\text{ are nef}.
\end{equation}

We record compatibility with rational normal factors, used in
the next section.  At connected level, write $G=H\times Q$.
Then $\mathbb A=A_H\otimes A_Q$ and
$E_0=E_H\otimes E_Q$.  Splitting the compact marks accordingly
gives
\begin{align}\label{eq:marked-factor-lines}
 \mathcal H_H
   &=(\rk E_Q)\det E_H+\mathcal H_{\mathrm u,H},&
 \mathcal H_Q
   &=(\rk E_H)\det E_Q+\mathcal H_{\mathrm u,Q},&
 \mathcal H&=\mathcal H_H+\mathcal H_Q.
\end{align}
These are nef lines on suitable smooth models.  Apply
\Cref{lem:period-descent} to the full $Q$-period and its
compact marks, obtaining $q:S\to S_Q$ after modification.
The $Q$-line is $q^*\underline{\mathcal H}_Q$ for the same
construction on $S_Q$.  The restricted congruence condition
holds on the lattice intersected with each rational ideal, so
the descended $Q$-system has unipotent boundary monodromy.
Consequently its extensions, flag lines, and these identities
remain compatible on a common resolution.

Mark the reduced SNC union $D_0$ of those boundary prime
divisors of $S$ whose local monodromy on $\mathfrak g$ is
not the identity.

\begin{lemma}[Initial adjoint bigness]\label{lem:initial-adjoint}
For the marked-period quotient just constructed,
\[
 K_S+D_0+c\mathcal H_{\mathrm u}
 \quad\text{is big for all sufficiently large }c.
\]
\end{lemma}

\begin{proof}
Let $d=\dim S$ and work first on a smooth unipotent level
$p:\widetilde S\to S$, resolving further when necessary.
On its open part the joint differential has image
\[
 \mathcal J\subseteq
 \Gr_F^{-1}\mathfrak g_{\C}\oplus
 j_{\mathrm u}^*T_{J_{\mathrm u}}.
\]
The differential here is $(t,dj_{\mathrm u})$, where the
second component is the vertical derivative in flat flag
charts.  It is globally defined because the transition maps
are flat.  Its rank is $d$ by \Cref{lem:period-descent}.
Resolve the map to the corresponding relative Grassmannian
so that $\mathcal J$ extends as a subbundle of the indicated
extended ambient bundle.

Equip it with the subbundle metric.  We give the curvature
calculation because it uses the full period differential,
whereas the nef line $\mathcal H$ only uses the marks $j$.
Fix a tangent vector $v$, and choose an orthonormal frame
$(a_\mu,b_\mu)$ of $\mathcal J$ at the point.  Each
$a_\mu$ lies in the image of $t$, so integrability gives
$[t(v),a_\mu]=0$.  The Hodge curvature formula on the
degree $-1$ Lie space consequently gives
\[
 \sum_\mu
 \bigl\langle\Theta_{\Gr_F^{-1}\mathfrak g}(v,\bar v)
          a_\mu,a_\mu\bigr\rangle
 =-\sum_\mu\bigl\|[t(v)^*,a_\mu]\bigr\|^2.
\]
The convention is the curvature normalization of
\eqref{eq:prelim-hodge-curvature}.  The flag tangent curvature is
bounded above by $C\omega_{\mathrm u}(v,\bar v)$ times
the identity, where $\omega_{\mathrm u}$ is the positive
compact-flag form pulled back along the section.  Such a
uniform $C$ exists because the flag variety is compact and
the flat transition actions are isometric.  Subbundle
curvature subtracts a squared second fundamental form.
After increasing $C$ by the fixed rank $d$, this proves
\begin{equation}\label{eq:marked-curvature-estimate}
 \Theta_{-\det\mathcal J+c\widetilde{\mathcal H}_{\mathrm u}}
       (v,\bar v)
 \geq \sum_\mu\bigl\|[t(v)^*,a_\mu]\bigr\|^2
          +(c-C)\omega_{\mathrm u}(v,\bar v).
\end{equation}
The right side is positive whenever the joint differential
is nonzero.  Indeed, if its compact component is nonzero,
the last term is positive for $c>C$.  Otherwise
$t(v)\ne0$.  Since $(t(v),dj_{\mathrm u}(v))$ belongs
to $\mathcal J$, vanishing of all the commutators in
the first term would imply $[t(v)^*,t(v)]=0$.
But $t(v)$ acts as a nonzero lowering nilpotent in the
faithful adjoint representation.  A normal nilpotent
endomorphism is zero, a contradiction.  Thus the line
\[
 R=-\det\mathcal J+c\widetilde{\mathcal H}_{\mathrm u}
\]
has semipositive curvature and full generic curvature rank.
The extended ambient Hodge metrics and their duals have
logarithmic bounds, the compact metric is smooth, and the
resolved subbundle metric has the same bounds.  By
\Cref{lem:log-metric}, $R$ is nef with numerical dimension
$d$, hence big.

It remains to compare $R$ with the claimed adjoint on $S$.
The differential gives a generically invertible map
\begin{equation}\label{eq:marked-log-tangent-map}
 p^*T_S(-\log D_0)\dashrightarrow\mathcal J.
\end{equation}
It is regular at every prime upstairs which dominates a
prime of $S$.  At a marked prime, the finite cover is
generically a power substitution $t=z^e$.  The lift of
$t\partial_t$ is $e^{-1}z\partial_z$.
The logarithmic Higgs field extends on the unipotent
canonical extension, so its Lie component is regular;
one recovers the Lie-valued $t$ from the adjoint Higgs
field by the split adjoint inclusion into endomorphisms.
The resolved compact flag is regular upstairs, so its
derivative on $z\partial_z$ is regular as well.  The
tangential generators cause no poles.  At an unmarked
prime the monodromy level is generically unramified,
since it was defined using only the projected adjoint
monodromy.  The local Lie monodromy is the identity and
the extended period has no logarithmic term.  The same
regularity assertion follows using the ordinary tangent
generator there.

These regular ambient maps land in $\mathcal J$ at the
primes in question.  Indeed, $\mathcal J$ is a subbundle,
so its quotient in the ambient bundle is torsion-free;
a regular ambient section whose generic value lies in
$\mathcal J$ has zero image in that quotient.  Taking
determinants in \eqref{eq:marked-log-tangent-map} and
dualizing gives a generically nonzero map
\begin{equation}\label{eq:marked-adjoint-comparison}
 R\dashrightarrow p^*(K_S+D_0+c\mathcal H_{\mathrm u})
\end{equation}
regular above every prime of $S$.

Possible poles on divisors exceptional over $S$ do not
affect the conclusion, as follows.  Factor $p$ through
the finite normalization $S'\to S$ of this level.  Choose
an ample line $A_S$ on $S$ and a positive rational
$\epsilon$ sufficiently small that
$R-\epsilon p^*A_S$ is big.  For a sufficiently divisible
positive integer $m$, it has a nonzero section.  Its image
under \eqref{eq:marked-adjoint-comparison} is a rational
section of
\[
 p^*m(K_S+D_0+c\mathcal H_{\mathrm u}-\epsilon A_S)
\]
regular at every divisor of the normal variety $S'$.
Normality extends it over $S'$.  Taking its field norm
to $S$ gives a nonzero section of a positive multiple
of $K_S+D_0+c\mathcal H_{\mathrm u}-\epsilon A_S$.
Adding the ample term proves bigness.  Since
$\mathcal H_{\mathrm u}$ is nef, the conclusion persists
when $c$ is increased.
\end{proof}

This proves all construction and descent assertions of
\Cref{prop:marked-period}, as well as its initial adjoint
bigness.  The next section removes $D_0$ and proves the
tail-rank identity without any hypothesis that $u_*$ span
the whole of $E_0$.

\section{Boundary rank and Higgs tails}\label{sec:boundary-rank}

This section proves the boundary rank loss and the weighted Higgs-tail
comparison that turn the marked-period construction into adjoint
bigness and, later, base nonvanishing.

We retain the marked data of \Cref{sec:marked-periods}.  In particular,
$G$ is the connected rational adjoint monodromy group, $\mathbb A$ is an
irreducible representation of $G_{\C}$ with the compatible polarized
Hodge grading, and $E_0$ is its whole highest Hodge step.  The marks are
$j=(E_0,j_{\rm u})$.  Their orbit $J$ is a product of projective flag
varieties, and each rational simple factor of $G$ has at least one
complex simple factor acting nontrivially on $J$.  The compact marks
occur on factors with an invariant positive Hermitian metric.  On an
appropriate smooth connected unipotent level the line $\mathcal H$ is
$\det E_0+\mathcal H_{\rm u}$; it is understood downstairs as a rational
line.  Its curvature rank is the generic rank of $dj$, whereas the
quotient $S$ remembers the full Lie period and the compact marks.
These two ranks need not be identified in advance.

\subsection{The flat-connection estimate}

To compare numerical ranks, let $T$ be a local fiber of a flag map in
flat coordinates and let $Z$ be its algebraic closure.  Frames in which
the flag has its fixed value form an algebraic incidence over $Z$.
Its fiber dimension is the stabilizer dimension, whereas it contains
the horizontal lift of $T$.  The following estimate bounds the closure
of that lift from below.  In the orbit calculation these two dimensions
will force the differential on $Z$ to fill the monodromy orbit.

The estimate is the semisimple case of the connection Ax--Schanuel
framework of \cite[Theorem~A, Theorem~3.6, and Remark~3.7]{BCFN}.
We give the argument, following the connection-form proof in
\cite[\S3.1]{BCFN}.

\begin{lemma}\label{lem:connection-dimension}
Let $p:\mathcal P\to Z$ be an algebraic principal bundle for a connected
semisimple complex algebraic group $H$, where $Z$ is smooth and
irreducible.  Suppose $\mathcal P$ carries an algebraic flat principal
connection whose monodromy is Zariski dense in $H$.  Let $U$ be an
irreducible analytic germ contained in a horizontal leaf, and suppose
$p(U)$ is Zariski dense in $Z$.  If $R$ is the Zariski closure of $U$ in
$\mathcal P$, then
\begin{equation}\label{eq:prelim-horizontal-dimension}
 \dim R\geq\dim U+\dim H.
\end{equation}
\end{lemma}

\begin{proof}
Put $\mathfrak h=\Lie H$, and let
$\xi\in H^0(\mathcal P,\Omega^1_{\mathcal P}\otimes\mathfrak h)$
be the absolute connection form.  It has kernel the horizontal
distribution, is the identity on fundamental vertical vectors, and
satisfies
\begin{equation}\label{eq:prelim-maurer-cartan}
 d\xi+\tfrac12[\xi,\xi]=0,\qquad
 R_g^*\xi=\operatorname{Ad}(g^{-1})\xi\quad(g\in H).
\end{equation}
Restrict to a nonempty smooth open subset $R^\circ\subset R$ on which
$\xi|_{TR}$ has constant rank $r$.  The germ $U$ meets this subset, because
it is Zariski dense in $R$.  We may replace $U$ by a small connected
smooth germ there without changing its Zariski closure.  Set
$\mathcal K=\ker(\xi|_{TR^\circ})$, and consider the algebraic
Grassmann map
\[
 \beta:R^\circ\longrightarrow\operatorname{Gr}(r,\mathfrak h),
 \qquad x\longmapsto\xi(T_xR).
\]
For a local holomorphic vector field $V$ in $\mathcal K$, Cartan's
formula and \eqref{eq:prelim-maurer-cartan} give
\[
 \mathcal L_V(\xi|_{R^\circ})
 =\iota_Vd\xi+d(\iota_V\xi)=0.
\]
The local flow of $V$ consequently preserves both the restricted
connection form and its image subspace; hence $d\beta(V)=0$.
Since $TU\subset\mathcal K|_U$, the map $\beta$ is constant on $U$.
Its algebraicity and the Zariski density of $U$ make it constant on
$R^\circ$.  Denote its value by $\mathfrak b\subseteq\mathfrak h$.

This fixed subspace is a Lie subalgebra.  In fact, for $b_1,b_2\in
\mathfrak b$, locally choose tangent fields $V_1,V_2$ on $R^\circ$ with
$\xi(V_i)=b_i$.  Such lifts exist because $\xi|_{TR^\circ}$ is a
surjection onto the constant bundle $\mathfrak b$.  Evaluating
\eqref{eq:prelim-maurer-cartan} yields
\[
 \xi([V_1,V_2])=[b_1,b_2],
\]
so $[b_1,b_2]\in\mathfrak b$.

Suppose $\mathfrak b\ne\mathfrak h$.  Let $A$ be the connected analytic
subgroup of $H$ with Lie algebra $\mathfrak b$, and let $B$ be its
Zariski closure.  Then $B$ is a connected algebraic subgroup and is
proper in $H$.  To prove the latter point, if $B=H$, the fact that $A$
preserves $\mathfrak b$ under its adjoint action would imply that $H$
preserves $\mathfrak b$: the stabilizer of this subspace is algebraic.
Thus $\mathfrak b$ would be an ideal of the semisimple Lie algebra
$\mathfrak h$.  Every such ideal is a sum of simple factors and
integrates to a closed connected normal algebraic subgroup of $H$.
That subgroup is $A$, by uniqueness of a connected analytic subgroup
with a given Lie algebra.  A proper $\mathfrak b$ therefore cannot have
Zariski-dense $A$, a contradiction.

Consider the irreducible algebraic subset
\[
 R_B=\overline{R\cdot B}\subseteq\mathcal P.
\]
It dominates $Z$, since $p(U)$ is dense, and it is invariant under the
right action of $B$.  The equivariance in
\eqref{eq:prelim-maurer-cartan} shows that, at a general smooth point,
\begin{equation}\label{eq:prelim-saturated-tangents}
 \xi(T R_B)\subseteq\Lie B.
\end{equation}
Indeed, on the image of the multiplication map $R^\circ\times B\to
\mathcal P$, tangent vectors coming from $R^\circ$ have connection
values in $\operatorname{Ad}(B)\mathfrak b\subseteq\Lie B$, and
tangent vectors coming from $B$ have values in $\Lie B$.  Generic
smoothness of this dominant map onto $R_B$ proves
\eqref{eq:prelim-saturated-tangents}.

The $B$-orbits in a fiber have dimension $\dim B$.  Domination of $Z$
therefore gives
$\dim R_B\geq\dim Z+\dim B$.  Conversely,
\eqref{eq:prelim-saturated-tangents} and
$\dim\ker\xi=\dim Z$ give the reverse inequality.  Equality holds,
and at general smooth points
\[
 TR_B=\xi^{-1}(\Lie B).
\]
In particular the whole horizontal distribution is tangent to $R_B$.
This is an algebraic tangency statement on all of $R_B$: if a local
algebraic horizontal vector field differentiates an equation of $R_B$,
the result vanishes on its dense smooth locus and hence on $R_B$.
Thus the ideal sheaf of $R_B$ is preserved by horizontal vector fields.
Their local holomorphic flows preserve $R_B$ in both directions.

Choose a nonempty smooth open $Z_0\subseteq Z$ over which the general
fiber of $R_B\to Z$ is nonempty of dimension $\dim B$.  Horizontal
transport along a path in $Z_0$ is defined throughout the path: locally
it is the fundamental solution of a holomorphic linear system after a
faithful representation of $H$.  The preceding ideal invariance shows
that it preserves $R_B$.  In particular the monodromy of $Z_0$ preserves
the nonempty closed subset $(R_B)_z$ of the $H$-torsor $\mathcal P_z$.
Deleting a proper algebraic subset of a smooth complex variety does not
decrease the original monodromy image: the map
$\pi_1(Z_0,z)\to\pi_1(Z,z)$ is surjective, since loops can be perturbed
off a subset of real codimension at least two.  This monodromy is
therefore Zariski dense in $H$.  A closed subset of an $H$-torsor
invariant under a Zariski-dense translation subgroup is invariant under
$H$, and is either empty or the whole torsor.  But $(R_B)_z$ is nonempty
and has dimension $\dim B<\dim H$.  This contradiction proves
$\mathfrak b=\mathfrak h$.

Finally the generic kernel of $\xi|_{TR}$ contains $TU$ and hence has
dimension at least $\dim U$.  Its generic image has dimension $\dim H$,
so rank--nullity gives \eqref{eq:prelim-horizontal-dimension}.
\end{proof}


\subsection{The orbit calculation}

In this subsection, a differential of a period, flag, or projective
line is computed in local flat frames.  A local fiber is a connected
smooth analytic fiber germ on the locus of constant maximal rank.
The phrase Hodge-generic refers to the integral rational Lie variation,
including its tensor constructions, as in \Cref{lem:marked-monodromy}.

\begin{lemma}[Orbit rank]\label{lem:orbit-rank}
Let $B^\circ$ be a smooth connected algebraic base carrying the marked
data above, with Zariski-dense connected monodromy $G$.  The same
conclusions hold for their pullback by a dominant algebraic map with
connected generic fiber.  Consider either the joint flag map $j$, or
an algebraic projective line section $\ell$ of a flat representation
of $G$.

\emph{Local orbit filling.}
Let $T$ be a local fiber at a Hodge-generic point and let $Z$ be its
irreducible Zariski closure.  After passing to a finite level on a
smooth open of $Z$, its connected monodromy is a rational normal
factor $H$ of $G$.  Write $G=H\times Q$.  The full $Q$ period is
constant on the lifted $Z$.  In the flag case the $Q$ compact marks
are constant there as well.  At suitable smooth generic points of
$T$, writing $a$ for $j$ or $\ell$, one has
\begin{equation}\label{eq:boundary-orbit-fill}
 \dim Z-\dim T=\dim(H_{\C}\cdot a),\qquad
 da(T_zZ)=T_{a(z)}(H_{\C}\cdot a(z)).
\end{equation}
In the flag case write $J=J_H\times J_Q$ for the products of its $H$
and $Q$ flag factors, and $j=(j_H,j_Q)$ for the corresponding
projections in local flat frames.  The orbit in
\eqref{eq:boundary-orbit-fill} is then $J_H$.

\emph{Flag ranks and lines.}
For the flag map on a level of $S$, let $q:B\to S_Q$ denote a resolved
connected-fiber quotient for the full $Q$ period together with its
compact marks, and set $s_Q=\dim S_Q$ and $d_H=\dim J_H$.  If
$r=\rk(dj)$, then
\begin{equation}\label{eq:boundary-rank-split}
 \rk(dj_Q)=s_Q,\qquad r=d_H+s_Q.
\end{equation}
On suitable smooth projective models, with $\rho:B\to S$ the level
map, the nef lines split as
\begin{equation}\label{eq:boundary-line-split}
 \rho^*\mathcal H=\mathcal H_H+\mathcal H_Q,
 \qquad \nu(\mathcal H_H)=d_H,
 \qquad \mathcal H_Q=q^*\underline{\mathcal H}_Q,
 \qquad \nu(\underline{\mathcal H}_Q)=s_Q.
\end{equation}
Here all line equalities are rational line-bundle equalities.
\end{lemma}

\begin{proof}
We may choose the initial point outside the countably many exceptional
tensor loci and in the maximum-rank loci of the maps associated with
every rational normal-factor subset of $G$.  There are only finitely
many such subsets.  The restriction to $Z$ has the same generic
Mumford--Tate group as the ambient variation: a tensor condition
holding on $Z$ holds at the chosen Hodge-generic point.  The connected
monodromy normality theorem, in the pure polarized integral setting
specified in \Cref{lem:marked-monodromy}, makes $H$ normal in that
Mumford--Tate group.  Since $H\subseteq G$ and $G$ is semisimple
adjoint, $H$ is a product of rational simple factors and has the
complementary factor $Q$.

The restricted $Q$ monodromy is finite.  Kill it by a finite cover of
a smooth open of $Z$.  The fixed-part theorem then says that the full
$Q$ variation is constant.  Its flat frames are algebraic frames:
the algebraic connection has regular singularities and its trivial
local system is the trivial regular-singular connection.  The
$Q$ compact marks are algebraic in these frames.  They are constant
on the lifted $T$, since $j$ was constant there, and hence on the
lifted $Z$.  Indeed, a lift of a nonempty open subgerm of $T$ is
Zariski dense in an irreducible component of this finite cover.  Its
closure has dimension at least $\dim Z$, because its finite image
contains a Zariski-dense subset of $Z$; an irreducible component of
the cover also has dimension $\dim Z$.  This proves the required
density assertion after finite levels.  In the line case we only
need the constancy of the $Q$ period.

Use the algebraic flat $H_{\C}$-frame torsor over this smooth open of
$Z$.  It is the reduction specified by the parallel tensor
invariants; its connection has Zariski-dense monodromy $H$.  Set
$a_0=a(T)$ in a flat trivialization along $T$.  Consider the
algebraic incidence consisting of frames in which $a$ has value
$a_0$.  Over a point where it is nonempty its fiber is a coset of
the stabilizer of $a_0$ in $H_{\C}$.  Its image contains $T$ and is
constructible, so contains a dense open of $Z$.  If
$d=\dim(H_{\C}\cdot a_0)$, the incidence over that open therefore
has dimension
\[
                  \dim Z+\dim H-d.
\]
It contains the horizontal lift of $T$, whose Zariski closure therefore
has at most this dimension.  Applying
\Cref{lem:connection-dimension} to that horizontal germ gives
\[
 \dim T+\dim H\leq\dim Z+\dim H-d,
 \quad\hbox{or equivalently}\quad
 \dim Z-\dim T\geq d.
\]
This use of the incidence works equally for a locally closed
projective-line orbit; projectivity of the orbit is not needed.

Move within $T$ to a smooth point of $Z$ lying over the dense open
just obtained.  Nonempty open subgerms are still Zariski dense in
$Z$.  In flat frames the values of $a$ on a neighborhood in $Z$
lie in $H_{\C}\cdot a_0$.  Moreover
\[
 \ker(da|_{T_zZ})=T_zZ\cap\ker da=T_zT,
\]
because the ambient map has constant rank near $T$.  Thus
$\dim Z-\dim T\leq d$.  Equality proves
\eqref{eq:boundary-orbit-fill}.  For the joint marks, the
factor-by-factor description of $J$ identifies this orbit with
$J_H$.

We now work with $a=j$ on a level of $S$.  The $Q$ joint data is
constant on $Z$, so $T_zZ\subseteq\ker dq$, while
$\ker dj=T_zT\subseteq\ker dq$ at the chosen points.  The orbit
calculation gives
\[
                   dj_H(\ker dq)=T_{j_H(z)}J_H.
\]
Since $j_Q$ descends through $q$, one has
$\ker dq\subseteq\ker dj_Q$.  Conversely, if $dj_Q(v)=0$, choose
$w\in\ker dq$ with $dj_H(w)=dj_H(v)$.  Then $dj(v-w)=0$, hence
$v-w\in\ker dq$, and therefore $v\in\ker dq$.  It follows that
$\ker dj_Q=\ker dq$.  Computing the rank first along $\ker dq$
and then on its quotient proves \eqref{eq:boundary-rank-split}.
Our initial choice in the maximum-rank loci makes these generic
rank identities, even though the argument was carried out at
points of $T$.

The representation decomposes at connected level as
$\mathbb A=A_H\otimes A_Q$, and its highest step as
$E_0=E_H\otimes E_Q$.  The factors are the compatible Hodge
representations constructed from the separate adjoint Lie tensor
systems; a constant scalar shift of a grading has no effect on
the following determinant calculation.  If $e_H=\rk E_H$ and
$e_Q=\rk E_Q$, then
\[
 \det E_0=(\det E_H)^{\otimes e_Q}
                    \otimes(\det E_Q)^{\otimes e_H}.
\]
Add the compact-mark line on the appropriate factors to define
$\mathcal H_H$ and $\mathcal H_Q$.  Their semipositive curvature
forms have kernels $\ker dj_H$ and $\ker dj_Q$.  The first has
rank $d_H$, by orbit filling and the upper bound $\dim J_H$;
the second has rank $s_Q$.  \Cref{lem:log-metric} gives their
numerical dimensions.

The $Q$ representation and flags descend through $q$ by
\Cref{lem:period-descent}, applied to precisely the rational
normal factors in $Q$.  Its integral lattice inherits the
congruence-level condition, so its quasi-unipotent boundary
monodromies are already unipotent.  Choose smooth projective
models resolving $q$ and all the flag sections.  Unipotent
canonical extension commutes with this pullback: in a monomial
boundary chart the new logarithmic residues are integral sums
of the commuting old nilpotent residues, and the untwisted
filtration pulls back.  Thus the equality of the $Q$ lines on
the open set extends as the equality in
\eqref{eq:boundary-line-split}.  The metric rank on the quotient
is $s_Q$, so the same metric lemma gives
$\nu(\underline{\mathcal H}_Q)=s_Q$.
\end{proof}

\subsection{Strict loss of numerical dimension at the boundary}

The reduced divisor $D_0$ on $S$ consists of the boundary components
whose monodromy on the retained rational Lie variation
$\mathfrak g$ is not the identity.  In particular finite
nonidentity monodromy is included.

\begin{proposition}[Boundary drop]\label{prop:boundary-drop}
For every component $D_i$ of $D_0$,
\begin{equation}\label{eq:boundary-strict-drop}
                 \nu(\mathcal H|_{D_i})<\nu(\mathcal H).
\end{equation}
If $\nu(\mathcal H)=0$, then $D_0$ is empty.
\end{proposition}

\begin{proof}
Put $r=\nu(\mathcal H)=\rk dj$ and use
\Cref{lem:orbit-rank}.  Write $B$ for its smooth unipotent level,
and put $n=\dim B=\dim S$.  On $B$ choose a
smooth boundary component $D'$ mapping generically finitely onto
$D_i$.  Further resolutions and restrictions to the strict
transform are allowed.  For any nef line $L$ on a smooth
projective variety, restriction to a divisor cannot increase
its numerical dimension: choose an ample $A$ with $mA-D'$
effective and compare the nonnegative intersections
$L^kD'A^{n-k-1}$ and $mL^kA^{n-k}$.
Consequently
\[
                 \nu(\mathcal H_H|_{D'})\leq d_H.
\]
If $D'$ does not dominate $S_Q$, then
\[
 \nu(\mathcal H_Q|_{D'})\leq\dim q(D')<s_Q.
\]
The binomial expansion of intersections of nef classes gives
$\nu(P+Q)\leq\nu(P)+\nu(Q)$ for nef $P,Q$.  Applied on $D'$,
this proves the strict inequality in this case.  We may therefore
assume that $D'$ dominates $S_Q$.  At its generic point the
$Q$ system is pulled back from the interior of $S_Q$.

\smallskip\noindent
\emph{Infinite local monodromy.}
The logarithm $N$ of a sufficiently divisible power of the
original monodromy is then nonzero and belongs to the rational
Lie algebra of $H$.  Its $Q$ component is zero because the
$Q$ variation extends from the interior at the generic point
of $D'$.  It acts nontrivially on the $H$ orbit of $E_H$.
To check this last assertion, a nonzero rational element in
a rational simple factor has nonzero projection to every
Galois-conjugate complex simple factor.  A unipotent isometry
of a positive Hermitian space is the identity.  Thus a rational
factor detected only by compact marks cannot support a
nonzero component of $N$.  Every other retained rational factor
has a visible first-kind flag factor, and the action of a
complex simple adjoint group on a nontrivial flag variety is
faithful.  A nonzero component of $N$ therefore moves such a
flag.

Let $p_H$ be the highest Hodge index of $A_H$ and set $p_0=e_Hp_H$.
There is no restriction on the rank $e_H$ of $E_H$.  Form
\[
           \mathbb B=\bigwedge^{e_H}A_H,
           \qquad L_H=\det E_H=F^{p_0}\mathbb B,
           \qquad F^{p_0+1}\mathbb B=0.
\]
The highest filtration step of $\mathbb B$ is exactly the
one-dimensional line $L_H$.  Denote the weight of $\mathbb B$
by $w$.  On the open stratum of $D'$, untwist by
$\exp(-\log(z)N/(2\pi\sqrt{-1}))$ in a local normal coordinate
$z$ and let $[v]$ be the limiting highest line.  Schmid's
nilpotent-orbit theorem and limiting mixed Hodge theorem
\cite[Theorems~4.12 and~6.16]{Schmid} apply to the polarized
unipotent variation.  The same assertions hold in $\mathbb B$:
it is a Hodge-compatible summand of rational adjoint tensor
systems, whose flat projectors commute with $N$ and preserve
the two limiting Hodge filtrations and the weight filtration.

Write $W_\bullet=W(N)_\bullet$, centered at $w$.  The line
$F^{p_0}$ occupies a single summand $I^{p_0,q}$ of the Deligne
splitting, since there are no summands of first index greater
than $p_0$.  Set $p_0+q=w+k$.  The primitive decomposition of
the monodromy weight grades shows that $k\geq0$ and that the
class of $v$ in $\Gr^W_{w+k}$ is primitive.  Here is the
specific reason.  A nonprimitive term of first Hodge index
$p_0$ has the form $N^j u$ with $j>0$ and with $u$ of first
index $p_0+j$, which is impossible.  Likewise a weight grade
below $w$ is entirely obtained from positive powers of $N$
on higher primitive grades, and so cannot contain the highest
line.  Monodromy Lefschetz gives a nonzero $N^k[v]$ in
$\Gr^W_{w-k}$ and zero $N^{k+1}[v]$.  In fact
\begin{equation}\label{eq:boundary-actual-primitive}
                         N^kv\ne0,\qquad N^{k+1}v=0
\end{equation}
for the actual vector: $N$ has bidegree $(-1,-1)$ in the
Deligne splitting, so each $N^jv$ lies in a single summand;
its vanishing on the corresponding weight grade is its
vanishing as a vector.  The primitive decomposition and
Lefschetz isomorphisms used here are those of the monodromy
weight filtration, as in \cite[Lemma~6.4 and Theorem~6.16]{Schmid}.

Let $O_H$ be the closed orbit of the highest line of
$\mathbb B$, namely the exterior-power embedding of the
$H$ orbit of $E_H$.  The untwisting acts through $H_{\C}$,
so $[v]\in O_H$.  By \eqref{eq:boundary-actual-primitive},
\[
             \lim_{\tau\to\infty}\exp(\tau N)[v]=[N^kv]
                         \in O_H\cap\mathbf P(\ker N).
\]
The intersection on the right is a proper closed subset of
the irreducible variety $O_H$, because $N$ acts nontrivially
on $O_H$.  Project $N^kv\in W_{w-k}$ to $\Gr^W_{w-k}$ and
then apply the fixed inverse of
\[
                N^k:\Gr^W_{w+k}\xrightarrow{\ \sim\ }
                                      \Gr^W_{w-k}.
\]
The resulting projective point is the highest graded line
$[v]_{\Gr}$.  These maps are fixed linear maps in local flat
limiting coordinates.  Projection is taken on the open set
where its value is nonzero; this set contains all the points
under consideration by Lefschetz.  The local variation rank
of the highest graded line is therefore at most
$\dim O_H-1$.  Together with all the $H$ compact marks its
rank is at most
\begin{equation}\label{eq:boundary-graded-rank}
                    \dim O_H-1+\dim J_{{\rm u},H}=d_H-1.
\end{equation}

\smallskip\noindent
\emph{From limiting rank to numerical dimension.}
\phantomsection\label{proof:limiting-weight-comparison}
We next justify that the graded rank just obtained bounds the numerical
dimension of the restricted line, including at crossings.
One cannot simply restrict the original singular Hodge
metric to $D'$.  Instead, the pure weight grades of the
limiting mixed Hodge structures give polarized pure
variations on the open stratum of $D'$
\cite[Proposition~2.10]{CK}.  Their filtrations
are holomorphic and transverse, by the extended filtration
and the tangential logarithmic connection.  Their primitive
parts have the monodromy polarizations, and their remaining
parts are polarized through the Lefschetz decomposition.

There is a canonical extension description on the whole
component.  In a Deligne frame at a crossing
$z_1\cdots z_a=0$, with $D'=(z_1=0)$, the transverse residue
is the constant nilpotent $N_1$ and the remaining residues
$N_2,\ldots,N_a$ commute with it.  The filtration $W(N_1)$
therefore consists of subbundles of the restricted extended
flat bundle.  Every remaining residue preserves this
filtration and induces a nilpotent operator on each weight
grade.  Changing the normal coordinate or the tangential
lift changes the tangential connection by a multiple of
$N_1$; this acts trivially on the weight grades.  Their flat connections are thus well defined. The remaining logarithmic
exponentials are exactly the Deligne extensions of these flat graded
variations. Separately, each polarized pure graded variation has its
Schmid-extended Hodge subbundles inside that canonical flat bundle.
We do not identify a possibly jumping intersection with a weight step
with an extended Hodge subbundle. Instead we use the independent Hodge
extension, and below map into it by continuation from the open stratum.
This distinction proves the needed comparison also at crossings.

On the stratum, $L_H|_{D'}$ lies in $W_{w+k}$ and maps
nontrivially to its highest graded line.  Inclusion in the
subbundle $W_{w+k}$ extends everywhere because it holds
generically.  Projection to $\Gr^W_{w+k}$ is regular.  The highest line of the graded variation extends as its independent
Schmid Hodge subbundle. The projection lands in it: its image in the
locally free quotient of the canonical flat graded bundle vanishes on
the dense open stratum and hence vanishes everywhere. The resulting
map can have zeros at further boundary strata; these produce an
effective difference in the direction needed below.  If necessary resolve
the additional algebraic filtration-rank loci on $D'$.
The loops newly omitted from the interior have trivial
monodromy, and all boundary monodromies remain unipotent;
the preceding description is preserved on this resolution.
We obtain a generically nonzero morphism of line bundles
\begin{equation}\label{eq:boundary-line-comparison}
                 L_H|_{D'}\longrightarrow L_{H,\Gr},
                 \qquad L_{H,\Gr}-L_H|_{D'}\ \text{effective}.
\end{equation}
Here and below restriction notation includes the pullback
to that resolution when one was needed.

Both sides of \eqref{eq:boundary-line-comparison} are nef.
For the right side this follows from the pure highest-line
metric and \Cref{lem:log-metric}; for the left side it follows
by restricting the original nef line.  Tensor the comparison
by $e_Q$ and add the same compact-mark line on $D'$.  If
$P,Q$ are the resulting nef lines, then $Q-P$ is effective,
and for an ample $A$ on this $(n-1)$-dimensional variety,
\[
 (Q^k-P^k)A^{n-1-k}
  =(Q-P)\sum_{a=0}^{k-1}Q^{k-1-a}P^a A^{n-1-k}\geq0.
\]
Thus $\nu(P)\leq\nu(Q)$.  The metric lemma and
\eqref{eq:boundary-graded-rank} now give
\[
                    \nu(\mathcal H_H|_{D'})<d_H.
\]
Adding $\mathcal H_Q|_{D'}$, whose numerical dimension is
at most $s_Q$, proves the desired strict drop upstairs.

\smallskip\noindent
\emph{Finite nonidentity local monodromy.}\label{proof:finite-boundary}
Let $\gamma\ne1$ denote the original monodromy.  On the
unipotent level its logarithm is zero, so the pure period
extends ordinarily across a general point of $D'$.  Its
limiting flags and compact marks are fixed by $\gamma$.
Indeed, the local cover is $t=z^e$; the continuation around
$t=0$ acts on its sheets by $z\mapsto e^{2\pi\sqrt{-1}/e}z$.
Equivariance of the extended period and marks then gives
the fixed-point assertion at $z=0$ in common flat frames.
The comparison in \eqref{eq:boundary-line-comparison}, now
with a pure ordinary limit, shows that the numerical
dimension on $D'$ is bounded by the differential rank of
these limiting marks.  That rank is at most $r$: in the
ordinary extending flat frame, all $(r+1)$-minors of the
mark differential vanish on the interior and hence on
its extension, including after restriction to $D'$.

Suppose that rank were $r=d_H+s_Q$.  Since $D'$ dominates
$S_Q$, the limiting $Q$ data comes from the interior of
$S_Q$.  In particular $dj_Q$ on $D'$ has rank $s_Q$ and
kernel $\ker(dq|_{D'})$.  Equality of the total rank forces
the $H$ flags along a local general fiber of $q|_{D'}$ to
fill an open subset of $J_H$.  The $Q$ flags are fixed
on this fiber.

The rational automorphism $\gamma$ normalizes the connected
group, and acts on the product $J$ factorwise up to
permutation of simple factors.  An open subset with all
$H$ flag coordinates varying independently cannot be fixed
if any visible $H$ coordinate is moved to a different
coordinate.  If the destination were a $Q$ coordinate it
would be constant; if it were another $H$ coordinate the
fixed-point equations would impose a relation between two
independent coordinates.  Thus every visible complex
$H$ factor is preserved and $\gamma$ acts identically on
its whole flag variety.  Equivariance and faithfulness of
the simple adjoint action imply identity on its Lie
algebra.  Rationality then implies identity on every
Galois-conjugate complex factor, hence on all of $H$.
Consequently $\gamma$ preserves $Q$.  Its action on $Q$
fixes a generic lifted full $Q$ period, because $D'$ dominates
$S_Q$.  \Cref{lem:generic-lie-stabilizer} gives identity on
$Q$ also, a contradiction.  The limiting rank is strictly
less than $r$ in this case as well.

Pullback by the generically finite map $D'\to D_i$ preserves
the numerical dimension of a nef line, so the strict drop
descends to $D_i$.  Finally, if $r=0$, a local fiber of $j$
is open in the base, so its Zariski closure is the entire
base and $H=G$.  Orbit filling gives $\dim J=0$.
Visibility on every retained rational factor forces
$G=1$; the full projected adjoint action then has no
nonidentity boundary monodromy.  Thus $D_0$ is empty.
\end{proof}

\subsection{Removing the marked boundary}

\begin{lemma}[Section estimate]\label{lem:boundary-remove}
Let $M$ be a smooth projective $n$-fold, let $H$ be a nef
rational line, and let $T=\sum_i T_i$ be a reduced divisor
with smooth components.  Suppose $r=\nu(H)>0$ and
$\nu(H|_{T_i})<r$ for every $i$.  If $A$ is ample and
$e>0$ is an integer, then
\[
                         A+tH-eT
\]
is big for all sufficiently large divisible integers $t$.
\end{lemma}

\begin{proof}
Choose fixed ample integral lines $A_i$ on $T_i$ sufficiently
positive that all successive restriction terms obtained
while imposing order $le$ along each $T_j$ inject into
\[
 H^0\bigl(T_i,l(A|_{T_i}+tH|_{T_i}+A_i)\bigr).
\]
This choice is uniform in $l$ and $t$.  Explicitly, choose
ample $C_{ij}$ so that $C_{ij}$ and $C_{ij}+T_j|_{T_i}$
have nonzero sections, and put $A_i=e\sum_j C_{ij}$,
enlarging it if needed.  For $0\leq n_j\leq el$ the line
\[
 lA_i+\sum_j n_jT_j|_{T_i}
 =\sum_j\bigl((el-n_j)C_{ij}
                  +n_j(C_{ij}+T_j|_{T_i})\bigr)
\]
has a nonzero section.  Multiplication by it supplies
exactly the asserted injections, including the normal
twists when $j=i$.

The restriction exact sequences, applied one copy of a
component at a time, give
\begin{align*}
 h^0\bigl(M,l(A+tH-eT)\bigr)
 &\geq h^0\bigl(M,l(A+tH)\bigr)\\
 &\quad-le\sum_i
 h^0\bigl(T_i,l(A|_{T_i}+tH|_{T_i}+A_i)\bigr).
\end{align*}
For each fixed divisible $t$, all the positive lines in
this formula are ample.  Their Hilbert polynomials give,
in the coefficient of $l^n$, a positive supply
\[
                         \frac{(A+tH)^n}{n!}
\]
and a loss at most
\[
 \frac{e}{(n-1)!}\sum_i
                (A|_{T_i}+A_i+tH|_{T_i})^{n-1}.
\]
The supply is a polynomial in $t$ of degree exactly $r$
with positive leading coefficient.  Each term in the
loss has degree at most $r-1$.  For $n=1$ the restriction
terms are simply section dimensions on points and the
same statement holds.  Choose $t$ sufficiently large
that the difference of these leading coefficients is
positive, and then let $l\to\infty$.  This proves bigness.
No asymptotic estimate uniform in both $l$ and $t$ is used.
\end{proof}

\begin{proposition}[Adjoint bigness]\label{prop:adjoint-bigness}
If $\dim S>0$, then $K_S+b\mathcal H$ is big for all
sufficiently large rational $b$.
\end{proposition}

\begin{proof}
By \Cref{lem:initial-adjoint},
$K_S+D_0+c\mathcal H_{\rm u}$ is big for sufficiently
large $c$.  Since $\mathcal H-\mathcal H_{\rm u}$ is nef,
$K_S+D_0+c\mathcal H$ is big.  Choose an integer $e>0$
clearing denominators and an ample integral $A$ with
\[
                   e(K_S+c\mathcal H+D_0)-A
\]
effective.  \Cref{prop:boundary-drop,lem:boundary-remove}
show that $A+t\mathcal H-eD_0$ is big for large divisible
$t$.  Adding the displayed effective divisor proves
bigness of $eK_S+(ec+t)\mathcal H$.  When $\nu(\mathcal H)=0$,
\Cref{prop:boundary-drop} says $D_0=0$, and the initial
bigness already gives the conclusion.  Finally the sum
of a big line and a nef line is big, so every larger
rational $b$ works as well.
\end{proof}

We have proved adjoint bigness on the parameter space.  To use it for a
chosen vector, we need a nef line generated by its Higgs images whose
numerical dimension does not increase on adding the marked line.
The next proposition supplies this second input to the removal of the
period twist in \Cref{sec:nonvanishing}.

\subsection{The rank of the Higgs-tail line}

\begin{proposition}[Higgs tails]\label{prop:tail-rank}
Let $h:W\to S$ be any dominant morphism of smooth
connected projective varieties with connected generic
fiber.  On a dense algebraic open of $W$, suppose that
$\ell\subset h^*E_0$ is an algebraic line subbundle.
Require that the pointwise instability-parabolic
construction of \Cref{lem:marked-flags}, applied to
$\ell$, gives exactly the pulled-back compact marks
$h^*j_{\rm u}$.  This compatibility is part of the
hypothesis.

Take a smooth connected unipotent level
$\pi:\widehat W\to W$ dominating the level used to
define $\mathcal H$, and resolve its boundary and the
generated subspace maps.  Over the variation locus,
let $G_i$ be the span of all contractions of
$\theta^i(\ell)$ in
$\pi^*h^*\Gr_F^{p_*-i}\mathbb A$, with $G_0=\ell$.
Use the resulting subbundles of the extended Hodge
grades on $\widehat W$.  Put
\begin{equation}\label{eq:tail-line-definition}
 U_j=\bigoplus_{i\geq j}G_i,
 \qquad
 P=-\sum_{j\geq0}\det U_j
   =-\sum_{i\geq0}(i+1)\det G_i.
\end{equation}
Only finitely many summands are nonzero.  Then $P$ is
nef and
\begin{equation}\label{eq:tail-rank-equality}
              \nu\bigl(P+\pi^*h^*\mathcal H\bigr)=\nu(P).
\end{equation}
The assertion applies to every such $h$ and every such
compatible line; $\ell$ need not be the original vector
used to construct $S$.
\end{proposition}

The curvature argument will show that a null direction for $P$ fixes
the chosen line and its compact marks.  Passing from that line to the
whole highest Hodge space requires the following representation lemma.
We state it here, use it in the proof of the proposition,
and give its root-theoretic proof afterward.

\begin{lemma}[Degree-one roots]\label{lem:degree-one-root}
Let $\mathfrak g=\bigoplus_k\mathfrak g_k$ be a complex
semisimple Lie algebra with an integral grading and an
adjoint operation $*$ carrying $\mathfrak g_k$ to
$\mathfrak g_{-k}$, arising from a compatible compact
form.  Let $\mathfrak h$ be a graded ideal stable under
$*$, and let $V$ be an irreducible finite-dimensional
representation with compatible grading and highest
graded piece $E$.  The metric on $V$ is chosen so that
representation adjoints are the actions of Lie adjoints.
For $0\ne s\in E$ and $x\in\mathfrak h_1$, suppose
\begin{equation}\label{eq:root-annihilation-hypothesis}
                         [x,\mathfrak h_{-1}]s=0.
\end{equation}
Then $x^*E=0$.
\end{lemma}

\begin{proof}[Proof of Proposition~\ref{prop:tail-rank}]
The connected-fiber hypothesis ensures that, after
shrinking to a smooth locally trivial locus, the
monodromy image of the pullback is the same as that
on $S$.  Passing to its connected level therefore
leaves Zariski-dense monodromy $G$.  The variation
and the orbit calculation of \Cref{lem:orbit-rank}
apply on $\widehat W$.

Each $U_j$ is Higgs invariant because
$\theta G_i\subseteq G_{i+1}\otimes\Omega^1$; this
follows from Higgs integrability and the definition
by all contractions.  On the open locus the Hodge
decomposition is orthogonal.  Write $\mathcal D$ for
its Chern connection.  For a tangent vector $v$ put
$A=\theta(v)$, let $\Pi_j$ be orthogonal projection
onto $U_j$, and set
\[
             B_j(v)=(1-\Pi_j)\mathcal D'_v|_{U_j}.
\]
The Hodge curvature equation and the subbundle
curvature formula give, in a common positive
normalization,
\begin{equation}\label{eq:tail-curvature}
 \Theta_{-\det U_j}(v,\bar v)
   =\bigl\|(1-\Pi_j)A^*|_{U_j}\bigr\|_{\rm HS}^2
                      +\|B_j(v)\|_{\rm HS}^2.
\end{equation}
To see the sign directly, use $A U_j\subseteq U_j$
and write $A$ in upper block-triangular form for
$U_j\oplus U_j^\perp$.  The traces of the two
internal products cancel.  The remaining trace is
minus the square of the off-diagonal block of
$A^*|_{U_j}$ for $\det U_j$, and subbundle curvature
subtracts the square of $B_j(v)$.  Negating the
determinant gives \eqref{eq:tail-curvature}.

The metrics on the extended subbundles have the
two-sided logarithmic norm estimates needed in
\Cref{lem:log-metric}.  Indeed they are restrictions
of the extended Hodge metrics; after resolution
their frames have full rank in the extended Hodge
grades, so the dual estimates also restrict.
Thus every $-\det U_j$, and hence $P$, is nef.

Suppose now that $\Theta_P(v,\bar v)=0$.  Every
nonnegative summand in \eqref{eq:tail-curvature}
vanishes.  Applying its first term for $U_i$ to
$G_i$, the vector $A^*G_i$ belongs to the preceding
Hodge grade, which is orthogonal to $U_i$.  Hence
\begin{equation}\label{eq:tail-zero-actions}
                A^*G_i=0\quad\hbox{for every }i.
\end{equation}
Since $AG_i\subseteq G_{i+1}$, for $g\in G_i$ one
then has $\|Ag\|^2=\langle g,A^*Ag\rangle=0$.
Thus $AG_i=0$ as well.  Moreover $B_0(v)=0$ and
$\mathcal D'$ preserves the Hodge grading, so
$\mathcal D'_v$ preserves $G_0=\ell$.  The flat
connection satisfies
$\nabla^{1,0}=\mathcal D'+\theta$ in the smooth
Hodge splitting.  Its projective derivative on
$\ell$ is consequently zero.  In other words,
\begin{equation}\label{eq:tail-null-line}
                         v\in\ker d\ell.
\end{equation}
The compact marks, which by hypothesis are the
same pointwise functions of $\ell$ as before,
are constant on a local $\ell$-fiber.  At its
constant-rank points \eqref{eq:tail-null-line}
also gives $dj_{\rm u}(v)=0$.

It remains to prove $dE_0(v)=0$.  Apply the line
version of \Cref{lem:orbit-rank} to a local
$\ell$-fiber $T$, with Zariski closure $Z$ and
decomposition $G=H\times Q$.  Work at one of the
smooth generic points $z$ supplied by that lemma.
The preceding calculation places every
$v\in\ker\Theta_{P,z}$ in $T_zT$; fix this same
$T$, $Z$, and $H$ for all such vectors.  The full
$Q$ period is constant on $Z$.  Hence the
infinitesimal period element of every such $v$ satisfies
\[
                    t(v)\in\mathfrak h^{-1,1}.
\]
Write $\ell_z=\C s$.  Evaluating the tangent map
of $\ell$ at $s$ and projecting to the next Hodge
grade sends $d\ell(T_zZ)$ into $G_1$: its flat
derivative is $\mathcal D'_w s+\theta(w)s$, and
only the second term contributes to that grade.
Orbit filling in \eqref{eq:boundary-orbit-fill}
therefore gives
\begin{equation}\label{eq:tail-orbit-action}
                         \mathfrak h^{-1,1}s\subseteq G_1.
\end{equation}
Use the metric supplied by the adjoint tensor
construction, so that representation adjoints
agree with Lie adjoints.  Set $x=t(v)^*$.
By \eqref{eq:tail-zero-actions},
$xG_1=0$.  Also $xs=0$ because $s$ is in the
highest Hodge grade and $x$ raises that grade.
It follows from \eqref{eq:tail-orbit-action} that
\[
                         [x,\mathfrak h^{-1,1}]s=0.
\]
\Cref{lem:degree-one-root}, with
$\mathfrak g_k=\mathfrak g^{k,-k}$, gives
$t(v)E_0=0$.  The flat-chart differential of the
whole highest step is precisely Higgs evaluation
on $E_0$, so $dE_0(v)=0$.

We have proved, at the points where the orbit
calculation applies, the kernel inclusion
\begin{equation}\label{eq:tail-kernel-inclusion}
                  \ker\Theta_P\subseteq
                          \ker\Theta_{\pi^*h^*\mathcal H}.
\end{equation}
For completeness, these points can be chosen to
compute the global curvature ranks.  The Hodge
metrics and subspace data are real analytic on
the variation locus.  Choose the initial point
in the open maximum-rank loci of $\Theta_P$ and
$\Theta_P+\Theta_{\pi^*h^*\mathcal H}$, as well
as the finitely many relevant period, flag, and
projective-line rank loci and the Hodge-generic locus.  Moving
slightly within its $\ell$-fiber reaches the
smooth Zariski-open locus of $Z$ required by
\Cref{lem:orbit-rank}, while staying in the two
curvature-rank opens.  Indeed an open subgerm of
that fiber is Zariski dense in $Z$.  The argument
then establishes \eqref{eq:tail-kernel-inclusion}
at a point where both curvature ranks are
maximal.  Semipositivity implies equality of
those ranks.  Apply \Cref{lem:log-metric} to
$P$ and $P+\pi^*h^*\mathcal H$ to obtain
\eqref{eq:tail-rank-equality}.
\end{proof}

\begin{proof}[Proof of Lemma~\ref{lem:degree-one-root}]
Choose a Cartan in the compact form adapted to the grading,
and choose positive roots so that each simple root has
nonnegative degree.  Metric adjoints exchange opposite
root spaces.  Let $\Delta_0$ be the degree-zero simple
roots, and let $v_\lambda$ be a highest-weight vector
for $V$.  Write $\mathfrak g_0^-$ for the sum of the
negative root spaces of degree zero.  The highest graded piece is
\begin{equation}\label{eq:root-top-module}
                  E=U(\mathfrak g_0^-)v_\lambda.
\end{equation}
In particular it is irreducible for the full degree-zero
reductive algebra.  Indeed, reaching the extremal grading
from the highest weight allows only degree-zero
lowerings; the highest-weight theory for that Levi
algebra then gives \eqref{eq:root-top-module}.

Every degree-one positive root contains exactly one
degree-one simple root $\alpha$, with coefficient one,
and otherwise only simple roots of degree zero.  Let
$\mathfrak b^+_\alpha$ be the sum of its root spaces and
$\mathfrak b^-_\alpha$ the sum of the opposite root
spaces.  These are $\mathfrak g_0$-modules, and $*$
exchanges them.  Roots of degree one in $\mathfrak h$
are exactly the blocks belonging to its simple factors.
For $\alpha\ne\beta$,
\begin{equation}\label{eq:root-cross-bracket}
                 [\mathfrak b^+_\alpha,
                         \mathfrak b^-_\beta]=0.
\end{equation}
Indeed, the difference of the relevant roots has a
positive coefficient at $\alpha$ and a negative one at
$\beta$, and hence is neither a root nor zero.

Write $x=\sum_\alpha x_\alpha$ using these blocks in
$\mathfrak h_1$.  For every nonzero $x_\alpha$,
\eqref{eq:root-annihilation-hypothesis} and
\eqref{eq:root-cross-bracket} imply
$[x_\alpha,\mathfrak b^-_\alpha]s=0$.  Hence the incidence
\[
 I_\alpha=
 \left\{([z],[u])\in
  \mathbf P(\mathfrak b^+_\alpha)\times\mathbf P(E):
             [z,y]u=0\text{ for every }y\in\mathfrak b^-_\alpha\right\}
\]
is nonempty, closed, and projective.  It is stable under
the full connected degree-zero group.  Apply the Borel
fixed-point theorem to the Borel of this group determined
by the chosen positive roots and containing the full Cartan.
The theorem used here says that a
connected solvable algebraic group acting on a nonempty
complete variety has a fixed point
\cite[Corollary~17.3]{Milne}.  A fixed pair in $I_\alpha$
has the form $([e_\rho],[v_\lambda])$: the first
coordinate is a root line since the full Cartan acts,
and the second is the unique highest line of the
irreducible degree-zero module $E$.  Irreducibility of
$\mathfrak b^+_\alpha$ is not required.

Let $C_\alpha$ be the connected component containing
$\alpha$ in the Dynkin subdiagram on
$\Delta_0\cup\{\alpha\}$.  Connectedness of root support
gives $\operatorname{Supp}(\rho)\subseteq C_\alpha$.
If the inclusion were strict, there would be a missing
zero-degree node $\delta$ adjacent to the support.
Its coefficient in $\rho$ is zero, while some adjacent
support coefficient is positive, so
$\langle\rho,\delta^\vee\rangle<0$.  The root-string
property implies that $\rho+\delta$ is a root.  Thus
$e_\delta$ does not preserve the line $\C e_\rho$,
contrary to the Borel-fixed condition.  Therefore
\[
                     \operatorname{Supp}(\rho)=C_\alpha.
\]
Taking $y=e_{-\rho}$ in the incidence equation gives
$\langle\lambda,\rho^\vee\rangle=0$.  The simple-coroot
expansion of $\rho^\vee$ has strictly positive
coefficients on its support.  Dominance of $\lambda$
therefore gives
$\langle\lambda,\delta^\vee\rangle=0$ for every
$\delta\in C_\alpha$.  Every root indexing
$\mathfrak b^+_\alpha$ has support in $C_\alpha$, so
its coroot also pairs trivially with $\lambda$.
Finite-dimensional root $\mathfrak{sl}_2$ theory implies
that every vector in $\mathfrak b^-_\alpha$ kills
$v_\lambda$.  Since $\mathfrak b^-_\alpha$ is a
$\mathfrak g_0$-module, commuting it past the
degree-zero lowerings in \eqref{eq:root-top-module}
shows that it kills all of $E$.  These root-string and
highest-weight facts are the usual semisimple
representation theory; see \cite[\S\S9--10,20--21]{Humphreys}.
Finally $x_\alpha^*\in\mathfrak b^-_\alpha$ for every
nonzero component, and summing proves $x^*E=0$.
\end{proof}

\section{The relative Iitaka base and its section lattice}
\label{sec:relative-iitaka}

Our goal is to construct an intermediate base $W$ over the marked-period
quotient $S$ and a divisor on $W$ whose complete divisible section spaces
are those of the original logarithmic base. The rank-one generator on
the generic Iitaka fiber will give this exact comparison.

The coefficient estimates are now available, as are the marked quotient
and both of its positivity comparisons. We will first obtain a nonzero
logarithmic form on the generic period fiber. Only after this
nonnegativity has been established will we form its relative Iitaka
fibration.

We use the marked quotient of \Cref{prop:marked-period}.  Resolve its
compactification and the rational quotient map, obtaining a morphism
\(q:Y_0\to S\) of smooth projective varieties.  Its geometric generic
fiber is connected.  Every source model used below is resolved together
with its reduced boundary, so that the strict transform of the preceding
boundary plus all reduced exceptional divisors is SNC.  A higher source
model always means such a log resolution.  We keep the subscript \(0\) for a fixed such source
model and denote its reduced log boundary by \(D_{Y_0}\); subsequent
models are denoted by \(Y\).  In this section \(E_0\)
denotes the highest Hodge bundle of the descended tensor system on a dense
open subset of \(S\).  The homogeneous vector of degree \(a>0\) is viewed
as a rational tensor
\[
 u_*\in aK_{Y_0}\otimes q^*E_0,
\]
with the logarithmic pole and higher-model properties of
\Cref{prop:logarithmic-vector}.  This notation for a rational tensor
allows the poles prescribed by \(aD_{Y_0}\).

We first record why these source modifications preserve the spaces of
logarithmic pluriforms.  If \(p:Y'\to Y_0\) is a log resolution and
\(D'\) is strict transform plus reduced exceptional divisor, then
\begin{equation}\label{eq:iitaka-log-discrepancy}
 K_{Y'}+D'=p^*(K_{Y_0}+D_{Y_0})+R_p,
 \qquad R_p\ge0\quad\text{and}\quad p_*R_p=0.
\end{equation}
For a blowup of a smooth center of codimension \(c\) contained in
\(e\) boundary branches, the new coefficient in \(R_p\) is \(c-e\ge0\).
For a general log resolution the formula follows from the log canonical
discrepancy inequality for the log smooth pair: each exceptional
coefficient is one plus its discrepancy and is therefore nonnegative.
A rational section of a line bundle on the smooth variety \(Y_0\) is regular if it
has no pole at any prime divisor.  Consequently an effective exceptional
divisor introduces no new rational sections, and, for every \(l\ge0\),
\begin{equation}\label{eq:iitaka-log-sections}
 H^0\bigl(Y',l(K_{Y'}+D')\bigr)
   =H^0\bigl(Y_0,l(K_{Y_0}+D_{Y_0})\bigr).
\end{equation}
These are identifications of rational pluriforms and therefore preserve
their ratios.  The same argument applies over the generic point of a
base.

\subsection{Constructing the relative Iitaka fibration}

Over the generic point of $S$, a nonzero scalar coordinate of $u_*$
will supply a logarithmic pluriform.  The relative Iitaka map then
separates its section ratios, following the classical relative
construction of \cite[\S6, Remark~1]{Iitaka1971}.  On its generic fiber
the nonzero degree-$a$ section space has rank one, so a single generator
factors the entire vector.  The exact comparison with the original
source's complete section spaces will require the valuation argument
that follows this construction.

\begin{lemma}[The relative Iitaka diagram]
\label{lem:relative-iitaka-construction}
For \(q:(Y_0,D_{Y_0})\to S\) and \(u_*\) as above, there are smooth
projective varieties \(Y,W\), a birational morphism \(p:Y\to Y_0\),
and morphisms \(x:Y\to W\) and \(h:W\to S\) with geometrically
connected generic fibers, such that \(h\circ x=q\circ p\).
Put \(D=p_*^{-1}D_{Y_0}+\operatorname{Exc}(p)_{\rm red}\) and
\(k=\kappa(Y_{0,\eta_S},K_{Y_{0,\eta_S}}+D_{Y_0,\eta_S})\).
Then \(k\ge0\), \(\dim W-\dim S=k\), and the geometric generic
fiber of \(x\) has logarithmic Iitaka dimension zero.  Write
\(J=Y_{\eta_W}\) for its generic fiber over \(\C(W)\).

The degree \(a\) used for \(u_*\) itself has a nonzero generator
\(\omega\in H^0(J,a(K_J+D_J))\).  There is a rational tensor
\(u_W\in aK_W\otimes h^*E_0\) for which
\begin{equation}\label{eq:iitaka-vector-factorization}
 u_*=\omega\,x^*u_W.
\end{equation}
\end{lemma}

\begin{figure}[ht]
\centering
\begin{tikzcd}[column sep=large,row sep=large]
Y \arrow[r,"x"] \arrow[d,"p"'] & W \arrow[r,"h"] & S\\
Y_0 \arrow[rru,"q"',bend right=15] &&
\end{tikzcd}
\caption{The relative logarithmic Iitaka diagram over the marked-period
quotient $S$, with $p$ birational.  The generic $x$-fiber $J$ has
$\kappa(J,K_J+D_J)=0$; its degree-$a$ generator $\omega$ gives
$u_*=\omega x^*u_W$.  The generic $h$-fiber has dimension
$k=\kappa(Y_{0,\eta_S},K_{Y_{0,\eta_S}}+D_{Y_0,\eta_S})$.}
\label{fig:intermediate-base}
\end{figure}


\begin{proof}
Over the generic point of \(S\) the coefficient bundle of \(u_*\) is
constant.  At any divisor horizontal over \(S\), its lattice is the
ordinary pullback lattice from the interior of \(S\).  The bounds of
\Cref{prop:logarithmic-vector} consequently give regular logarithmic
coefficients on \(Y_{0,\eta_S}\).  If the bounds were verified after
ramification, regularity descends: the order of a pulled-back ordinary
coefficient is the ramification index times its original order.
Choosing a basis of \((E_0)_{\eta_S}\) and trivializing the base
canonical line therefore gives a nonzero element of
\[
 H^0\bigl(Y_{0,\eta_S},a(K_{Y_{0,\eta_S}}+D_{Y_0,\eta_S})\bigr)
       \otimes_{\C(S)}(E_0)_{\eta_S}.
\]
At least one scalar component is nonzero, proving \(k\ge0\).

Work temporarily over \(K=\C(S)\).  Choose a complete logarithmic
pluricanonical system whose image has dimension \(k\), and resolve its
base ideal.  If its degree is \(b\), its morphism \(\phi:Z\to Z_b\)
satisfies
\begin{equation}\label{eq:iitaka-movable-system}
 b(K_Z+D_Z)\sim\phi^*H+F_b,
 \qquad F_b\ge0,
\end{equation}
where \(H\) is the hyperplane divisor on its projective image.
The extra logarithmic discrepancy in a source resolution is effective,
so the indicated inequality remains valid after further resolution.
Take the Stein factor of \(\phi\), resolve it and the source, and
spread the resulting diagram over a dense open subset of \(S\).
Compactification and resolution give \(Y\xrightarrow{x}W
\xrightarrow{h}S\).  The field \(\C(W)\) is the relative algebraic
closure of the system-image field in \(\C(Y)\); in particular it is
algebraically closed in \(\C(Y)\).  Since \(\C(S)\) is algebraically
closed in \(\C(Y)\), it is also algebraically closed in \(\C(W)\).
In characteristic zero these facts give geometrically connected
generic fibers for both maps.

Here is the section-ratio argument for \(\kappa(J,K_J+D_J)=0\).
The restriction of \(u_*\) still gives a nonzero section in degree
\(a\), exactly as above.  Suppose two independent logarithmic
pluriforms existed on \(J\) in some common degree \(c\).
Their ratio \(r\) would not lie in \(\C(W)\), and would be
transcendental over that field, by its relative algebraic closure.
Spread the two sections to rational sections of
\(c(K_Z+D_Z)\), trivializing the canonical directions from the base
rationally.  Their pole divisors are vertical over \(W_{\eta_S}\).
There are finitely many such prime divisors; the images of all of them
in \(Z_b\) are proper, since \(W_{\eta_S}\to Z_b\) is generically
finite.  Hypersurfaces on \(Z_b\) containing these images, with
sufficient multiplicity, give a section of some \(nH\) whose pullback
cancels all their poles.  Multiply also by the effective section of
\(nF_b\) in \eqref{eq:iitaka-movable-system}.  The two resulting
sections belong to
\(H^0(Z,(c+nb)(K_Z+D_Z))\) and retain ratio \(r\).
Together with the \(k\) independent ratios of the original system,
they give \(k+1\) independent ratios of logarithmic pluriforms.
All finitely many ratios can be placed in one degree by multiplying
numerators and denominators by the other denominators.  This
contradicts the definition of \(k\).

The same conclusion holds for the geometric generic fiber.  In every
fixed degree, proper cohomology in degree zero commutes with field
extension; two sections over an algebraic closure would therefore
already give section-space dimension at least two over \(\C(W)\).
Adjunction and generic smoothness identify the fiber line with
\(K_J+D_J\), and \eqref{eq:iitaka-log-sections} justifies all the log
source replacements.  This proves the claimed geometric Iitaka
dimension without an abundance assumption.

The degree-\(a\) space on \(J\) is now one-dimensional and nonzero.
Choose its generator \(\omega\) and extend it as a rational relative
canonical tensor.  Divide the vector \(u_*\) by it on the generic
fiber.  Its scalar coefficients are in \(\C(W)\), so they define
\(u_W\) and prove \eqref{eq:iitaka-vector-factorization} in the stated
degree \(a\).  The Higgs operator is \(\OO\)-linear; thus its iterates
satisfy the same factorization under differential pullback.  There is
no derivative of \(\omega\) in this assertion.  In particular the
projective direction \([u_W]\) gives precisely the same compact flags
as \([u_*]\) at general corresponding points.
\end{proof}

\subsection{The exact section lattice}

The factorization has so far been obtained over the function field.
To recover the complete logarithmic section spaces on $Y_0$, we must
determine which rational tensors on $W$ multiply by powers of $\omega$
to become regular on the source.  We must test divisors on the fixed
original source even when their images have codimension at least two
on a temporary base. The next lemma extracts precisely those finitely many
base valuations; subsequent source modifications add no new sections.

\begin{lemma}[Extraction of the source valuations]
\label{lem:valuation-preparation}
Let \(x_1:Y_1\to W_1\) be a dominant projective morphism of smooth
projective varieties with geometrically connected generic fiber, and fix
a reduced simple-normal-crossing divisor
\(D_1\) on \(Y_1\).  Write \(J=Y_{1,\eta_{W_1}}\), and suppose that its geometric
logarithmic Iitaka dimension is zero.  Choose, over \(\C(W_1)\),
\[
 0\ne\omega\in H^0\bigl(J,a(K_J+D_J)\bigr)
\]
for some \(a>0\).  Regard \(\omega\) as a rational section of
\(aK_{Y_1/W_1}\).  There is a smooth projective modification of \(W_1\)
with the following property.  On any further smooth base model \(W\),
take a smooth resolution \(p:Y\to Y_1\) of the main component, with
\(x:Y\to W\), and set \(D=p_*^{-1}D_1+\operatorname{Exc}(p)_{\rm red}\).
For a prime \(A\subset W\), put
\begin{equation}\label{eq:iitaka-minimum}
 \begin{aligned}
 t_A&=\min_{I\,\mapsto\,A}
       \frac{a^{-1}\operatorname{ord}_I(\omega)+d_I}{m_I},\\
 m_I&=\operatorname{ord}_I(x^*A),
       \qquad d_I=\operatorname{coeff}_I D,\\
 T_W&=\sum_A t_AA,\qquad L_W=K_W+T_W.
 \end{aligned}
\end{equation}
Here \(I\) runs through prime divisors dominating \(A\), and the order
of \(\omega\) is its order as a section of \(a(K_Y-x^*K_W)\).
The sum defining \(T_W\) is finite.  For every sufficiently divisible
\(l>0\), multiplication gives an isomorphism
\begin{equation}\label{eq:iitaka-exact-sections}
 \begin{split}
 H^0(W,lL_W)&\xrightarrow{\ \simeq\ }
       H^0\bigl(Y_1,l(K_{Y_1}+D_1)\bigr),\\
 \sigma&\longmapsto \omega^{l/a}x^*\sigma.
 \end{split}
\end{equation}
The right side is interpreted on the fixed model \(Y_1\), by rational
identification.  The isomorphisms respect multiplication in a common
Veronese subring and preserve ratios.
\end{lemma}

\begin{proof}
Write \(n=\dim Y_1\), \(w=\dim W_1\), and \(r=n-w\).  We explain the
valuation assertion needed for the modification.  Let \(I_1\) be a prime
of \(Y_1\) whose image has codimension at least two in \(W_1\), and let
\(v=\operatorname{ord}_{I_1}\).  Its restriction to \(\C(W_1)\) is
nontrivial: a local function vanishing on the image has positive order
at \(I_1\).  The nonzero value subgroup is \(m\Z\) for an integer
\(m>0\).  Let \(v_0\) be the normalized restricted valuation and let
\(\kappa(v)\) and \(\kappa(v_0)\) denote their residue fields.

Residues algebraically independent over \(\kappa(v_0)\) lift to elements
algebraically independent over \(\C(W_1)\).  Indeed, in a proposed
polynomial relation among lifts of value zero, divide all coefficients
by a coefficient of smallest valuation.  Reduction then gives a
nonzero polynomial relation among their residues.  It follows that
\[
 \operatorname{trdeg}_{\kappa(v_0)}\kappa(v)\le r,
 \qquad
 \operatorname{trdeg}_{\C}\kappa(v_0)\ge(n-1)-r=w-1.
\]
The reverse inequality is elementary: lifts of independent residues,
together with one element of positive valuation, are algebraically
independent over \(\C\), by examining the term of smallest valuation
in a polynomial relation.  Thus the residue transcendence degree is
exactly \(w-1\).

Choose \(w-1\) such independent residues and rational-function lifts on
\(W_1\).  Resolve the graph of their map to \((\mathbf P^1)^{w-1}\).
The center of \(v_0\) on this proper model has dimension at least
\(w-1\), since its image contains the generic point of the product;
it is proper since the valuation is nontrivial.  It is therefore a
prime divisor.  Resolving further makes the base smooth and does not
contract that center.  This proves that the restriction of \(v\) is a
positive integral multiple of a divisorial valuation of the base.

Only finitely many primes \(I_1\) need this treatment.  Such a prime
lies in the locus where the local fiber dimension of \(x_1\) is at
least \(r+1\).  This is a proper closed subset of \(Y_1\), and a prime
divisor contained in it must be one of its finitely many divisorial
components.  Extract all the restricted valuations simultaneously.
After any further base modification their centers are still divisors.
On a compatible source resolution, the strict transform of each
\(I_1\) survives and dominates its extracted center.  Every prime on
the fixed model \(Y_1\) now either dominates the base or has strict
transform dominating a base prime.

There are finitely many divisors in
\(\operatorname{div}_{aK_{Y/W}}(\omega)+aD\).  Outside the images of
their vertical components, all the numerators in
\eqref{eq:iitaka-minimum} vanish.  This proves finiteness of \(T_W\).
For a nonzero section \(s\) on the right of
\eqref{eq:iitaka-exact-sections}, pullback and
\eqref{eq:iitaka-log-discrepancy} first make \(s\) a logarithmic
pluriform on \((Y,D)\).  If \(a\mid l\), the section space on \(J\)
in degree \(l\) is exactly the line generated by \(\omega^{l/a}\):
it contains that element, and two independent sections would have a
nonconstant ratio, contradicting \(\kappa(J,K_J+D_J)=0\).
Hence, at the function-field level, there is a unique rational
\(l\)-canonical tensor \(\sigma\) on \(W\) such that
\(s=\omega^{l/a}x^*\sigma\).

Fix a base prime \(A\), take a regular local frame of \(lK_W\) at its
generic point, and let \(c\) be the coefficient of \(\sigma\) in that
frame.  For every \(I\) dominating \(A\), regularity of \(s\) says
\begin{equation}\label{eq:iitaka-order-test}
 m_I\operatorname{ord}_A(c)
       +\frac la\operatorname{ord}_I(\omega)+l d_I\ge0.
\end{equation}
These inequalities are equivalent to
\(\operatorname{ord}_A(c)+lt_A\ge0\).  After clearing denominators,
they say precisely that \(\sigma\in H^0(W,lL_W)\).

Conversely, suppose \(\sigma\in H^0(W,lL_W)\), and form the rational
pluriform \(s=\omega^{l/a}x^*\sigma\).  A horizontal prime of the fixed
source \(Y_1\) satisfies the logarithmic regularity test because
\(\omega^{l/a}\) is a logarithmic pluriform on \(J\).  The strict
transform of every other prime of \(Y_1\) dominates a prime \(A\) of
\(W\), by the extraction just performed.  The minimum rule implies
\eqref{eq:iitaka-order-test} there.  The boundary coefficient and the
order of an absolute rational pluriform at a surviving prime agree
with those on \(Y_1\).  Thus \(s\) is regular as a section of
\(l(K_{Y_1}+D_1)\) at every prime of the fixed model.  Smoothness, or
equivalently the normality extension criterion for a line bundle,
extends it across the remaining codimension-two subset.

This argument does not discard the primes contracted by \(x_1\): their
tests are exactly the reason for extracting their restricted valuations.
It also explains why no indefinite extraction of all new source
exceptional divisors is needed.  Once a section is regular on \(Y_1\),
\eqref{eq:iitaka-log-discrepancy} makes its pullback regular on every
higher log source.  Both maps in \eqref{eq:iitaka-exact-sections} are
inverse rational identifications, which proves the last assertions.
\end{proof}

\begin{proposition}[The rank-one base]
\label{prop:rank-one-base}
Take the diagram and degree-\(a\) generator \(\omega\) of
\Cref{lem:relative-iitaka-construction}.  After further smooth
birational preparation, with the reduced source boundary
\(D=p_*^{-1}D_{Y_0}+\operatorname{Exc}(p)_{\rm red}\), the divisor
\(L_W\) defined by \eqref{eq:iitaka-minimum} is big over \(S\).
In every sufficiently divisible degree it satisfies the exact
section comparison
\begin{equation}\label{eq:iitaka-fixed-source}
 H^0(W,lL_W)\simeq
 H^0\bigl(Y_0,l(K_{Y_0}+D_{Y_0})\bigr).
\end{equation}
Moreover, the models can be chosen so that
\begin{equation}\label{eq:iitaka-bundle-formula}
 L_W=K_W+T_W=K_W+B_W+M_W,
 \qquad 0\le B_W\le1,
 \qquad M_W\ \text{nef},
\end{equation}
where \(B_W\) has simple-normal-crossing support and the equality uses
compatible representatives of the rational divisor classes.  If \(I\)
attains the minimum at a prime \(A\), then
\begin{equation}\label{eq:iitaka-discriminant-bound}
 \operatorname{coeff}_A B_W
       \ge 1-\frac{1-d_I}{m_I}\ge0.
\end{equation}

After any further birational preparation of the diagram, the minimum
rule can be recomputed with the higher reduced log boundary.  The
section comparison with the same fixed \((Y_0,D_{Y_0})\) still holds, and
the moduli divisor is the trace of the same b-divisor.  Once that
b-divisor descends to the chosen model, all these higher traces are
pullbacks of \(M_W\).  A further log resolution restores the
simple-normal-crossing condition on \(B_W\).
\end{proposition}

\begin{proof}
\emph{The exact lattice and relative bigness.}
Fix an initial smooth source model \(Y_1\to Y_0\) on which the just
constructed Iitaka map is a morphism.  Apply
\Cref{lem:valuation-preparation} to \(Y_1\), and thereafter use only
models dominating its extracted base.  That lemma, composed with
\eqref{eq:iitaka-log-sections}, proves
\eqref{eq:iitaka-fixed-source} with the original fixed \(Y_0\).
Its order argument also applies over \(\eta_S\).  The logarithmic
systems on \(Y_{\eta_S}\) have image dimension \(k\), and their
ratios descend to systems of \(L_W|_{W_{\eta_S}}\).  After replacing
a degree by a multiple to clear denominators, those descended ratios
still have transcendence degree \(k=\dim W_{\eta_S}\).  Thus \(L_W\)
is big over \(S\).  This remains true when \(k=0\), with the usual
meaning of bigness on a zero-dimensional generic fiber.

\emph{The auxiliary subpair.}
On the current smooth diagram define
\begin{equation}\label{eq:iitaka-auxiliary-pair}
 \Delta=-\frac1a\operatorname{div}_{aK_{Y/W}}(\omega)+x^*T_W.
\end{equation}
Choose compatible rational canonical forms, and write \(\varphi\) for
the rational function expressing \(\omega\) in the corresponding
relative frame.  Then
\begin{equation}\label{eq:iitaka-trivialization}
 K_Y+\Delta+\frac1a\operatorname{div}(\varphi)=x^*L_W.
\end{equation}
At a horizontal prime the inequality \(\Delta\le D\) follows from
the logarithmic regularity of \(\omega|_J\).  At every prime
dominating a fixed base prime \(A\), it is exactly the minimum
inequality \eqref{eq:iitaka-minimum}.  Divisors mapping into
codimension two disappear after shrinking around \(\eta_A\).
It follows that \(\Delta\le D\) near \(x^{-1}(\eta_A)\) for every
\(A\), and also over the generic point of \(W\).  Since the reduced
log smooth pair \((Y,D)\) is log canonical, these inequalities imply
that \((Y,\Delta)\) is sub-log-canonical there.  Negative coefficients
of \(\Delta\) cause no difficulty.

We verify the discrepancy-sheaf rank condition rather than inferring
it from generic log canonicity.  On a resolved geometric generic
fiber \(\rho:J'\to J\), use the higher reduced boundary \(D_{J'}\)
and put
\[
 Z_{J'}=\operatorname{div}_{aK_{J'}}(\omega)+aD_{J'}\ge0.
\]
The crepant transform \(\Delta_{J'}\) of the auxiliary subpair is
\[
 \Delta_{J'}=D_{J'}-\frac1a Z_{J'}.
\]
Moreover \(Z_{J'}\sim a(K_{J'}+D_{J'})\), so
\(\kappa(J',Z_{J'})=0\).  For the discrepancy b-divisor
\(\mathbf A^*\), the coefficient \(-1\) is omitted.  On this
resolution its rounded trace
\(R=\lceil\mathbf A^*_{J'}\rceil\) is effective and is supported
on \(\Supp Z_{J'}\): where the logarithmic zero divisor vanishes,
the auxiliary coefficient is either \(0\) or \(1\), and its
contribution to \(R\) is zero.  Therefore \(R\le N Z_{J'}\) for
some integer \(N\), and
\[
 h^0(J',\OO(R))\le h^0(J',\OO(NZ_{J'}))=1.
\]
If \(Z_{J'}=0\), this simply says \(R=0\).  Sections of the
discrepancy sheaf satisfy in particular the order conditions on this
resolution, so the displayed bound is also an upper bound for its
generic rank.  Conversely, on every higher model sub-log-canonicity
gives discrepancies at least \(-1\); after omitting \(-1\), their
ceilings are nonnegative.  The constant function \(1\) consequently
belongs to the discrepancy sheaf.  We have proved
\begin{equation}\label{eq:iitaka-rank-condition}
 \operatorname{rank}x_*\OO_Y
       \bigl(\lceil\mathbf A^*(Y,\Delta)\rceil\bigr)=1.
\end{equation}

The morphism \(x\) has connected fibers: its Stein factor is finite
birational over the normal variety \(W\), since its geometric generic
fiber is connected.  Thus \(x\) is projective and surjective with connected fibers,
its source and base are normal, \(\Delta\) is rational, the subpair is
sub-log-canonical generically, and
\eqref{eq:iitaka-trivialization} and
\eqref{eq:iitaka-rank-condition} give its remaining lc-trivial
conditions.  Thus it is an lc-trivial fibration in the precise sense of
\cite[Definition~3.2]{FG}.  By \cite[Theorem~3.6]{FG}, applied with
the structural morphism \(W\to\operatorname{Spec}\C\), its moduli
b-divisor is b-nef.  This application requires neither effectivity of
the auxiliary subpair nor abundance or semiampleness of the moduli
part.

\emph{The discriminant and higher models.}
Following \cite[\S3.4]{FG}, define
\[
 c_A=\sup\{c\in\Q:(Y,\Delta+c x^*A)
                  \text{ is sub-log-canonical over }\eta_A\},
 \qquad b_A=1-c_A.
\]
The comparison with \(D\) already proved gives \(c_A\ge0\), hence
\(b_A\le1\).  If \(I\) attains the minimum in
\eqref{eq:iitaka-minimum}, its coefficient in \(\Delta\) is exactly
\(d_I\).  The discrepancy condition at that prime forces
\(d_I+c_A m_I\le1\).  Therefore
\[
 1\ge b_A\ge1-\frac{1-d_I}{m_I}\ge0,
\]
where the last inequality uses \(d_I\in\{0,1\}\) and \(m_I\ge1\).
Put \(B_W=\sum_A b_AA\) and \(M_W=T_W-B_W\), using the
normalization \eqref{eq:iitaka-trivialization}.  This is the moduli
trace for that normalization.

It remains to reconcile b-nefness on a higher model with recomputation
of our minimum rule.  Let \(\beta:W'\to W\) be a smooth birational
base modification, take a compatible source resolution
\(\rho:Y'\to Y\), and use the higher reduced boundary \(D'\).
Write \(L'=K_{W'}+T_{W'}\) for the recomputed minimum divisor and
\(\Delta'\) for \eqref{eq:iitaka-auxiliary-pair} on that diagram.
Let \(\Delta^{\rm cr}\) be the crepant transform, defined by
\(K_{Y'}+\Delta^{\rm cr}=\rho^*(K_Y+\Delta)\).  Comparing
\eqref{eq:iitaka-trivialization} on the two models, with the same
rational trivialization, gives
\begin{equation}\label{eq:iitaka-normalization-change}
 \Delta'=\Delta^{\rm cr}+x'^*Q,
 \qquad Q=L'-\beta^*L_W.
\end{equation}
This identity includes the change of the relative canonical frame;
omitting that change would give an incorrect formula for \(T_{W'}\).

Let \(B^{\rm cr}_{W'}\) be the discriminant of the crepant induced
fibration, as in \cite[\S3.3]{FG}.  At a prime of \(W'\), adding
\(x'^*Q\) decreases its threshold by the local coefficient of
\(Q\).  Thus the recomputed discriminant and moduli divisors satisfy
\begin{equation}\label{eq:iitaka-moduli-compatibility}
 \begin{split}
 B_{W'}&=B^{\rm cr}_{W'}+Q,\\
 M_{W'}&=L'-K_{W'}-B_{W'}
       =\beta^*L_W-K_{W'}-B^{\rm cr}_{W'}.
 \end{split}
\end{equation}
The latter is precisely the old moduli b-divisor's trace.  Hence we
may first dominate a model on which it is nef and descends, and then
make any necessary further modifications; its new trace is its nef
pullback.  Repeating the minimum and threshold argument on each new
diagram still proves \(0\le B_{W'}\le1\).

For completeness, at a base prime not exceptional for \(\beta\) the
minimum does not change.  Indeed, the crepant transform of the old
auxiliary divisor is bounded by \(D'\) over its generic point, by
\eqref{eq:iitaka-log-discrepancy}; this gives the old minimum as a
lower bound for all new candidates.  An old minimizing prime survives
strictly and supplies equality.  The discriminant is likewise
unchanged there.  Thus any new discriminant support is contained in
the exceptional locus together with the transform of the old support.
Resolve that union.  Recomputing the minimum on the resulting model
therefore gives simple-normal-crossing support and keeps the moduli
part nef.  The extraction in \Cref{lem:valuation-preparation} survives
all these operations, so the exact section comparison continues to
refer to the same fixed \(Y_0\).  This proves all the assertions.
\end{proof}

All arguments in this section use the logarithmic vector constructed
above and the published lc-trivial b-nefness theorem.  The logarithmic
Iitaka lower bound of \Cref{thm:log-iitaka} is not used.

\section{Base nonvanishing and removal of the period twist}
\label{sec:nonvanishing}

It remains to prove the section construction stated in the introduction.
The marked quotient gives two cases. A point quotient supplies scalar
coordinates directly; for a positive-dimensional quotient the exact
section lattice of \Cref{sec:relative-iitaka} lets us work on its
intermediate base. There we remove the divisor added by period
positivity. Every argument here is independent of logarithmic
subadditivity.

\begin{proposition}\label{prop:nonvanishing}
Let $f:(X,E)\to(Y,D)$ satisfy the hypotheses of
\Cref{thm:upper-bound}, and let
\[
  0\ne s\in H^0\bigl(X,m(K_X+E)\bigr),\qquad m>0.
\]
Then $\kappa(Y,L_Y)\ge0$, where $L_Y=K_Y+D$.
More precisely, apply the homogeneous highest-vector and marked-period
constructions to $s$, and denote their connected quotient by $S$.
If $\dim S>0$, then the relative Iitaka space $W$ of
\Cref{lem:relative-iitaka-construction,prop:rank-one-base} satisfies
\[
                 \kappa(Y,L_Y)\ge\dim W\ge\dim S>0.
\]
If $S$ is a point, there is an integer $a>0$ and a nonzero section
$\tau\in H^0(Y,aL_Y)$.  Moreover, for every $y\in V$ at which
$s|_{X_y}=0$, this section can be chosen with $\tau(y)=0$.
\end{proposition}

The construction is attached to the chosen section $s$; in particular,
its quotient $S$ may depend on $s$.  We first complete the point-quotient
case, using the original coefficient lattice.  The rest of the section
proves that a positive-dimensional quotient forces positive base
Kodaira dimension.

\subsection{The point-quotient case}

\begin{proof}[Proof of \Cref{prop:nonvanishing} when $S$ is a point]
For a point base there is a nonzero constant base section and the zero
test has no nonzero input satisfying its premise.  Assume $\dim Y>0$.
Apply \Cref{prop:logarithmic-vector} and
\Cref{prop:marked-period} to the chosen $s$, obtaining a homogeneous
vector with source $-aL_Y$.

\phantomsection\label{proof:point-nonvanishing}
Since $S$ is a point, the descended coefficient system is constant.  In
a constant frame its coordinates are rational sections of $aL_Y$.
The divisorial lattice bounds make each coordinate regular at every
prime of $Y$: this can be tested after finite ramified substitution,
since a negative order remains negative after multiplying it by the
ramification index.  On the smooth variety $Y$, regularity in
codimension one gives global regularity.  At least one coordinate
is nonzero because the homogeneous vector is nonzero; choose it as
$\tau$.

If $s|_{X_y}=0$ for $y\in V$, \Cref{lem:zero-test} says that all
these coordinates have positive order in the pulled-back original
$aL_Y$ frame along the divisor of the blowup of $y$ (along $y$ itself
if $Y$ is a curve).  Finite ramification multiplies orders by its
positive index.  Hence each original regular coefficient has positive
exceptional order, equal to its vanishing order at $y$.  Thus the same nonzero coordinate
vanishes at $y$, as asserted.
\end{proof}

\subsection{The numerical target}

Assume henceforth that $\dim S>0$.
After the birational preparations in
\Cref{lem:relative-iitaka-construction,prop:rank-one-base}, write
\[
  Y\xrightarrow{x}W\xrightarrow{h}S,\qquad
  u_*=\omega\,x^*u_W,
  \qquad L_W=K_W+T_W\sim_{\Q}K_W+B_W+M_W.
\]
Here $W,S$ are smooth and projective, both maps have geometrically
connected generic fibers, $0\le B_W\le1$ has SNC support, and $M_W$ is
nef.  The divisor $L_W$ is big over $S$.  The highest vector has source
$-aL_Y$, and $\omega$ is the generator in degree $a$ on the generic
$x$-fiber.  All references to $Y$ in this diagram use its prepared model;
the actual pluriform comparison in \Cref{prop:rank-one-base} returns
sections to the original pair.

\begin{proposition}[Removal of a period twist]\label{prop:twist-removal}
Let $h:W\to S$ be a surjective morphism of smooth projective
varieties, with $\dim S>0$. Let $L$ be a rational divisor on $W$
and $H$ a nef rational divisor on $S$. Suppose:
\begin{enumerate}
\item $L$ is big on the generic fiber of $h$, and
      $L-h^*K_S$ is pseudoeffective;
\item $K_S+bH$ is big for some rational $b\geq0$;
\item there are a dominant generically finite morphism
      $\pi:\widehat W\to W$ from a smooth projective variety,
      a nef rational divisor $P$ on $\widehat W$, and a rational
      $Q>0$ such that, with $\widetilde H=\pi^*h^*H$,
      \[
      \nu(P+\widetilde H)=\nu(P),\qquad Q\pi^*L\ge_{\mathrm{pe}} P.
      \]
\end{enumerate}
Then $L$ is big.
\end{proposition}

The first two hypotheses make $L+b h^*H$ big. The numerical-dimension
condition in the third says that the added period directions contribute
no positivity beyond $P$. Its fixed comparison with $L$ allows us to
remove the added divisor. We prove this numerical implication first;
the geometric comparisons that follow will establish its hypotheses
for $L_W$ and $\mathcal H$.

\begin{lemma}[Nef subtraction]\label{lem:nef-morse}
Let $Z$ be smooth projective of dimension $n>0$, let $A$ be an ample
rational divisor, and let $P,H$ be nef rational divisors satisfying
$\nu(P+H)=\nu(P)$.  For every fixed rational $b\ge0$, the divisor
$A+tP-bH$ is big for all sufficiently large rational $t$.
\end{lemma}

\begin{proof}
Put $r=\nu(P)$.  The polynomial
\[
                      (A+tP)^n
\]
has degree $r$ and strictly positive leading coefficient.  In the
polynomial $n b(A+tP)^{n-1}H$, all coefficients of $t^i$ for $i\ge r$
vanish.  Indeed
\[
 (P+H)^{i+1}A^{n-1-i}=0\quad(r\le i\le n-1),
\]
and each mixed nef intersection in this expansion is nonnegative;
in particular $P^iHA^{n-1-i}=0$.  When $r=n$, this index range is empty.
For $r=0$ this says that the entire
second polynomial is zero.  Thus, for large $t$,
\begin{equation}\label{eq:nonvanishing-morse}
                   (A+tP)^n>n(A+tP)^{n-1}bH.
\end{equation}

We recall directly why this strict nef Morse inequality implies
bigness here.  Write $U=A+tP$ and $V=bH$.  The first divisor is ample
and the second is nef.  Replace $V$ by $V+\delta A$, with rational
$\delta>0$ sufficiently small that the strict inequality persists.
After multiplying by a common positive integer, $U$ is integral ample
and the new $V$ is very ample.  Choose a general member $H_V\in|V|$.
Successive restriction exact sequences give
\[
 h^0\bigl(Z,l(U-V)\bigr)
   \ge h^0(Z,lU)
      -\sum_{j=0}^{l-1}
             h^0\bigl(H_V,(lU-jV)|_{H_V}\bigr).
\]
Multiplication by a nonzero section of $jV|_{H_V}$ bounds each
summand by $h^0(H_V,lU|_{H_V})$.  On a curve choose this section
nonvanishing at the finitely many points of $H_V$; in higher dimension
take $H_V$ smooth and irreducible.  Ample Hilbert asymptotics therefore
give
\[
 h^0\bigl(Z,l(U-V)\bigr)
 \ge\frac{U^n-nU^{n-1}V}{n!}\,l^n+O(l^{n-1}).
\]
The coefficient is positive.  Hence $U-V$ is big, and adding back
$\delta A$ proves bigness before perturbation as well.
\end{proof}

\begin{proof}[Proof of Proposition~\ref{prop:twist-removal}]
Choose $b\geq0$ with $K_S+bH$ big and put $D_b=L+b h^*H$.
By hypothesis \textup{(i)} and bigness of $K_S+bH$, there are an ample rational $A_S$ and
$\delta>0$ rational such that $D_b\ge_{\mathrm{pe}}\delta h^*A_S$.
The divisor $D_b$ is also big over $S$. Hence
$D_b+c h^*A_S$ is big for some rational $c>0$.

To justify the latter assertion, fix an ample $A_W$ on $W$.
Relative bigness gives a nonzero generic-fiber section of
$lD_b-A_W$ for some divisible $l$. Its direct image is nonzero;
after twisting by $kA_S$ for large $k$ it has a nonzero global
section. Thus $lD_b-A_W+k h^*A_S$ is effective. Taking $c=k/l$
proves the assertion. Now
\[
 D_b=\frac{\delta}{c+\delta}(D_b+c h^*A_S)
      +\frac{c}{c+\delta}(D_b-\delta h^*A_S)
\]
expresses $D_b$ as a positive multiple of a big class plus a
pseudoeffective class. Therefore $D_b$ is big.

Pull back and choose an ample rational $A'$ on $\widehat W$ with
$\pi^*L+b\widetilde H\ge_{\mathrm{pe}} A'$. For every $t>0$,
hypothesis \textup{(iii)} gives
\[
 (1+tQ)\pi^*L\ge_{\mathrm{pe}} A'+tP-b\widetilde H.
\]
The right side is big for large rational $t$ by
the numerical-dimension equality in \textup{(iii)} and
Lemma~\ref{lem:nef-morse}. Thus
$\pi^*L$, and hence $L$, is big. Bigness descends under a dominant
generically finite projective map, for example by the norm of a
section after subtracting a small ample class. No equality of
numerical and Iitaka dimensions for a general pseudoeffective
divisor is being assumed.

\end{proof}

\subsection{The fixed determinant lines and their cotangent maps}

Choose the connected unipotent level and equivariant resolutions of
\Cref{prop:tail-rank}, and write their composite as
$\pi:\widehat W\to W$.  This map is generically finite and Galois at the
function-field level.  On $\widehat W$ let $G_i$ be the extended subbundles
generated by contractions of $\theta^i(u_W)$, put $r_i=\rk G_i$, and retain
only indices with $r_i>0$.  Set
\begin{equation}\label{eq:nonvanishing-tail}
 P=-\sum_i(i+1)\det G_i,\qquad
 \widetilde{\mathcal H}=\pi^*h^*\mathcal H.
\end{equation}
Both lines are nef and
\begin{equation}\label{eq:nonvanishing-rank}
             \nu(P+\widetilde{\mathcal H})=\nu(P).
\end{equation}
The projective line $[u_W]$ gives precisely the compact-factor marks
required in that proposition, by the factorization
$u_*=\omega x^*u_W$.

Fix a rational frame of each $\det G_i^\vee$ over the original function
field $\C(W)$.  Its divisor in the extended line upstairs is denoted by
$\widehat D_i^G$.  Define rational divisors on $W$ by
\begin{equation}\label{eq:nonvanishing-determinants}
 D_i^G=\frac{\pi_*\widehat D_i^G}{\deg\pi},\qquad
 P_W=\sum_i(i+1)D_i^G.
\end{equation}
The superscript distinguishes these determinant divisors from the
boundary $D$ on $Y$.  The rational frames and these divisors are fixed
throughout the ample perturbations below.

We define the cotangent lattice before writing the transfer map. For
the rational SNC boundary $B_W=\sum_A b_AA$, its adapted cotangent
lattice $\Omega^1(W,B_W)$ has, at a general point of $A=(t=0)$,
transverse generator $t^{1-b_A}dt/t$ and ordinary tangential generators.
These are interpreted on a cover $t=v^e$ with $eb_A$ integral. Thus
$b_A=0$ retains the pullback of $dt$, and $b_A=1$ retains the pullback
of $dt/t$. An adapted cover is chosen with these divisible ramification
indices along the boundary; any auxiliary branch divisors have boundary
coefficient zero and retain the pullback of ordinary forms. Tensor maps
into this lattice are statements on such covers, in codimension one.
The proof will also verify their extension across the remaining locus.

\begin{lemma}[Orbifold transfer]\label{lem:orbifold-transfer}
The divisor $P_W$ is pseudoeffective and
\begin{equation}\label{eq:nonvanishing-pushdown}
                    \pi^*P_W\ge_{\mathrm{pe}}P.
\end{equation}
For each $i$, the determinant of the $i$th Higgs map gives a generically
nonzero rational map
\begin{equation}\label{eq:nonvanishing-orbifold-map}
 \OO_W(D_i^G-ar_iL_W)
       \longrightarrow \bigl(\Omega^1(W,B_W)\bigr)^{\otimes ir_i}.
\end{equation}
This notation means the following precise assertion: after a fixed
tensor power clearing the rational source divisor, the pullback map is
regular on every sufficiently divisible adapted cover, using the
orbifold cotangent lattice in codimension one.  The same tensor power
works after replacing $B_W$ by any effective SNC boundary
$B_W+\Lambda$ with coefficients at most one.
\end{lemma}

\begin{proof}
The extending lattices and their determinant divisors are invariant
under the Galois group.  For every prime $A_0$ of $W$, the orders of the
pullback of $D_i^G$ therefore agree with those of $\widehat D_i^G$ at all
primes over $A_0$.  Indeed, if $e$ is the common ramification index and
$v$ their common order, then the coefficient downstairs obtained by
pushforward and division by $\deg\pi$ is $v/e$.

Factor $\pi$ through the finite normalization $W^{\mathrm{fin}}\to W$.
The difference
\[
                    Z=\pi^*P_W-\sum_i(i+1)\widehat D_i^G
\]
is exceptional over $W^{\mathrm{fin}}$.  Its negative is relatively nef:
the pullback term has degree zero on contracted curves, and the remaining
term represents the nef line $P$.  The negativity lemma
\cite[Lemma~3.39]{KollarMori} gives $Z\ge0$.  This proves
\eqref{eq:nonvanishing-pushdown}.  Pushforward of the pseudoeffective
class $P$ gives $\deg(\pi)P_W$, proving the first assertion.

We prove the lattice assertion at an arbitrary prime $A_0\subset W$.
Use the notation of \Cref{prop:rank-one-base}:
\[
 t_{A_0}=\min_{I\to A_0}
       \frac{a^{-1}\operatorname{ord}_I(\omega)+d_I}{m_I},
 \quad m_I=\operatorname{ord}_I(x^*A_0),
 \quad d_I=\operatorname{coeff}_I(D).
\]
Choose a component $I$ achieving this minimum.  The discriminant
coefficient $b_{A_0}=\operatorname{coeff}_{A_0}(B_W)$ satisfies
\begin{equation}\label{eq:nonvanishing-same-minimizer}
 t_{A_0}=\frac{a^{-1}\operatorname{ord}_I(\omega)+d_I}{m_I},
 \qquad b_{A_0}\ge1-\frac{1-d_I}{m_I}.
\end{equation}
Both statements use this same $I$.

At general points of $A_0$ and $I$, choose a parameter $t$ downstairs
and a parameter $z$ upstairs.  After an etale local change absorbing a
unit, $t=z^{m_I}$; the other base coordinates pull back to independent
coordinates along $I$.  Work on a local cover $t=v^e$ sufficiently
divisible to make the Hodge monodromy unipotent and clear all relevant
denominators.  Write a coefficient $c$ of $\theta^i(u_W)$ in Schmid
frames, using an ordinary frame of $aK_W$ and a cotangent tensor having
$j$ factors $dt/t$, with the other factors tangential.  Define
$\operatorname{ord}(c)$ using the parameter $v$.

The identity $u_*=\omega x^*u_W$ holds for Higgs iterates as well:
the Higgs field is linear over functions, so there is no derivative of
$\omega$.  Make a further divisible substitution in $z$ dominating
the chosen $t$-cover.  The lattice bounds of
\Cref{prop:logarithmic-vector,lem:weight-extension} then apply to each
coefficient in the independent pulled-back coordinate tensors.
Each transverse factor pulls back as $m_I\,dz/z$ and costs
$1-d_I$ units of $I$-order.  Consequently
\[
 m_I\frac{\operatorname{ord}(c)}{e}
       +\operatorname{ord}_I(\omega)+ad_I-j(1-d_I)\ge0.
\]
The order of $\omega$ is measured in $aK_{Y/W}$, so the canonical
change of variables is already included.  Using
\eqref{eq:nonvanishing-same-minimizer}, we obtain
\begin{equation}\label{eq:nonvanishing-orbifold-order}
 \frac{\operatorname{ord}(c)}e
       \ge-at_{A_0}+\frac{j(1-d_I)}{m_I}
       \ge-at_{A_0}+j(1-b_{A_0}).
\end{equation}
This is exactly the required regularity after twisting the source by
$-aL_W$: the transverse orbifold cotangent frame is
$t^{1-b_{A_0}}dt/t$, interpreted after ramification, and the tangential
frames are ordinary differentials.  For example, when $b_{A_0}=0$
this frame is $dt$, and when $b_{A_0}=1$ it is $dt/t$.

Generically the contractions of this map surject onto $G_i$.  On the
unipotent resolution $G_i$ has its subbundle lattice inside the Schmid
grade, so the regularity just proved holds with this target lattice.
Dualize, take the top exterior power of $G_i^\vee$, and use the natural
inclusion of an exterior power in a tensor power.  The resulting source
is $\det G_i^\vee-ar_i\pi^*L_W$, which gives
\eqref{eq:nonvanishing-orbifold-map} by the order identification for
$D_i^G$ at primes over $W$.

Here is also a global interpretation of the cover assertion.  Fix an
integer $N_0>0$ clearing all the source divisors
$D_i^G-ar_iL_W$.  On a smooth adapted cover $\rho:W_{\rm ad}\to W$
for $B_W+\Lambda$, the $N_0$th tensor power is a rational map from
$\rho^*\OO_W(N_0(D_i^G-ar_iL_W))$ to the tensor power of degree
$N_0ir_i$ of its orbifold cotangent bundle.  At each prime it is regular
after a further common local ramification that is also unipotent for
the Hodge system.  Fractional cotangent frames pull back compatibly
under that ramification; an inequality of integral orders after such
a finite substitution implies the original inequality.  Thus the map
was already regular on $W_{\rm ad}$.  Regularity in codimension one
extends across codimension two because source and target are locally
free on its smooth model.

Increasing the boundary coefficient enlarges the cotangent lattice:
the original transverse frame is the new one multiplied by
$t^{\operatorname{coeff}_{A_0}(\Lambda)}$.  At auxiliary branch
divisors with zero boundary coefficient the lattice is the pullback
of ordinary cotangent forms.  Hence the same argument applies at every
prime of any adapted cover for the larger boundary.  Only the
denominators of this cover change with $\Lambda$; the integer $N_0$
is the fixed integer chosen above.
\end{proof}

\subsection{The uniform threshold and relative positivity}

We next turn these maps into a numerical comparison.  We use
\cite[Theorem~1.3]{CP} in its orbifold tensor form: for a smooth
projective SNC pair $(Z,\Delta)$ with rational coefficients in $[0,1]$
and $K_Z+\Delta$ pseudoeffective, every torsion-free quotient of an
orbifold cotangent tensor power on an adapted cover has nonnegative
degree against the pullback of every movable curve class on $Z$.
In particular, if a line injects into such a tensor power, its degree
is at most the degree of the determinant of that tensor power.  To
see the last assertion, saturate the image and apply the theorem to
the quotient; the effective divisor introduced by saturation only
strengthens the assertion.  For a rank-one tensor it follows directly
from the effective zero divisor of the map.

The threshold argument below follows the ample-perturbation method in
\cite[proof of Theorem~7.6]{CP}; here the fixed Higgs maps provide
constants independent of the varying boundary.

\begin{lemma}[A uniform pseudoeffective threshold]
\label{lem:pseff-threshold}
In the notation above, $L_W$ is pseudoeffective.  There is a constant
$Q>0$, depending only on the fixed Higgs maps, such that
\begin{equation}\label{eq:nonvanishing-domination}
                         Q\pi^*L_W\ge_{\mathrm{pe}}P.
\end{equation}
\end{lemma}

\begin{proof}
Fix an ample rational divisor $A$ on $W$ and write $n=\dim W>0$.
For a positive rational $\epsilon$ for which $L_W+\epsilon A$ is
pseudoeffective, the divisor $M_W+\epsilon A$ is ample.  Choose a
sufficiently divisible $k$ and a general smooth member of
$|k(M_W+\epsilon A)|$, transverse to every stratum of $B_W$, and divide
it by $k$.  This gives
\[
  \Lambda_\epsilon\sim_{\Q}M_W+\epsilon A,
  \quad 0\le B_W+\Lambda_\epsilon\le1,
  \quad K_W+B_W+\Lambda_\epsilon\sim_{\Q}L_W+\epsilon A.
\]
The support is SNC.  Thus the pair is log canonical and has
pseudoeffective adjoint, as required in \cite[Theorem~1.3]{CP}.

Choose a smooth Galois adapted cover for this pair, as in
\cite[Definition~5.1 and Lemma~5.2]{CP}.  Such covers can be
obtained by the covering construction: add general transverse
auxiliary divisors to make the root divisors divisible in the Picard
group, and normalize the resulting finite root covers.  Use sufficiently
many general auxiliary choices that, at each point and for each
prescribed boundary component, one auxiliary equation is a unit.
Locally the normalization then extracts roots of independent SNC
coordinates and is smooth.  The auxiliary ramifications are assigned
boundary coefficient zero.  On this cover the determinant of the
orbifold cotangent bundle is the pullback of
$K_W+B_W+\Lambda_\epsilon$.

Put $d_i=ir_i$.  For $d_i>0$, the tensor in
\Cref{lem:orbifold-transfer} has rank $n^{N_0d_i}$ and first Chern
class
\[
 N_0d_i\,n^{N_0d_i-1}
             \rho^*(L_W+\epsilon A).
\]
Apply the preceding quotient statement to its nonzero source line.
Divide the resulting inequality by $N_0\deg\rho$.  For every movable
curve class $\alpha$ on $W$ this gives
\[
 \bigl(c_i(L_W+\epsilon A)+ar_iL_W-D_i^G\bigr)\cdot\alpha\ge0,
 \qquad c_i=d_i n^{N_0d_i-1}.
\]
For $d_i=0$, the map is into the trivial line; take $c_i=0$ and use
its effective zero divisor.  All $c_i$ are independent of $\epsilon$
and of the chosen representative $\Lambda_\epsilon$.

The duality between the pseudoeffective divisor cone and the movable
curve cone on the smooth projective variety $W$
\cite[Theorem~2.2]{BDPP} now gives
\begin{equation}\label{eq:nonvanishing-uniform-inequality}
 C(L_W+\epsilon A)+RL_W\ge_{\mathrm{pe}}P_W,
 \qquad
 C=\sum_i(i+1)c_i\ge0,
 \quad R=a\sum_i(i+1)r_i>0.
\end{equation}
The strict positivity of $R$ follows from $G_0\ne0$.

For completeness, consider
\[
 \tau=\inf\{t\ge0:L_W+tA\text{ is pseudoeffective}\}.
\]
This is finite.  Since $P_W$ is pseudoeffective,
\eqref{eq:nonvanishing-uniform-inequality} implies, for every rational
$\epsilon>\tau$,
\[
        L_W+\frac{C}{C+R}\epsilon A
                         \text{ is pseudoeffective}.
\]
If $\tau>0$, closedness of the cone and $\epsilon\downarrow\tau$
would give $\tau\le C\tau/(C+R)<\tau$, a contradiction.  Thus
$\tau=0$ and $L_W$ is pseudoeffective.  We may now let
$\epsilon\downarrow0$ in
\eqref{eq:nonvanishing-uniform-inequality}; pullback and
\eqref{eq:nonvanishing-pushdown} yield
\eqref{eq:nonvanishing-domination} with $Q=C+R$.
\end{proof}

\Needspace{8\baselineskip}
\begin{lemma}[Relative pseudoeffectivity]\label{lem:nonvanishing-relative}
For the prepared morphism $h:W\to S$ one has
\begin{equation}\label{eq:nonvanishing-relative}
                         L_W-h^*K_S\ge_{\mathrm{pe}}0.
\end{equation}
\end{lemma}

\begin{proof}
If $h$ has relative dimension zero, it is birational by connectedness
of its generic fiber, and the conclusion follows from
$K_W-h^*K_S\ge0$, $B_W\ge0$, and nefness of $M_W$.
Assume henceforth that its relative dimension is positive.
Use $A$ and $\Lambda_\epsilon$ from the preceding proof.  We have
established pseudoeffectivity of
$K_W+B_W+\Lambda_\epsilon\sim_{\Q}L_W+\epsilon A$ for every rational
$\epsilon>0$.  Apply \cite[Theorem~3.4 and Remark~3.3]{CP} to the
fibration and this smooth SNC log canonical pair.  We spell out the
model issue because the fixed $W$ need not itself be the good model
used in that theorem.

Take a commutative birational diagram
\[
 \begin{array}{ccc}
  W'&\xrightarrow{h'}&S'\\
  {\scriptstyle p}\downarrow&&\downarrow{\scriptstyle q}\\
  W&\xrightarrow{h}&S,
 \end{array}
\]
with smooth projective varieties, resolved morphism, SNC discriminant
and reduced inverse discriminant, and such that every divisor of $W'$
contracted by $h'$ to codimension at least two is exceptional over
$W$.  These are the good-model requirements in
\cite[Remark~3.3]{CP}.  They may be arranged by extracting on $S$ the
restricted valuations of the finitely many divisors of $W$ with
exceptional image, as in \Cref{lem:valuation-preparation}, and then
resolving the map and the boundary.  Once those original valuations
have divisorial center, further base modifications preserve that
property for their strict transforms; all remaining offending source
divisors are exceptional over $W$.

One can require the discriminant here to be the actual critical-value
locus, rather than merely an SNC divisor containing it.  To check this
point in the preparation, first take an SNC divisor on the resolved
base containing the critical values and all the loci over which the
modifications occur; resolve its reduced inverse image, preserving
the smooth complement.  If a component of this base divisor is not
generically a critical-value divisor and the relative dimension is
positive, blow up a smooth component of its inverse image intersected
with a general very ample hypersurface.  Choose a component of this
intersection dominating the base divisor, which exists since its
generic fiber has positive dimension.  The center has normal crossings
with the other boundary components.  Locally it has equations
$t=z=0$, where $t$ is the base parameter and $z$ a fiber coordinate;
after the blowup a chart has $t=uv$.  The resulting critical locus
therefore dominates that base divisor.  Repeating for the finitely
many remaining components makes the actual critical-value locus
equal to the chosen SNC divisor.  All these new source divisors are
exceptional over $W$, and the reduced inverse image remains SNC.

Set $\Delta_\epsilon=B_W+\Lambda_\epsilon$ and take on $W'$ the
strict transform of $\Delta_\epsilon$ plus the reduced
$p$-exceptional divisor, resolving its support together with the
discriminant.  Denote this boundary by $\Delta'_\epsilon$.  Log
canonicity gives
\[
 K_{W'}+\Delta'_\epsilon
     =p^*(K_W+\Delta_\epsilon)+E_\epsilon,
               \qquad E_\epsilon\ge0
               \text{ exceptional over }W.
\]
Its adjoint is therefore pseudoeffective.  Theorem~3.4 of \cite{CP}
is applied with the full boundary $\Delta'_\epsilon$ and gives
pseudoeffectivity of
\[
 K_{W'/S'}+(\Delta'_\epsilon)_{\rm hor}-D(h'),
\]
where $D(h')$ counts ramification along primes dominating divisors of
$S'$.  Its generic-fiber hypothesis, as used in the proof of that
theorem, holds because the restriction of the pseudoeffective total
adjoint to a very general fiber is pseudoeffective.  The vertical
boundary restricts to zero there, so the fiber sees exactly
$(\Delta'_\epsilon)_{\rm hor}$.  This does not require
$K_{W'}+(\Delta'_\epsilon)_{\rm hor}$ to be pseudoeffective globally;
see also \cite[Remark~3.6]{CP}.
Adding the effective vertical boundary and $D(h')$ proves that
$T'_\epsilon=K_{W'/S'}+\Delta'_\epsilon$ is pseudoeffective.

Write $K_{S'}=q^*K_S+F$ with $F\ge0$ exceptional; this is the
canonical Jacobian divisor of a birational morphism between smooth
varieties.  Pushforward of the preceding pseudoeffective divisor gives
\[
 p_*T'_\epsilon
   =K_W+\Delta_\epsilon-h^*K_S-p_*h'^*F.
\]
Consequently $L_W+\epsilon A-h^*K_S$ is the sum of a
pseudoeffective class and the effective class $p_*h'^*F$.
Letting $\epsilon\downarrow0$ proves
\eqref{eq:nonvanishing-relative} on the original fixed $W$.
\end{proof}

\subsection{Return to the original base}

The geometric comparisons now supply all three hypotheses of the
numerical removal proposition.  Its conclusion gives a big divisor
on $W$; the exact section lattice then returns actual pluriforms to
the original base.

\begin{proposition}\label{prop:positive-period}
For the data attached to a nonzero total-space pluriform, if
$\dim S>0$ then $L_W$ is big and
\[
                      \kappa(Y,L_Y)\ge\dim W>0.
\]
\end{proposition}

\begin{proof}
Apply \Cref{prop:twist-removal} with $L=L_W$ and
$H=\mathcal H$. Relative bigness is
\Cref{prop:rank-one-base}; relative pseudoeffectivity is
\Cref{lem:nonvanishing-relative}; adjoint bigness is
\Cref{prop:adjoint-bigness}. The nef line $P$ and numerical-rank equality
are supplied by \Cref{prop:tail-rank}; its fixed comparison with $L_W$
is \Cref{lem:pseff-threshold}. Thus $L_W$ is big.

The divisible-degree comparison of \Cref{prop:rank-one-base} sends its
complete section spaces to actual sections on the original fixed base
and preserves their ratios. The images of sufficiently large divisible
systems therefore have dimension $\dim W$. Hence
$\kappa(Y,L_Y)\geq\dim W\geq\dim S>0$.
\end{proof}

\begin{proof}[Proof of \Cref{prop:nonvanishing}]
The point-quotient case, including its original-point zero assertion,
was proved at the start of the section.  For $\dim S>0$,
\Cref{lem:relative-iitaka-construction} supplies the nonnegative
generic-fiber Iitaka dimension and the diagram used above.
\Cref{prop:positive-period} then gives
$\kappa(Y,L_Y)\ge\dim W\ge\dim S>0$.  These two alternatives prove
all the assertions.
\end{proof}

\begin{proof}[Proof of \Cref{prop:section-construction}]
Apply \Cref{prop:nonvanishing} to the specified nonzero section $s$.
It gives $\kappa(Y,L_Y)\geq0$. If that dimension is zero, the
positive-dimensional quotient case is excluded. The point-quotient
case supplies a nonzero original base section vanishing at every
chosen $y\in V$ for which $s|_{X_y}=0$. This is the second assertion.
Together with the reduction in \Cref{sec:completion}, it completes the
independent upper inequality.
\end{proof}

\subsection{Equality from logarithmic subadditivity}
\label{sec:matched-equality}

We finally invoke the only companion input used in this paper. Its
statement is recorded here to make the equality's precise hypothesis
match explicit.

\begin{theorem}[Logarithmic subadditivity, {\cite[Corollary~\FoundationLogNumber]{Foundation}}]
\label{thm:log-iitaka}
Let $g\colon Z\to T$ be a surjective morphism with connected fibers between
smooth connected projective complex varieties. Let $B_Z,B_T$ be reduced
SNC divisors with
$\Supp(g^*B_T)\subseteq\Supp(B_Z)$, and let
$G$ be a very general smooth fiber with $B_G=B_Z|_G$. Then
\begin{equation}\label{eq:lower-bound}
 \kappa(Z,K_Z+B_Z)
 \geq \kappa(T,K_T+B_T)+\kappa(G,K_G+B_G).
\end{equation}
No smoothness outside $B_T$, abundance, or good-minimal-model hypothesis
is imposed.
\end{theorem}

To prove \Cref{cor:additivity}, apply logarithmic subadditivity
\cite[Corollary~\FoundationLogNumber]{Foundation}, stated in \Cref{thm:log-iitaka}, with
$Z=X$, $T=Y$, $g=f$, $B_Z=E$, and $B_T=D$. The source and base are smooth
connected projective complex varieties, the morphism is surjective with
connected fibers, the boundaries are reduced SNC, and
\eqref{eq:support-inclusion} is exactly its boundary hypothesis.
The very general smooth fiber and its restricted boundary are precisely
$F,E_F$. This theorem supplies the lower inequality; \Cref{thm:upper-bound}
supplies the upper inequality. Their infinity conventions agree, giving
\eqref{eq:additivity} in every case.


\bibliographystyle{amsplain}
\bibliography{references}
\end{document}
